% =====================================================================
%  CAPITULO 1 - ELETROSTATICA
%  Subsecao 1.1 - Cargas eletricas  (versao revisada)
% =====================================================================
\section{Eletrostática}

\begin{conceito}[Antes de começar]
\textbf{Quantização da carga:} $Q = n\,e$, com $e = \SI{1.6e-19}{\coulomb}$.\par
\textbf{Conservação da carga:} em um sistema isolado, a soma algébrica das cargas é constante.\par
\textbf{Contato entre esferas idênticas:} a carga total se divide igualmente,
$Q'_1 = Q'_2 = \dfrac{Q_1+Q_2}{2}$. Para esferas de raios diferentes, a divisão é
proporcional ao raio: $\dfrac{Q'_1}{R_1} = \dfrac{Q'_2}{R_2}$.\par
\textbf{Indução:} a aproximação de um corpo eletrizado separa cargas no condutor
neutro, que continua com carga total nula até ser aterrado.
\end{conceito}

\legendaniveis

\subsection{Cargas elétricas}

% ---------------------------------------------------------------
% NIVEL 1 - FIXACAO
% ---------------------------------------------------------------

\blocofixacao

\begin{questao}[1]{}
O princípio da conservação da carga elétrica estabelece que:
\begin{alternativas}
\item as cargas elétricas de mesmo sinal se repelem.
\item cargas elétricas de sinais opostos se atraem.
\item a soma algébrica das cargas elétricas é constante em um sistema eletricamente isolado.
\item a soma das cargas positivas e negativas é diferente de zero em um sistema neutro.
\item os elétrons livres se atraem.
\end{alternativas}
\resposta{(c)}
\resolucao{A conservação da carga afirma que, em um sistema isolado, cargas podem
ser transferidas de um corpo para outro, mas a soma algébrica permanece a mesma.
As alternativas (a) e (b) descrevem a \emph{lei de atração e repulsão}, e não a
conservação; (d) e (e) são falsas.}
\end{questao}

\begin{questao}[1]{}
Dois corpos de materiais diferentes, inicialmente neutros, são atritados entre si,
de forma isolada do meio. É correto afirmar que:
\begin{alternativas}
\item um fica eletrizado positivamente e o outro, negativamente.
\item um fica eletrizado negativamente e o outro permanece neutro.
\item um fica eletrizado positivamente e o outro permanece neutro.
\item ambos ficam eletrizados negativamente.
\item ambos ficam eletrizados positivamente.
\end{alternativas}
\resposta{(a)}
\resolucao{Na eletrização por atrito há apenas \emph{transferência} de elétrons de
um corpo para o outro. Pela conservação da carga, o que um perde o outro ganha:
as cargas finais têm o mesmo módulo e sinais opostos.}
\end{questao}

\begin{questao}[1]{}
Atrita-se um bastão de vidro com um pano de lã, ambos inicialmente neutros.
Pode-se afirmar que:
\begin{alternativas}
\item só a lã fica eletrizada.
\item só o bastão fica eletrizado.
\item o bastão e a lã se eletrizam com cargas de mesmo sinal.
\item o bastão e a lã se eletrizam com cargas de mesmo valor absoluto e sinais opostos.
\item nenhuma das anteriores.
\end{alternativas}
\resposta{(d)}
\resolucao{Mesma justificativa da questão anterior. Na série triboelétrica o vidro
cede elétrons à lã: o vidro fica positivo e a lã, negativa.}
\end{questao}

\begin{questao}[1]{}
A tabela apresenta parte de uma \textbf{série triboelétrica}: ao serem atritados,
o material que aparece \emph{mais acima} tende a ceder elétrons ao que aparece
mais abaixo.
\begin{table}[H]
\centering
\small\sffamily
\setlength{\arrayrulewidth}{0.5pt}
\begin{tabular}{@{}c l@{}}
\toprule
\textbf{Ordem} & \textbf{Material}\\
\midrule
1 & Pele de coelho\\
2 & Vidro\\
3 & Lã\\
4 & Seda\\
5 & Plástico (PVC)\\
6 & Teflon\\
\bottomrule
\end{tabular}
\caption{Série triboelétrica simplificada.}
\end{table}
Atritando-se um bastão de \textbf{PVC} com um pedaço de \textbf{lã}, os sinais das
cargas adquiridas pelo PVC e pela lã são, respectivamente:
\begin{alternativas}[2]
\item positivo e positivo.
\item positivo e negativo.
\item negativo e positivo.
\item negativo e negativo.
\item nulo e nulo.
\end{alternativas}
\resposta{(c)}
\resolucao{A lã (posição 3) está acima do PVC (posição 5), logo cede elétrons ao PVC.
O PVC \emph{ganha} elétrons (fica negativo) e a lã \emph{perde} elétrons (fica positiva).}
\end{questao}

\begin{questao}[1]{}
Duas cargas elétricas $Q_1$ e $Q_2$ atraem-se quando colocadas próximas uma da outra.
\begin{itensq}
\item O que se pode afirmar sobre os sinais de $Q_1$ e $Q_2$?
\item Sabe-se que $Q_1$ é repelida por uma terceira carga $Q_3$, positiva.
      Qual é o sinal de $Q_2$?
\end{itensq}
\resposta{(a) sinais opostos; (b) negativo}
\resolucao{(a) A atração ocorre entre cargas de sinais contrários (ou entre corpo
eletrizado e corpo neutro; como o enunciado fala em \emph{cargas}, os sinais são
opostos).\\
(b) $Q_1$ é repelida por $Q_3>0$, logo $Q_1>0$. Como $Q_1$ e $Q_2$ se atraem,
$Q_2<0$.}
\end{questao}

\begin{questao}[1]{}
Com relação à eletrização de um corpo, assinale \textbf{todas} as afirmativas corretas:
\begin{alternativas}
\item Um corpo eletricamente neutro que perde elétrons fica eletrizado positivamente.
\item Um corpo eletricamente neutro não possui cargas elétricas.
\item Um dos processos de eletrização consiste em retirar prótons do corpo.
\item Um corpo eletricamente neutro não pode ser atraído por um corpo eletrizado.
\item Friccionando-se dois corpos constituídos do mesmo material, um se eletriza
      positivamente e o outro, negativamente.
\end{alternativas}
\resposta{Apenas (a)}
\resolucao{(a) Correta. (b) Errada: o corpo neutro possui igual número de prótons
e elétrons, não ausência de cargas. (c) Errada: prótons estão presos ao núcleo;
nos sólidos só os elétrons se movem. (d) Errada: por indução, um corpo neutro é
atraído por um corpo eletrizado. (e) Errada: materiais iguais têm a mesma
tendência de ceder ou receber elétrons, e o atrito não produz eletrização
apreciável.}
\end{questao}

\begin{questao}[1]{}
Em uma cabine de \textbf{pintura eletrostática}, as gotículas de tinta são
eletrizadas negativamente e a peça metálica a ser pintada é ligada ao terminal
positivo da fonte. Esse arranjo é usado porque:
\begin{alternativas}
\item as gotículas se repelem entre si e são atraídas pela peça, reduzindo o desperdício.
\item as gotículas se atraem entre si, formando gotas maiores.
\item a tinta se torna condutora e adere melhor.
\item o processo elimina a necessidade de secagem.
\item a peça precisa estar eletrizada com o mesmo sinal da tinta.
\end{alternativas}
\resposta{(a)}
\resolucao{Gotículas de mesmo sinal se repelem, o que espalha o jato e evita
acúmulos; simultaneamente, todas são atraídas pela peça de sinal oposto,
inclusive contornando-a (efeito \emph{wrap-around}). O aproveitamento da tinta
passa de cerca de \SI{50}{\percent} para mais de \SI{90}{\percent}.}
\end{questao}

\begin{questao}[1]{}
Um corpo, inicialmente neutro, é eletrizado com carga $Q = \SI{48}{\micro\coulomb}$.
\dados{carga elementar $e = \SI{1.6e-19}{\coulomb}$.}
\begin{itensq}
\item O corpo ficou com falta ou com excesso de elétrons?
\item Qual é o número de elétrons retirado do corpo ou a ele fornecido?
\end{itensq}
\resposta{(a) falta de elétrons; (b) $n = \pot{3.0}{14}$ elétrons}
\resolucao{(a) A carga final é positiva, logo o corpo \emph{perdeu} elétrons
(ficou com falta).\\
(b) Da quantização da carga, $Q = n\,e$:
\[
n=\frac{Q}{e}=\frac{\SI{48e-6}{\coulomb}}{\SI{1.6e-19}{\coulomb}}
 = \pot{3.0}{14}\ \text{elétrons}.
\]}
\end{questao}

\begin{questao}[1]{}
Um bastão carregado positivamente atrai um objeto isolado, suspenso por um fio.
Sobre esse objeto, é correto afirmar que:
\begin{alternativas}
\item necessariamente possui elétrons em excesso.
\item é obrigatoriamente condutor.
\item trata-se de um isolante.
\item está carregado positivamente.
\item pode estar neutro.
\end{alternativas}
\resposta{(e)}
\resolucao{A atração é compatível com duas situações: objeto negativo, ou objeto
\emph{neutro} polarizado por indução. Como não há informação adicional, a única
afirmação segura é que ele \emph{pode} estar neutro --- por isso (a) e (d) estão
erradas. A polarização ocorre tanto em condutores quanto em isolantes, o que
elimina (b) e (c).}
\end{questao}

\blocoaplicacao

\begin{questao}[2]{}
Uma barra eletrizada negativamente é aproximada de um corpo metálico AB,
inicialmente neutro e apoiado sobre suportes isolantes, conforme a figura.

\begin{figure}[H]
\centering
\begin{tikzpicture}[scale=1,>=Stealth]
  % barra eletrizada
  \draw[fill=laranja!15,draw=laranja,thick,rounded corners=2pt]
        (-4.3,-0.3) rectangle (-2.6,0.3);
  \foreach \x in {-4.1,-3.7,-3.3,-2.9}
     \node[laranja,font=\bfseries\small] at (\x,0) {$-$};
  \node[below,font=\footnotesize\sffamily] at (-3.45,-0.45) {barra eletrizada};
  % corpo metalico
  \shade[left color=cinza!35,right color=cinza!15,draw=cinza,thick,rounded corners=6pt]
        (-1.9,-0.45) rectangle (2.4,0.45);
  \node[font=\bfseries] at (-1.55,0) {A};
  \node[font=\bfseries] at ( 2.05,0) {B};
  % suportes isolantes
  \foreach \x in {-1.2,1.7}{
    \draw[fill=cinza!25,draw=cinza] (\x-0.18,-1.1) rectangle (\x+0.18,-0.45);
    \draw[fill=cinza!45,draw=cinza] (\x-0.45,-1.25) rectangle (\x+0.45,-1.1);}
  \node[font=\footnotesize\sffamily,cinza] at (0.25,-1.5) {suportes isolantes};
  % linha de base
  \draw[cinza!60] (-2.6,-1.25) -- (2.9,-1.25);
\end{tikzpicture}
\caption{Barra negativa aproximada do condutor neutro AB.}
\label{fig:inducao-ab}
\end{figure}

Podemos afirmar que:
\begin{alternativas}
\item não haverá movimento de elétrons livres no corpo AB.
\item os elétrons livres do corpo AB deslocam-se para a extremidade A.
\item o sinal da carga que aparece em B é positivo.
\item ocorre no corpo metálico o fenômeno da indução eletrostática.
\item após a separação de cargas, a carga total do corpo é não nula.
\end{alternativas}
\resposta{(d)}
\resolucao{A barra negativa \emph{repele} os elétrons livres do metal, que migram
para a extremidade \emph{mais distante}, B. Logo B fica negativa e A, positiva
--- o que elimina (a), (b) e (c). Como não houve aterramento nem contato, a carga
total de AB continua nula, eliminando (e). O fenômeno descrito é exatamente a
indução eletrostática.}
\end{questao}

\begin{questao}[2]{}
Um corpo eletrizado \textbf{positivamente} é aproximado de um corpo metálico
neutro, cujas extremidades são X (mais próxima) e Y (mais afastada).

\begin{figure}[H]
\centering
\begin{tikzpicture}[scale=1,>=Stealth]
  \draw[fill=Vcol!12,draw=Vcol,thick,rounded corners=2pt] (-4.2,-0.3) rectangle (-2.7,0.3);
  \foreach \x in {-4.0,-3.6,-3.2,-2.9} \node[Vcol,font=\bfseries\small] at (\x,0) {$+$};
  \node[below,font=\footnotesize\sffamily] at (-3.45,-0.45) {corpo eletrizado};
  \shade[left color=cinza!35,right color=cinza!15,draw=cinza,thick,rounded corners=6pt]
        (-2.0,-0.45) rectangle (2.3,0.45);
  \node[font=\bfseries] at (-1.65,0) {X};
  \node[font=\bfseries] at ( 1.95,0) {Y};
  \foreach \x in {-1.3,1.6}{
    \draw[fill=cinza!25,draw=cinza] (\x-0.18,-1.1) rectangle (\x+0.18,-0.45);
    \draw[fill=cinza!45,draw=cinza] (\x-0.45,-1.25) rectangle (\x+0.45,-1.1);}
  \draw[cinza!60] (-2.7,-1.25) -- (2.8,-1.25);
\end{tikzpicture}
\caption{Indução eletrostática em um condutor neutro.}
\label{fig:inducao-xy}
\end{figure}

Podemos afirmar que:
\begin{alternativas}
\item não haverá movimentação de cargas negativas no corpo neutro.
\item a carga que aparece em X é positiva.
\item a carga que aparece em Y é negativa.
\item haverá força de interação elétrica entre os dois corpos.
\item todas as afirmativas acima estão erradas.
\end{alternativas}
\resposta{(d)}
\resolucao{O corpo positivo \emph{atrai} os elétrons livres: X fica negativa e Y,
positiva --- o que torna (a), (b) e (c) falsas. Como a extremidade negativa X está
\emph{mais próxima} da barra do que a positiva Y, a atração supera a repulsão e
existe força resultante de atração entre os corpos.}
\end{questao}

\begin{questao}[2]{}
Duas esferas condutoras descarregadas, $x$ e $y$, colocadas sobre suportes
isolantes, estão inicialmente em contato. Um bastão carregado positivamente é
aproximado da esfera $x$, como mostra a figura.

\begin{figure}[H]
\centering
\begin{tikzpicture}[scale=1]
  \draw[fill=Vcol!12,draw=Vcol,thick,rounded corners=2pt] (-4.6,-0.25) rectangle (-3.2,0.25);
  \foreach \x in {-4.4,-4.0,-3.6,-3.35} \node[Vcol,font=\bfseries\small] at (\x,0) {$+$};
  % esferas
  \shade[ball color=cinza!30] (-1.9,0) circle (0.75);
  \shade[ball color=cinza!30] ( 0.1,0) circle (0.75);
  \draw[cinza,thick] (-1.9,0) circle (0.75);
  \draw[cinza,thick] ( 0.1,0) circle (0.75);
  \node at (-1.9,0) {\bfseries $x$};
  \node at ( 0.1,0) {\bfseries $y$};
  % contato
  \draw[line width=2pt,cinza] (-1.15,0) -- (-0.65,0);
  % suportes
  \foreach \x in {-1.9,0.1}{
    \draw[fill=cinza!25,draw=cinza] (\x-0.12,-1.55) rectangle (\x+0.12,-0.75);
    \draw[fill=cinza!45,draw=cinza] (\x-0.42,-1.7) rectangle (\x+0.42,-1.55);}
  \draw[cinza!60] (-3.0,-1.7) -- (1.1,-1.7);
\end{tikzpicture}
\caption{Eletrização por indução com separação das esferas.}
\label{fig:inducao-xy2}
\end{figure}

Em seguida, a esfera $y$ é afastada de $x$, \textbf{mantendo-se o bastão em sua
posição}. Ao final, as cargas de $x$ e $y$ são, respectivamente:
\begin{alternativas}[2]
\item nula e positiva.
\item negativa e positiva.
\item nula e nula.
\item negativa e nula.
\item positiva e negativa.
\end{alternativas}
\resposta{(b)}
\resolucao{Com as esferas em contato, elas formam um único condutor: o bastão
positivo atrai os elétrons para $x$ e deixa $y$ com falta de elétrons.
\textbf{Separando as esferas com o bastão ainda presente}, essa separação de
cargas fica "aprisionada": $x$ permanece negativa e $y$, positiva.
(Se o bastão fosse afastado \emph{antes} da separação, ambas voltariam a ser
neutras.)}
\end{questao}

\begin{questao}[2]{}
A figura mostra um eletroscópio de folhas, inicialmente neutro. Um bastão
eletrizado negativamente é aproximado do disco, \textbf{sem tocá-lo}.

\begin{figure}[H]
\centering
\begin{tikzpicture}[scale=1]
  % bastao
  \draw[fill=laranja!15,draw=laranja,thick,rounded corners=2pt] (-2.9,1.9) rectangle (-1.3,2.3);
  \foreach \x in {-2.7,-2.35,-2.0,-1.6} \node[laranja,font=\bfseries\small] at (\x,2.1) {$-$};
  % disco
  \draw[fill=cinza!30,draw=cinza,thick] (-0.9,2.0) rectangle (0.9,2.2);
  % haste
  \draw[line width=1.6pt,cinza] (0,2.0) -- (0,0.55);
  % frasco de vidro
  \draw[azulprof!35,thick,fill=azulclaro!35]
        (-1.05,-1.5) -- (-1.05,1.35) .. controls (-1.05,1.65) and (-0.6,1.7) ..
        (-0.45,1.7) -- (0.45,1.7) .. controls (0.6,1.7) and (1.05,1.65) ..
        (1.05,1.35) -- (1.05,-1.5) -- cycle;
  \draw[fill=cinza!45,draw=cinza] (-1.25,-1.65) rectangle (1.25,-1.5);
  % folhas
  \draw[line width=1.1pt,laranja] (0,0.55) -- (-0.42,-0.75);
  \draw[line width=1.1pt,laranja] (0,0.55) -- ( 0.42,-0.75);
  \node[font=\footnotesize\sffamily,cinza,right] at (1.35,-0.2) {folhas metálicas};
  \node[font=\footnotesize\sffamily,cinza,right] at (1.35,2.1) {disco};
\end{tikzpicture}
\caption{Eletroscópio de folhas sob indução.}
\label{fig:eletroscopio}
\end{figure}

Sobre a situação descrita, é correto afirmar que:
\begin{alternativas}
\item as folhas permanecem fechadas, pois o eletroscópio continua neutro.
\item as folhas se abrem e o disco fica com carga positiva.
\item as folhas se abrem e o disco fica com carga negativa.
\item as folhas se abrem porque o eletroscópio adquiriu carga negativa total.
\item as folhas se fecham ainda mais.
\end{alternativas}
\resposta{(b)}
\resolucao{O bastão negativo repele os elétrons livres do disco, que descem para
as folhas. O disco fica \emph{positivo} e as folhas ficam \emph{negativas}; como
ambas têm o mesmo sinal, elas se repelem e se abrem. A carga total do
eletroscópio, porém, continua nula --- o que elimina (a) e (d).}
\end{questao}

\begin{questao}[2]{}
Três bolas metálicas podem ser carregadas eletricamente. Observa-se que cada uma
das três bolas \textbf{atrai uma e repele outra}. Considere as hipóteses:

\smallskip
\begin{tabular}{@{}ll@{}}
I   & Apenas uma das bolas está carregada.\\
II  & Duas das bolas estão carregadas.\\
III & As três bolas estão carregadas.\\
\end{tabular}
\smallskip

O fenômeno pode ser explicado:
\begin{alternativas}
\item somente pelas hipóteses II ou III.
\item somente pela hipótese I.
\item somente pela hipótese III.
\item somente pela hipótese II.
\item somente pelas hipóteses I ou II.
\end{alternativas}
\resposta{(c)}
\resolucao{A \emph{repulsão} só ocorre entre corpos eletrizados de mesmo sinal.
Como cada bola repele alguma outra, cada uma das três precisa estar eletrizada:
apenas a hipótese III explica o fenômeno. Um arranjo possível é $(+,+,-)$.}
\end{questao}

\begin{questao}[2]{}
Três esferas metálicas podem ser carregadas eletricamente. Aproximando-as duas a
duas, observa-se que \textbf{em todos os casos ocorre atração}. Considere:

\smallskip
\begin{tabular}{@{}ll@{}}
I   & Somente uma das esferas está carregada.\\
II  & Duas esferas estão carregadas.\\
III & As três esferas estão carregadas.\\
\end{tabular}
\smallskip

Quais hipóteses explicam o fenômeno descrito?
\begin{alternativas}
\item Apenas a hipótese I.
\item Apenas a hipótese II.
\item Apenas a hipótese III.
\item Apenas as hipóteses I e II.
\item Nenhuma das três hipóteses.
\end{alternativas}
\resposta{(d)}
\resolucao{Se as três estivessem carregadas, haveria pelo menos um par de mesmo
sinal (princípio da casa dos pombos com dois sinais e três esferas), gerando
repulsão --- III fica descartada. Com apenas uma carregada (I), ela atrai as duas
neutras por indução e as neutras se atraem mutuamente por polarização mútua.
Com duas carregadas de sinais opostos (II), elas se atraem entre si e cada uma
atrai a neutra. Logo I e II são possíveis.}
\end{questao}

\begin{questao}[2]{}
Uma estudante colocou sobre uma mesa horizontal três pêndulos eletrostáticos
idênticos, equidistantes entre si, como se cada um ocupasse o vértice de um
triângulo equilátero. As esferas metálicas dos pêndulos atraíram-se mutuamente.
A estudante pode concluir corretamente que:
\begin{alternativas}
\item as três esferas estavam eletrizadas com cargas de mesmo sinal.
\item duas esferas estavam eletrizadas com cargas de mesmo sinal e uma com sinal oposto.
\item duas esferas estavam eletrizadas com cargas de mesmo sinal e uma neutra.
\item duas esferas estavam eletrizadas com cargas de sinais opostos e uma neutra.
\item uma esfera estava eletrizada e duas, neutras.
\end{alternativas}
\resposta{(d)}
\resolucao{Cargas de mesmo sinal se repeliriam, o que elimina (a) e (c).
Em (b), as duas de mesmo sinal se repeliriam. Em (e), as duas esferas neutras
não interagiriam apreciavelmente entre si. Resta (d): as duas eletrizadas se
atraem por terem sinais opostos, e cada uma atrai a neutra por indução.}
\end{questao}

\begin{questao}[2]{}
Eletriza-se por atrito um bastão de plástico com um pedaço de papel. Em seguida,
o bastão é aproximado de um pêndulo eletrostático já eletrizado e observa-se
\textbf{repulsão}. Em qual alternativa a carga de cada elemento corresponde a
essa descrição?

\begin{table}[H]
\centering
\small\sffamily
\setlength{\arraystretch}{1.25}
\begin{tabular}{@{}c c c c@{}}
\toprule
 & \textbf{Papel} & \textbf{Bastão} & \textbf{Pêndulo}\\
\midrule
a) & positiva & positiva & positiva\\
b) & negativa & positiva & negativa\\
c) & negativa & negativa & positiva\\
d) & positiva & positiva & negativa\\
e) & positiva & negativa & negativa\\
\bottomrule
\end{tabular}
\caption{Alternativas da questão.}
\end{table}
\resposta{(e)}
\resolucao{Duas condições devem ser satisfeitas simultaneamente:
\textbf{(i)} bastão e papel, eletrizados por atrito, têm sinais \emph{opostos};
\textbf{(ii)} bastão e pêndulo se repelem, logo têm sinais \emph{iguais}.
Somente a alternativa (e) satisfaz ambas (papel $+$, bastão $-$, pêndulo $-$).
Em (a) o par bastão/papel tem sinais iguais; em (b) e (d) o bastão e o pêndulo
têm sinais opostos; em (c) o par bastão/papel tem sinais iguais.}
\end{questao}

\begin{questao}[2]{}
Duas esferas condutoras de dimensões iguais, A e B, estão eletrizadas com
$Q_A = \SI{+5}{\coulomb}$ e $Q_B = \SI{-1}{\coulomb}$, e são ligadas por um fio
condutor com uma chave C, inicialmente aberta.

\begin{figure}[H]
\centering
\begin{tikzpicture}[scale=1]
  \shade[ball color=Vcol!25] (-2.6,0) circle (0.8);
  \shade[ball color=Icol!25] ( 2.6,0) circle (0.8);
  \draw[Vcol,thick] (-2.6,0) circle (0.8);
  \draw[Icol,thick] ( 2.6,0) circle (0.8);
  \node[font=\bfseries] at (-2.6,0.18) {A};
  \node[font=\small] at (-2.6,-0.28) {$+\SI{5}{\coulomb}$};
  \node[font=\bfseries] at ( 2.6,0.18) {B};
  \node[font=\small] at ( 2.6,-0.28) {$-\SI{1}{\coulomb}$};
  % fio com chave
  \draw[line width=1.1pt,cinza] (-1.8,0) -- (-0.55,0);
  \draw[line width=1.1pt,cinza] ( 0.55,0) -- ( 1.8,0);
  \fill[cinza] (-0.55,0) circle (0.06);
  \fill[cinza] ( 0.55,0) circle (0.06);
  \draw[line width=1.1pt,cinza] (-0.55,0) -- (0.42,0.52);
  \node[font=\footnotesize\sffamily,cinza,above] at (0,0.62) {C (aberta)};
  % suportes
  \foreach \x in {-2.6,2.6}{
    \draw[fill=cinza!25,draw=cinza] (\x-0.12,-1.6) rectangle (\x+0.12,-0.8);
    \draw[fill=cinza!45,draw=cinza] (\x-0.42,-1.75) rectangle (\x+0.42,-1.6);}
  \draw[cinza!60] (-3.6,-1.75) -- (3.6,-1.75);
\end{tikzpicture}
\caption{Esferas idênticas ligadas por fio condutor com chave.}
\label{fig:esferas-chave-ab}
\end{figure}

Quando a chave é fechada, passam elétrons:
\begin{alternativas}
\item de A para B, e a nova carga de A é $\SI{+2}{\coulomb}$.
\item de A para B, e a nova carga de B é $\SI{-1}{\coulomb}$.
\item de B para A, e a nova carga de A é $\SI{+1}{\coulomb}$.
\item de B para A, e a nova carga de B é $\SI{-1}{\coulomb}$.
\item de B para A, e a nova carga de A é $\SI{+2}{\coulomb}$.
\end{alternativas}
\resposta{(e)}
\resolucao{As esferas são idênticas, então a carga total se divide igualmente:
\[
Q'_A=Q'_B=\frac{Q_A+Q_B}{2}=\frac{(+5)+(-1)}{2}=\SI{+2}{\coulomb}.
\]
A esfera A passou de $+5$ para $+2$: ela \emph{ganhou} \SI{3}{\coulomb} em
elétrons. Como os elétrons são negativos, eles se deslocam \emph{para} o corpo de
maior potencial, ou seja, de B para A.}
\end{questao}

\begin{questao}[2]{}
Duas esferas condutoras idênticas, K e L, com cargas
$Q_K = \SI{+3}{\micro\coulomb}$ e $Q_L = \SI{+1}{\micro\coulomb}$, são ligadas por
um fio condutor com uma chave C, inicialmente aberta.

\begin{figure}[H]
\centering
\begin{tikzpicture}[scale=1]
  \shade[ball color=Vcol!25] (-2.6,0) circle (0.8);
  \shade[ball color=Vcol!12] ( 2.6,0) circle (0.8);
  \draw[Vcol,thick] (-2.6,0) circle (0.8);
  \draw[Vcol,thick] ( 2.6,0) circle (0.8);
  \node[font=\bfseries] at (-2.6,0.18) {K};
  \node[font=\small] at (-2.6,-0.30) {$+\SI{3}{\micro\coulomb}$};
  \node[font=\bfseries] at ( 2.6,0.18) {L};
  \node[font=\small] at ( 2.6,-0.30) {$+\SI{1}{\micro\coulomb}$};
  \draw[line width=1.1pt,cinza] (-1.8,0) -- (-0.55,0);
  \draw[line width=1.1pt,cinza] ( 0.55,0) -- ( 1.8,0);
  \fill[cinza] (-0.55,0) circle (0.06);
  \fill[cinza] ( 0.55,0) circle (0.06);
  \draw[line width=1.1pt,cinza] (-0.55,0) -- (0.42,0.52);
  \node[font=\footnotesize\sffamily,cinza,above] at (0,0.62) {C (aberta)};
  \foreach \x in {-2.6,2.6}{
    \draw[fill=cinza!25,draw=cinza] (\x-0.12,-1.6) rectangle (\x+0.12,-0.8);
    \draw[fill=cinza!45,draw=cinza] (\x-0.42,-1.75) rectangle (\x+0.42,-1.6);}
  \draw[cinza!60] (-3.6,-1.75) -- (3.6,-1.75);
\end{tikzpicture}
\caption{Esferas idênticas com cargas de mesmo sinal.}
\label{fig:esferas-chave-kl}
\end{figure}

Quando a chave é fechada:
\begin{alternativas}
\item não ocorre fluxo de elétrons, pois ambas estão com cargas positivas.
\item passam elétrons de L para K, transportando carga de $\SI{+2}{\micro\coulomb}$.
\item passam elétrons de L para K, transportando carga de $\SI{-1}{\micro\coulomb}$.
\item passam elétrons de K para L, transportando carga de $\SI{+1}{\micro\coulomb}$.
\item passam elétrons de K para L, transportando carga de $\SI{-1}{\micro\coulomb}$.
\end{alternativas}
\resposta{(c)}
\resolucao{Carga final de cada esfera:
$Q' = \dfrac{3+1}{2} = \SI{+2}{\micro\coulomb}$.
A esfera K passou de $+3$ para $+2\,\si{\micro\coulomb}$: recebeu
$\SI{-1}{\micro\coulomb}$ (isto é, elétrons). Portanto os elétrons vão de L para K
transportando $\SI{-1}{\micro\coulomb}$. O que impulsiona o fluxo é a diferença de
\emph{potencial}, e não o sinal das cargas --- por isso (a) é falsa.}
\end{questao}

\begin{questao}[2]{}
Duas esferas metálicas idênticas estão carregadas com $5Q$ e $-3Q$.
Encostando-se uma na outra e separando-as em seguida, cada esfera ficará com
carga elétrica igual a:
\begin{alternativas}[3]
\item $8Q$
\item $5Q$
\item $3Q$
\item $2Q$
\item $Q$
\end{alternativas}
\resposta{(e)}
\resolucao{$Q' = \dfrac{5Q + (-3Q)}{2} = \dfrac{2Q}{2} = Q$ em cada esfera.}
\end{questao}

\begin{questao}[2]{}
Uma pequena esfera metálica está eletrizada com carga $\pot{8.0}{-8}\,\si{\coulomb}$.
Colocando-a em contato com outra idêntica, porém neutra, o número de elétrons que
passa de uma esfera para a outra é:
\dados{$e = \SI{1.6e-19}{\coulomb}$.}
\begin{alternativas}[3]
\item $\pot{4.0}{12}$
\item $\pot{4.0}{11}$
\item $\pot{4.0}{10}$
\item $\pot{2.5}{12}$
\item $\pot{2.5}{11}$
\end{alternativas}
\resposta{(e)}
\resolucao{Após o contato, cada esfera fica com
$Q' = \dfrac{\pot{8.0}{-8}+0}{2} = \pot{4.0}{-8}\,\si{\coulomb}$.
A carga transferida tem módulo $\Delta Q = \pot{4.0}{-8}\,\si{\coulomb}$, logo
\[
n=\frac{\Delta Q}{e}=\frac{\pot{4.0}{-8}}{\pot{1.6}{-19}}=\pot{2.5}{11}
\ \text{elétrons}.
\]
Note que a esfera eletrizada é positiva: os elétrons vão da esfera neutra para ela.}
\end{questao}

\begin{questao}[2]{}
Um estudante afirma ter medido, em um pequeno corpo eletrizado, a carga
$Q = \pot{2.4}{-19}\,\si{\coulomb}$. Sabendo que $e = \SI{1.6e-19}{\coulomb}$,
esse resultado é possível? Justifique.
\resposta{Não: a carga não é múltipla inteira de $e$.}
\resolucao{Pela quantização, toda carga deve satisfazer $Q=n\,e$ com $n$ inteiro.
Aqui,
\[
n=\frac{\pot{2.4}{-19}}{\pot{1.6}{-19}}=1{,}5,
\]
que não é inteiro. Não existe carga livre igual a $1{,}5\,e$, logo a medida está
incorreta (erro experimental ou de cálculo).}
\end{questao}

\blocoaprofundamento

\begin{questao}[3]{}
Três esferas metálicas A, B e C, idênticas e no vácuo, encontram-se inicialmente
com A carregada com $+Q$ e B e C neutras. A esfera A é colocada em contato com B
e, posteriormente, com C. Os valores finais das cargas de A, B e C são,
respectivamente:
\begin{alternativas}[2]
\item $\dfrac{Q}{3},\ \dfrac{Q}{3},\ \dfrac{Q}{3}$
\item $\dfrac{Q}{2},\ \dfrac{Q}{2},\ \dfrac{Q}{2}$
\item $\dfrac{Q}{4},\ \dfrac{Q}{4},\ \dfrac{Q}{4}$
\item $\dfrac{Q}{4},\ \dfrac{Q}{2},\ \dfrac{Q}{4}$
\item $\dfrac{Q}{2},\ \dfrac{Q}{2},\ \dfrac{Q}{4}$
\end{alternativas}
\resposta{(d)}
\resolucao{\textbf{1\textsuperscript{o} contato (A--B):}
$Q_A = Q_B = \dfrac{Q+0}{2} = \dfrac{Q}{2}$.\\
\textbf{2\textsuperscript{o} contato (A--C):}
$Q_A = Q_C = \dfrac{Q/2+0}{2} = \dfrac{Q}{4}$.\\
Situação final: $A=\dfrac{Q}{4}$, $B=\dfrac{Q}{2}$, $C=\dfrac{Q}{4}$.
O erro mais comum é esquecer que B \emph{conserva} a carga obtida no primeiro
contato.}
\end{questao}

\begin{questao}[3]{}
Quatro esferas metálicas idênticas estão isoladas entre si. As esferas A, B e C
estão neutras e a esfera D está eletrizada com carga $Q$. A esfera D é colocada
em contato sucessivamente com A, depois com B e, por fim, com C.
Ao final do processo, as cargas de C e D são, respectivamente:
\begin{alternativas}[2]
\item $\dfrac{Q}{8}$ e $\dfrac{Q}{8}$
\item $\dfrac{Q}{8}$ e $\dfrac{Q}{4}$
\item $\dfrac{Q}{4}$ e $\dfrac{Q}{8}$
\item $\dfrac{Q}{2}$ e $\dfrac{Q}{2}$
\item $Q$ e $-Q$
\end{alternativas}
\resposta{(a)}
\resolucao{A cada contato com uma esfera \emph{neutra} e idêntica, a carga de D
é dividida por 2:
\[
Q \xrightarrow{\text{A}} \frac{Q}{2}
  \xrightarrow{\text{B}} \frac{Q}{4}
  \xrightarrow{\text{C}} \frac{Q}{8}.
\]
No último contato, C fica com o mesmo valor de D: ambas com $\dfrac{Q}{8}$.
De modo geral, após $n$ contatos com esferas neutras idênticas,
$Q_D = \dfrac{Q}{2^{\,n}}$.}
\end{questao}

\begin{questao}[3]{}
Considere três esferas metálicas X, Y e Z, de diâmetros iguais. Y e Z estão fixas
e suficientemente afastadas para que efeitos de indução sejam desprezíveis. A
situação inicial é: X neutra, Y com carga $+Q$ e Z com carga $-Q$. As esferas não
trocam cargas com o ambiente. Fazendo-se X tocar \textbf{primeiro} Y e
\textbf{depois} Z, a carga final de X será:
\begin{alternativas}[2]
\item zero
\item $\dfrac{2Q}{3}$
\item $-\dfrac{Q}{2}$
\item $\dfrac{Q}{8}$
\item $-\dfrac{Q}{4}$
\end{alternativas}
\resposta{(e)}
\resolucao{\textbf{Contato X--Y:} $Q_X = Q_Y = \dfrac{0+Q}{2} = \dfrac{Q}{2}$.\\
\textbf{Contato X--Z:} agora X tem $+\dfrac{Q}{2}$ e Z tem $-Q$:
\[
Q_X=\frac{\dfrac{Q}{2}+(-Q)}{2}=\frac{-\dfrac{Q}{2}}{2}=-\frac{Q}{4}.
\]
Observe a inversão de sinal: X termina \emph{negativa}. A alternativa (a) é a
armadilha típica de quem soma as cargas iniciais de Y e Z.}
\end{questao}

\begin{questao}[3]{}
Duas esferas metálicas idênticas, A e B, estão carregadas com
$\SI{16}{\micro\coulomb}$ e $\SI{4}{\micro\coulomb}$, respectivamente. Uma terceira
esfera C, idêntica às anteriores, está inicialmente descarregada. Coloca-se C em
contato com A; desfeito o contato, C é colocada em contato com B. Supondo que não
haja troca de cargas com o exterior, a carga final de C é:
\begin{alternativas}[3]
\item $\SI{8}{\micro\coulomb}$
\item $\SI{6}{\micro\coulomb}$
\item $\SI{4}{\micro\coulomb}$
\item $\SI{3}{\micro\coulomb}$
\item nula
\end{alternativas}
\resposta{(b)}
\resolucao{\textbf{Contato C--A:}
$Q_C = Q_A = \dfrac{0+16}{2} = \SI{8}{\micro\coulomb}$.\\
\textbf{Contato C--B:}
$Q_C = Q_B = \dfrac{8+4}{2} = \SI{6}{\micro\coulomb}$.\\
Situação final: $A=\SI{8}{\micro\coulomb}$, $B=Q_C=\SI{6}{\micro\coulomb}$.
Confira a conservação: $16+4+0 = 8+6+6 = \SI{20}{\micro\coulomb}$.}
\end{questao}

\begin{questao}[3]{}
Três esferas condutoras A, B e C têm o mesmo diâmetro. A esfera A está
inicialmente neutra, e as outras duas estão carregadas com
$q_B = \SI{6}{\micro\coulomb}$ e $q_C = \SI{7}{\micro\coulomb}$. Com a esfera A,
toca-se primeiro B e depois C. As cargas de A, B e C após os contatos são,
respectivamente:
\begin{alternativas}
\item zero, zero e $\SI{13}{\micro\coulomb}$
\item $\SI{7}{\micro\coulomb}$, $\SI{3}{\micro\coulomb}$ e $\SI{5}{\micro\coulomb}$
\item $\SI{5}{\micro\coulomb}$, $\SI{3}{\micro\coulomb}$ e $\SI{5}{\micro\coulomb}$
\item $\SI{6}{\micro\coulomb}$, $\SI{7}{\micro\coulomb}$ e zero
\item todas iguais a $\SI{4.3}{\micro\coulomb}$
\end{alternativas}
\resposta{(c)}
\resolucao{\textbf{Contato A--B:}
$q_A = q_B = \dfrac{0+6}{2} = \SI{3}{\micro\coulomb}$.\\
\textbf{Contato A--C:}
$q_A = q_C = \dfrac{3+7}{2} = \SI{5}{\micro\coulomb}$.\\
Final: $A=\SI{5}{\micro\coulomb}$, $B=\SI{3}{\micro\coulomb}$,
$C=\SI{5}{\micro\coulomb}$.
Verificação: $0+6+7 = 5+3+5 = \SI{13}{\micro\coulomb}$.}
\end{questao}

\begin{questao}[3]{}
Na eletrização por contato entre duas esferas condutoras de raios $R$ e $2R$,
sabe-se que, após o contato, a razão entre a carga adquirida por cada esfera e o
seu respectivo raio é constante. Se, inicialmente, a esfera de raio $R$ possui
carga $3Q$ e a de raio $2R$ possui $2Q$, calcule a carga final de cada esfera.
\resposta{$Q_1 = \dfrac{5Q}{3}$ (raio $R$) e $Q_2 = \dfrac{10Q}{3}$ (raio $2R$)}
\resolucao{Sejam $Q_1$ e $Q_2$ as cargas finais. Aplicando as duas condições:
\[
\textbf{(i) conservação: } Q_1+Q_2 = 3Q+2Q = 5Q,
\qquad
\textbf{(ii) proporcionalidade: } \frac{Q_1}{R}=\frac{Q_2}{2R}
\ \Rightarrow\ Q_2 = 2Q_1 .
\]
Substituindo (ii) em (i): $Q_1 + 2Q_1 = 5Q \Rightarrow Q_1 = \dfrac{5Q}{3}$ e
$Q_2 = \dfrac{10Q}{3}$.
Fisicamente, a proporcionalidade decorre de as esferas ficarem com o \emph{mesmo
potencial} após o contato: $V = \dfrac{\ke Q}{R}$ igual para ambas.}
\end{questao}

\begin{questao}[3]{}
Uma esfera metálica isolada possui carga inicial $Q_0$. Ela é colocada em contato,
uma única vez e sucessivamente, com $n$ esferas idênticas a ela e inicialmente
neutras.
\begin{itensq}
\item Obtenha a expressão da carga final da esfera original em função de $Q_0$ e $n$.
\item A partir de quantos contatos a carga da esfera original cai abaixo de
      \SI{1}{\percent} de $Q_0$?
\end{itensq}
\resposta{(a) $Q_n = Q_0/2^{\,n}$; \quad (b) $n = 7$}
\resolucao{(a) Cada contato com uma esfera idêntica e neutra reparte a carga em
duas partes iguais, logo $Q_1 = Q_0/2$, $Q_2 = Q_0/4$ e, por indução,
\[
\boxed{\,Q_n = \frac{Q_0}{2^{\,n}}\,}.
\]
(b) Queremos $\dfrac{1}{2^{\,n}} < 0{,}01 \Rightarrow 2^{\,n} > 100$. Como
$2^6 = 64$ e $2^7 = 128$, são necessários $n = 7$ contatos.
Repare que a carga nunca se anula: ela apenas tende a zero, o que ilustra por
que o aterramento (contato com um condutor "infinito") é a forma prática de
descarregar um corpo.}
\end{questao}

\begin{questao}[3]{}
Um gerador de Van de Graaff transporta carga para sua cúpula a uma taxa
constante de $\SI{1}{\micro\ampere}$.
\begin{itensq}
\item Que carga é acumulada em $\SI{10}{\second}$?
\item Quantos elétrons isso representa?
\item Se a cúpula é uma esfera de raio $\SI{25}{\centi\meter}$, qual o potencial
      atingido? O ar rompe a cerca de $\SI{3}{\mega\volt\per\meter}$ --- haverá
      descarga?
\end{itensq}
\dados{$e=\SI{1.6e-19}{\coulomb}$; $\ke=\SI{9e9}{\newton\meter\squared\per\coulomb\squared}$.}
\resposta{(a) \SI{10}{\micro\coulomb}; (b) $\pot{6.25}{13}$;
(c) $V=\SI{360}{\kilo\volt}$, $E=\SI{1.44}{\mega\volt\per\meter}$ --- sem descarga}
\resolucao{(a) $Q = I\,t = \pot{1}{-6}\cdot10 = \SI{10}{\micro\coulomb}$.\\
(b) $n = \dfrac{Q}{e} = \dfrac{\pot{1}{-5}}{\pot{1.6}{-19}} = \pot{6.25}{13}$
elétrons --- \emph{removidos} da cúpula (que fica positiva).\\
(c) Para uma esfera condutora,
\[
V=\frac{\ke Q}{R}=\frac{9\cdot10^9\cdot\pot{1}{-5}}{0{,}25}=\SI{360}{\kilo\volt},
\qquad
E_{\text{sup}}=\frac{\ke Q}{R^{2}}=\frac{V}{R}=\frac{\pot{3.6}{5}}{0{,}25}
=\SI{1.44}{\mega\volt\per\meter}.
\]
Como $\SI{1.44}{} < \SI{3}{\mega\volt\per\meter}$, \textbf{não há ruptura} --- a
carga continua se acumulando.\\
\textbf{Limite do gerador:} a descarga ocorreria quando
$E_{\text{sup}} = \SI{3}{\mega\volt\per\meter}$, ou seja
$V_{\max}=E\,R = 3\cdot10^6\cdot0{,}25=\SI{750}{\kilo\volt}$.
\textbf{Consequência de projeto:} como $V_{\max} \propto R$, cúpulas
\emph{maiores} sustentam tensões maiores --- por isso os geradores de demonstração
têm esferas polidas e volumosas, sem pontas (onde $E$ dispararia).}
\end{questao}

\begin{questao}[3]{}
Um fio de cobre tem massa $\SI{1}{\gram}$. Sabendo que cada átomo de cobre
contribui com \emph{um} elétron livre:
\begin{itensq}
\item Calcule o número de elétrons livres no fio.
\item Determine a carga total desses elétrons.
\item Compare com a carga típica obtida por atrito ($\sim\SI{1}{\micro\coulomb}$)
      e interprete o resultado.
\end{itensq}
\dados{$N_A=\SI{6.02e23}{\per\mole}$; $M_{\text{Cu}}=\SI{63.5}{\gram\per\mole}$;
$e=\SI{1.6e-19}{\coulomb}$.}
\resposta{(a) $\approx\pot{9.5}{21}$; (b) $\approx\SI{1517}{\coulomb}$;
(c) razão $\sim10^9$}
\resolucao{(a) O número de mols é $n=\dfrac{1}{63{,}5}$ mol, logo
\[
N=\frac{1}{63{,}5}\cdot\pot{6.02}{23}\approx\pot{9.5}{21}\ \text{elétrons livres}.
\]
(b) $|Q| = N\,e = \pot{9.5}{21}\cdot\pot{1.6}{-19}\approx\SI{1517}{\coulomb}$.\\
(c) A razão entre essa carga e a obtida por atrito é
\[
\frac{1517}{\pot{1}{-6}}\approx\pot{1.5}{9},
\]
ou seja, cerca de \textbf{um bilhão e meio de vezes} maior.

\textbf{Interpretação --- por que isso importa.} Eletrizar um corpo por atrito
remove ou acrescenta uma fração absolutamente ínfima dos elétrons disponíveis
(da ordem de $10^{-9}$). O material permanece, para todos os efeitos,
\emph{eletricamente neutro em sua estrutura}; a "carga" que percebemos é apenas
um desequilíbrio superficial minúsculo.\\
Isso explica dois fatos:
\begin{itemize}[leftmargin=1.6em,itemsep=1pt,topsep=2pt]
\item por que um condutor pode fornecer corrente indefinidamente sem "acabar" os
      elétrons;
\item por que forças elétricas macroscópicas são tão intensas --- se
      conseguíssemos separar \emph{toda} a carga de $\SI{1}{\gram}$ de cobre e
      colocá-la a $\SI{1}{\meter}$ de distância de outra igual, a força seria da
      ordem de $10^{16}$\,N, milhões de vezes o peso de um transatlântico.
\end{itemize}}
\end{questao}

\begin{questao}[3]{}
Três esferas condutoras isoladas têm raios $R$, $2R$ e $3R$. A primeira está
carregada com $Q$ e as demais estão neutras. Sabe-se que, ao se tocarem, duas
esferas de raios diferentes repartem a carga de modo que o \emph{potencial} de
ambas fique igual, o que implica $\dfrac{Q_1}{R_1}=\dfrac{Q_2}{R_2}$.

Encosta-se a esfera $A$ (raio $R$) na esfera $B$ (raio $2R$); separa-se; em
seguida, encosta-se $A$ na esfera $C$ (raio $3R$).
\begin{itensq}
\item Determine a carga final de cada esfera.
\item Verifique a conservação da carga total.
\item Generalize: qual a carga final de $A$ após tocar sucessivamente esferas de
      raios $2R, 3R, \dots, nR$?
\end{itensq}
\resposta{(a) $Q_A=\dfrac{Q}{12}$, $Q_B=\dfrac{2Q}{3}$, $Q_C=\dfrac{Q}{4}$;
(b) soma $=Q$; (c) $Q_A^{(n)}=Q\displaystyle\prod_{j=2}^{n}\frac{1}{j+1}$}
\resolucao{\textbf{Contato A--B.} A carga total $Q$ se reparte proporcionalmente
aos raios ($R$ e $2R$), ou seja, na razão $1:2$:
\[
Q_A=\frac{R}{R+2R}\,Q=\frac{Q}{3},
\qquad
Q_B=\frac{2R}{3R}\,Q=\frac{2Q}{3}.
\]
\textbf{Contato A--C.} Agora $A$ traz $\dfrac{Q}{3}$ e encontra $C$ (raio $3R$,
neutra). A repartição é na razão $R:3R = 1:3$:
\[
Q_A=\frac{1}{4}\cdot\frac{Q}{3}=\frac{Q}{12},
\qquad
Q_C=\frac{3}{4}\cdot\frac{Q}{3}=\frac{Q}{4}.
\]
\textbf{(b) Conservação:}
$\dfrac{Q}{12}+\dfrac{2Q}{3}+\dfrac{Q}{4}
=\dfrac{Q+8Q+3Q}{12}=\dfrac{12Q}{12}=Q$. \checkmark

\textbf{(c) Generalização.} Ao tocar uma esfera de raio $jR$, a esfera $A$ retém a
fração $\dfrac{R}{R+jR}=\dfrac{1}{1+j}$ da carga que possuía. Logo, após a
sequência $2R, 3R, \dots, nR$:
\[
Q_A^{(n)} = Q\prod_{j=2}^{n}\frac{1}{1+j}
= Q\cdot\frac{1}{3}\cdot\frac{1}{4}\cdots\frac{1}{n+1}
= \frac{2\,Q}{(n+1)!}.
\]
Confira para $n=3$: $\dfrac{2Q}{4!}=\dfrac{2Q}{24}=\dfrac{Q}{12}$. \checkmark\\
\textbf{Por que isso importa:} a carga de $A$ decai \emph{fatorialmente} --- muito
mais rápido que o decaimento $2^{-n}$ das esferas idênticas. É o princípio do
\emph{gerador de Van de Graaff}: transferir carga para um condutor de raio muito
maior esvazia quase completamente o menor.}
\end{questao}

% BEGIN BANCO COMPLEMENTAR Cargas elétricas
% =====================================================================
%  Banco complementar --- Cargas Eletricas  (REVISADO)
%  Substitui a versao gerada por template (2 clones em N3 e 8 questoes
%  conceituais marcadas como N4).
%  RESTRICAO SEVERA: o cap01 e extenso (31 questoes) e ja cobre
%  praticamente todos os temas das N4 antigas: contato entre esferas
%  identicas (varias vezes), contatos sucessivos, esferas de raios
%  diferentes, eletroscopio, inducao com aterramento, verificacao de
%  quantizacao (Q = 2,4e-19 C impossivel), serie triboeletrica,
%  gerador de Van de Graaff e contagem de eletrons em fio de cobre.
%  Este banco explora o que restou: limite de eletrizacao pela rigidez
%  dieletrica, poder das pontas, ordem de grandeza do atrito, forca
%  eletrica x peso, conservacao de carga em processos nucleares,
%  gaiola/copo de Faraday e a serie geometrica dos contatos sucessivos.
%  Distribuicao mantida: 2 N3 + 8 N4 = 10
% =====================================================================

\subsubsection*{Banco complementar --- Cargas elétricas}

\begin{atencao}[Organização do banco]
Dois princípios sustentam o tópico: a carga é \textbf{conservada} (em sistema
isolado, a soma algébrica não muda) e é \textbf{quantizada} (todo valor é
múltiplo inteiro de $e=\pot{1.6}{-19}\,\si{\coulomb}$). O N4 explora as
consequências quantitativas menos óbvias --- há um teto para a carga que um
corpo pode reter, o atrito envolve uma fração ínfima dos átomos, e a
distribuição em condutores é governada pela curvatura.
\end{atencao}

\blocoaprofundamento

\begin{questao}[3]{}
Um sistema isolado contém $5$ prótons e $3$ elétrons livres. Após uma
reorganização interna, dois elétrons passam a orbitar dois dos prótons,
formando átomos neutros.
\begin{itensq}
\item Determine a carga total antes e depois.
\item Justifique o resultado.
\end{itensq}
\dados{$e=\pot{1.6}{-19}\,\si{\coulomb}$}
\resposta{$Q=+2e=\pot{3.2}{-19}\,\si{\coulomb}$ nos dois instantes}
\resolucao{(a) \textbf{Antes:} $Q=5(+e)+3(-e)=+2e=\pot{3.2}{-19}\,\si{\coulomb}$.\\
\textbf{Depois:} dois pares próton--elétron formam átomos neutros, restando
$3$ prótons livres e $1$ elétron livre:
\[
Q=3(+e)+1(-e)=+2e=\pot{3.2}{-19}\,\si{\coulomb}.
\]
(b) \textbf{Nada mudou} --- e não poderia mudar. A reorganização apenas
\emph{redistribuiu} as cargas dentro do sistema; nenhuma foi criada nem
destruída, e o sistema é isolado.\\
\textbf{Ponto conceitual:} um corpo \emph{neutro} não é um corpo \emph{sem
carga} --- ele tem tantas cargas positivas quanto negativas. Agrupar cargas em
átomos neutros esconde-as, mas não as elimina do balanço.}
\end{questao}

\begin{questao}[3]{}
Um bastão isolante adquire, por atrito, carga de $\SI{-5}{\nano\coulomb}$.
\begin{itensq}
\item Determine o número de elétrons transferidos.
\item Determine a massa correspondente e compare com a massa do bastão
      ($\SI{20}{\gram}$).
\end{itensq}
\dados{$e=\pot{1.6}{-19}\,\si{\coulomb}$;
$m_e=\pot{9.11}{-31}\,\si{\kilogram}$}
\resposta{$n=\pot{3.1}{10}$ elétrons; $\Delta m=\pot{2.8}{-20}\,\si{\kilogram}$}
\resolucao{(a)
\[
n=\frac{|Q|}{e}=\frac{\pot{5}{-9}}{\pot{1.6}{-19}}=\pot{3.125}{10}.
\]
Trinta bilhões de elétrons --- número enorme em termos absolutos.\\
(b)
\[
\Delta m=n\,m_e=\pot{3.125}{10}\times\pot{9.11}{-31}
=\pot{2.85}{-20}\,\si{\kilogram}.
\]
Comparando com os $\SI{20}{\gram}$ do bastão:
\[
\frac{\Delta m}{m}=\frac{\pot{2.85}{-20}}{\pot{2}{-2}}=\pot{1.4}{-18}.
\]
\textbf{Conclusão:} a eletrização por atrito \emph{não} altera a massa de forma
mensurável --- o desvio está $15$ ordens de grandeza abaixo da melhor balança
analítica. É por isso que a massa não serve como indicador de eletrização, e
por que o fenômeno passou séculos sendo descrito apenas por seus efeitos de
força.}
\end{questao}

\blocodesafio

\begin{questao}[4]{}
\textbf{(Série geométrica dos contatos.)} Uma esfera $A$ com carga $Q$ é
tocada sucessivamente por esferas idênticas \emph{neutras}
$B_1,B_2,\dots,B_n$, uma de cada vez, sendo cada uma retirada após o contato.
\begin{itensq}
\item Determine a carga de $A$ após $n$ contatos.
\item Determine a carga de cada esfera $B_k$.
\item Verifique a conservação da carga total.
\end{itensq}
\resposta{$Q_A=Q/2^{n}$; $Q_{B_k}=Q/2^{k}$; a soma reconstitui $Q$}
\resolucao{(a) Cada contato entre esferas \emph{idênticas} reparte igualmente a
carga presente. Como cada $B_k$ chega neutra, $A$ perde metade do que tem:
\[
Q\to\frac{Q}{2}\to\frac{Q}{4}\to\cdots
\;\Longrightarrow\;
Q_A=\frac{Q}{2^{n}}.
\]
(b) A esfera $B_k$ leva metade do que $A$ possuía imediatamente antes, ou seja
$\dfrac{1}{2}\cdot\dfrac{Q}{2^{k-1}}=\dfrac{Q}{2^{k}}$.\\
(c) \textbf{Conservação:}
\[
Q_A+\sum_{k=1}^{n}Q_{B_k}
=\frac{Q}{2^{n}}+Q\sum_{k=1}^{n}\frac{1}{2^{k}}
=\frac{Q}{2^{n}}+Q\left(1-\frac{1}{2^{n}}\right)=Q.\ \checkmark
\]
A soma geométrica $\sum 1/2^{k}$ tende a $1$, e o termo residual de $A$ fecha
exatamente a diferença.\\
\textbf{Leitura:} nenhuma carga se perde --- ela apenas se dilui. Com
$n=10$, $A$ retém menos de $0{,}1\%$ de $Q$, mas os $99{,}9\%$ restantes estão
distribuídos pelas dez esferas retiradas.\\
\textbf{Uso prático:} é assim que se \emph{descarrega} progressivamente um corpo
sem aterrá-lo, e também como se obtêm frações controladas de uma carga de
referência.}
\end{questao}

\begin{questao}[4]{}
\textbf{(Limite de eletrização.)} Uma esfera condutora de raio
$R=\SI{10}{\centi\meter}$ é carregada no ar, cuja rigidez dielétrica é
$E_{rup}=\SI{3e6}{\volt\per\meter}$.
\begin{itensq}
\item Determine a carga máxima que ela pode reter.
\item Determine o potencial correspondente.
\item Explique o que ocorre se tentarmos ultrapassar esse valor.
\end{itensq}
\dados{$\ke=\SI{9e9}{\newton\meter\squared\per\coulomb\squared}$}
\resposta{$Q_{\max}=\SI{3.33}{\micro\coulomb}$; $V=\SI{300}{\kilo\volt}$}
\resolucao{(a) O campo na superfície de uma esfera condutora vale
$E=\ke Q/R^{2}$. Impondo $E=E_{rup}$:
\[
Q_{\max}=\frac{E_{rup}R^{2}}{\ke}=\frac{\pot{3}{6}\times0{,}01}{\pot{9}{9}}
=\pot{3.33}{-6}\,\si{\coulomb}=\SI{3.33}{\micro\coulomb}.
\]
(b)
\[
V=\frac{\ke Q_{\max}}{R}=E_{rup}\,R=\pot{3}{6}(0{,}10)=\SI{300}{\kilo\volt}.
\]
\textbf{Relação notável:} $V_{\max}=E_{rup}R$ --- o potencial máximo é
\emph{proporcional ao raio}. Dobrar o raio dobra a tensão suportável, e é por
isso que geradores de Van de Graaff têm cúpulas grandes e lisas.\\
(c) \textbf{Acima do limite}, o ar ioniza e passa a conduzir: a carga escapa por
descarga corona (o brilho azulado e o chiado característicos) ou por faísca.
A esfera \emph{não consegue} reter mais carga --- ela se autolimita.\\
\textbf{Consequência de projeto:} para armazenar mais carga é preciso aumentar
$R$, usar meio de maior rigidez (óleo, $\mathrm{SF_6}$, vácuo) ou eliminar
irregularidades da superfície, onde o campo local é muito maior --- assunto da
questão seguinte.}
\end{questao}

\begin{questao}[4]{}
\textbf{(Poder das pontas.)} Duas esferas condutoras de raios
$R=\SI{10}{\centi\meter}$ e $r=\SI{1}{\centi\meter}$ são ligadas por um fio
longo e carregadas até $V=\SI{30}{\kilo\volt}$.
\begin{itensq}
\item Determine a razão entre as cargas e entre as densidades superficiais.
\item Determine o campo na superfície de cada uma.
\item Explique o efeito e cite duas aplicações.
\end{itensq}
\resposta{$Q_R/Q_r=10$; $\sigma_r/\sigma_R=10$;
$E_R=\SI{300}{\kilo\volt\per\meter}$ e $E_r=\SI{3}{\mega\volt\per\meter}$}
\resolucao{(a) Ligadas por fio, as duas ficam no \emph{mesmo potencial}:
\[
V=\frac{\ke Q}{R}\Rightarrow Q\propto R
\;\Longrightarrow\;
\frac{Q_R}{Q_r}=\frac{R}{r}=10.
\]
A densidade superficial é $\sigma=Q/(4\pi R^{2})\propto R/R^{2}=1/R$:
\[
\frac{\sigma_r}{\sigma_R}=\frac{R}{r}=10.
\]
\textbf{A esfera maior tem mais carga, mas a menor tem carga mais
concentrada.}\\
(b) $E=V/R$ na superfície de cada uma:
\[
E_R=\frac{\pot{30}{3}}{0{,}10}=\SI{300}{\kilo\volt\per\meter};
\qquad
E_r=\frac{\pot{30}{3}}{0{,}01}=\SI{3}{\mega\volt\per\meter}.
\]
(c) \textbf{O campo na esfera pequena atingiu exatamente a rigidez do ar.} A
ionização começa \emph{ali}, embora as duas estejam ao mesmo potencial ---
regiões de menor raio de curvatura concentram campo. É o \textbf{poder das
pontas}.\\
\textbf{Aplicações:}
\begin{itemize}
\item \textbf{Para-raios:} a ponta afiada ioniza o ar e estabelece o caminho
preferencial da descarga, protegendo a estrutura.
\item \textbf{Impressão e pintura eletrostática:} eletrodos pontiagudos ionizam
o ar para carregar tinta ou papel.
\end{itemize}
\textbf{E a contrapartida:} em equipamentos de alta tensão evita-se qualquer
ponta ou aresta viva --- daí os terminais arredondados e as blindagens
toroidais.}
\end{questao}

\begin{questao}[4]{}
\textbf{(Ordem de grandeza do atrito.)} Um bastão de vidro de $\SI{20}{\gram}$
adquire $\SI{5}{\nano\coulomb}$ por atrito. A massa molar média do vidro é
cerca de $\SI{60}{\gram\per\mole}$.
\begin{itensq}
\item Estime o número de átomos do bastão.
\item Determine a fração de átomos que perdeu um elétron.
\item Interprete o resultado.
\end{itensq}
\dados{$N_A=\pot{6.02}{23}\,\si{\per\mole}$}
\resposta{$\approx\pot{2}{23}$ átomos; fração de $\pot{1.6}{-13}$}
\resolucao{(a)
\[
N=\frac{m}{M}N_A=\frac{20}{60}\times\pot{6.02}{23}\approx\pot{2.0}{23}.
\]
(b) Já se sabe que $\pot{3.1}{10}$ elétrons foram transferidos:
\[
f=\frac{\pot{3.1}{10}}{\pot{2.0}{23}}=\pot{1.6}{-13},
\]
ou cerca de \textbf{um átomo em seis trilhões}.\\
(c) \textbf{Duas leituras que se complementam.}\\
\emph{Do ponto de vista microscópico}, a eletrização por atrito é um fenômeno
de superfície absolutamente marginal --- envolve uma fração desprezível da
matéria, e apenas a camada atômica mais externa.\\
\emph{Do ponto de vista macroscópico}, esses mesmos
$\pot{3.1}{10}$ elétrons produzem forças facilmente visíveis, capazes de
levantar pedaços de papel contra o peso.\\
\textbf{O contraste explica a intensidade da força elétrica:} ela é tão maior
que a gravitacional que um desequilíbrio de uma parte em $10^{13}$ já produz
efeitos observáveis. E explica também por que a matéria comum é neutra com
altíssima precisão --- qualquer desequilíbrio maior seria violentamente
desfeito.}
\end{questao}

\begin{questao}[4]{}
\textbf{(Força elétrica contra peso.)} Duas esferas idênticas de
$\SI{1}{\gram}$ estão a $\SI{5}{\centi\meter}$ uma da outra, cada uma com
carga $q$.
\begin{itensq}
\item Determine a força para $q=\SI{5}{\nano\coulomb}$ e para
      $q=\SI{50}{\nano\coulomb}$.
\item Determine a carga necessária para que a força iguale o peso.
\item Comente a viabilidade.
\end{itensq}
\dados{$g=\SI{9.8}{\meter\per\second\squared}$}
\resposta{$\SI{90}{\micro\newton}$ e $\SI{9}{\milli\newton}$;
$q\approx\SI{52}{\nano\coulomb}$}
\resolucao{(a)
\[
F=\frac{\ke q^{2}}{d^{2}}:\quad
q=\SI{5}{\nano\coulomb}\Rightarrow\SI{9.0e-5}{\newton};
\qquad
q=\SI{50}{\nano\coulomb}\Rightarrow\SI{9.0e-3}{\newton}.
\]
O peso vale $P=mg=\pot{1}{-3}(9{,}8)=\SI{9.8}{\milli\newton}$.\\
(b) Impondo $F=P$:
\[
q=d\sqrt{\frac{P}{\ke}}=0{,}05\sqrt{\frac{\pot{9.8}{-3}}{\pot{9}{9}}}
=0{,}05\times\pot{1.04}{-6}=\SI{52}{\nano\coulomb}.
\]
(c) \textbf{Perfeitamente viável.} Cinquenta nanocoulombs correspondem a
$\pot{3.3}{11}$ elétrons --- ordem de grandeza obtida rotineiramente por
atrito, como mostrou a questão anterior.\\
\textbf{É por isso que a eletrostática é observável a olho nu:} um simples
canudo atritado sustenta pedaços de papel contra a gravidade de todo o planeta.
\\
\textbf{Verificação do limite:} essas esferas suportariam essa carga? Para
$R=\SI{5}{\milli\meter}$, o campo na superfície seria
$\ke q/R^{2}=\SI{1.9}{\mega\volt\per\meter}$ --- abaixo da rigidez do ar, mas
já na mesma ordem. Cargas muito maiores escapariam por descarga.}
\end{questao}

\begin{questao}[4]{}
\textbf{(Conservação em processos nucleares.)} Verifique a conservação da carga
nos processos abaixo.
\begin{itensq}
\item Decaimento beta: $n\to p+e^{-}+\bar\nu_e$.
\item Aniquilação: $e^{-}+e^{+}\to2\gamma$.
\item Discuta o alcance do princípio.
\end{itensq}
\resposta{Conservada em ambos: $0=+e-e+0$ e $-e+e=0$}
\resolucao{(a) \textbf{Decaimento beta.} Antes: nêutron, carga $0$. Depois:
próton ($+e$), elétron ($-e$) e antineutrino ($0$):
\[
0=(+e)+(-e)+0=0.\ \checkmark
\]
Note que \emph{o número de partículas mudou} e que partículas foram
\emph{criadas} --- o elétron não estava dentro do nêutron. Ainda assim a carga
total se conserva.\\
(b) \textbf{Aniquilação.} Antes: $-e+(+e)=0$. Depois: dois fótons, ambos de
carga nula. \checkmark\\
Aqui a \emph{matéria} desaparece por completo, convertida em radiação --- e a
carga continua conservada, porque já era nula.\\
(c) \textbf{Alcance do princípio.} A conservação da carga é uma das leis mais
robustas da física:
\begin{itemize}
\item Vale em processos \textbf{químicos}, \textbf{nucleares} e de
\textbf{física de partículas}, sem exceção conhecida.
\item Vale mesmo quando partículas são criadas ou destruídas --- o que a
distingue da conservação do número de partículas, que \emph{não} é uma lei.
\item Experimentalmente, o decaimento do elétron (que violaria a conservação) foi
buscado e não observado; o limite atual para sua vida média é superior a
$10^{26}$ anos.
\end{itemize}
\textbf{Origem teórica:} a conservação da carga decorre de uma simetria
fundamental do eletromagnetismo, do mesmo modo que a conservação da energia
decorre da invariância no tempo. Ela não é um fato empírico isolado, mas
consequência da estrutura da teoria.}
\end{questao}

\begin{questao}[4]{}
\textbf{(Copo de Faraday.)} Um corpo carregado é introduzido, sem tocar as
paredes, no interior de um recipiente metálico fechado ligado a um eletrômetro.
\begin{itensq}
\item Explique o que o instrumento indica e por quê.
\item Explique por que a leitura não depende da posição do corpo dentro do
      recipiente.
\item Descreva o que ocorre se o corpo tocar a parede interna.
\end{itensq}
\resposta{O eletrômetro indica exatamente a carga do corpo, independentemente da
posição}
\resolucao{(a) \textbf{Indução total.} O corpo de carga $+Q$ induz $-Q$ na
superfície \emph{interna} do recipiente --- exatamente $-Q$, porque todas as
linhas de campo que saem do corpo terminam nela. Por conservação, sobra
$+Q$ na superfície \emph{externa}, e é essa carga que o eletrômetro mede.\\
\textbf{Resultado:} a leitura é $+Q$, o valor exato da carga introduzida.\\
(b) \textbf{Independência da posição.} O campo elétrico dentro do metal é nulo
em equilíbrio. Aplicando a lei de Gauss a uma superfície inteiramente contida no
metal, o fluxo é zero --- logo a carga interna total (corpo mais superfície
interna) é nula, \emph{qualquer que seja a posição do corpo}. A carga externa
fica sempre $+Q$, distribuída de acordo apenas com a geometria \emph{externa}
do recipiente.\\
\textbf{Consequência prática:} não é preciso centralizar o corpo nem controlar
sua orientação --- o que torna o copo de Faraday o método padrão de medição
absoluta de carga.\\
(c) \textbf{Se o corpo tocar a parede:} a carga $+Q$ neutraliza a
$-Q$ induzida, o corpo sai neutro e o recipiente fica com $+Q$ na superfície
externa. \textbf{A leitura do eletrômetro não muda} --- continua $+Q$.\\
Isso permite \emph{transferir} integralmente a carga de um corpo para o
recipiente, e é a base do método clássico de acumulação sucessiva de carga em
geradores eletrostáticos.}
\end{questao}

\begin{questao}[4]{}
\textbf{(Blindagem eletrostática.)} Uma malha metálica aterrada envolve um
instrumento sensível.
\begin{itensq}
\item Explique por que o campo externo não penetra.
\item Determine se a blindagem funciona também no sentido inverso.
\item Discuta a diferença entre malha aterrada e não aterrada.
\end{itensq}
\resposta{Blinda o interior contra campos externos; o aterramento é necessário
para blindar o exterior contra fontes internas}
\resolucao{(a) \textbf{Campo externo não penetra.} As cargas livres do condutor
se redistribuem até anular o campo em seu interior --- é a condição de
equilíbrio eletrostático. Qualquer campo externo induz uma distribuição
superficial cujo campo cancela exatamente o campo aplicado na região interna.\\
Isso vale \emph{mesmo sem aterramento}: a gaiola de Faraday isolada blinda o
interior. É o efeito que protege passageiros de um carro atingido por raio.\\
(b) \textbf{Sentido inverso: só com aterramento.} Uma carga $+Q$ colocada
\emph{dentro} da gaiola induz $-Q$ na face interna e $+Q$ na externa --- e essa
carga externa produz campo do lado de fora. \textbf{A blindagem falha.}\\
Com a malha \textbf{aterrada}, a carga $+Q$ da face externa escoa para a terra,
e o campo externo desaparece. Só então a blindagem funciona nos dois sentidos.\\
(c) \textbf{Resumo da diferença.}
\begin{center}
\begin{tabular}{lcc}
\hline
& não aterrada & aterrada\\
\hline
Protege o interior de campos externos & sim & sim\\
Impede que fontes internas afetem o exterior & não & sim\\
\hline
\end{tabular}
\end{center}
\textbf{Limitações práticas:} a malha só é eficaz para aberturas muito menores
que o comprimento de onda do sinal a blindar --- em eletrostática (campo
constante), qualquer malha fina basta; em radiofrequência, a exigência é
severa. E a impedância da conexão de terra importa: um fio longo de aterramento
tem indutância e degrada a blindagem em altas frequências.}
\end{questao}

% END BANCO COMPLEMENTAR

\gabarito[1.1 Cargas elétricas]

% ---------------------------------------------------------------
% Subsecoes seguintes do Capitulo 1
% =====================================================================
%  CAPITULO 1 - continuacao
%  1.2 Lei de Coulomb  |  1.3 Campo eletrico
% =====================================================================

% =====================================================================
\subsection{Lei de Coulomb}
% =====================================================================

\begin{conceito}[Antes de começar]
\textbf{Lei de Coulomb} --- o módulo da força entre duas cargas puntiformes é
\[
F = \ke\,\frac{|q_1||q_2|}{d^2},\qquad
\ke = \SI{9e9}{\newton\meter\squared\per\coulomb\squared}\ \text{(no vácuo)}.
\]
A força é \emph{atrativa} para cargas de sinais opostos e \emph{repulsiva} para
sinais iguais, sempre na direção da reta que une as cargas.\par
\textbf{Princípio da superposição:} a força resultante sobre uma carga é a
\emph{soma vetorial} das forças que cada uma das demais exerce sobre ela.\par
\textbf{Dependências úteis:} $F \propto \dfrac{1}{d^2}$ (dobrar a distância divide
a força por 4) e $F \propto q_1 q_2$.
\end{conceito}

\legendaniveis

\blocofixacao

\begin{questao}[1]{}
Duas cargas $q_1 = \SI{+2}{\micro\coulomb}$ e $q_2 = \SI{-3}{\micro\coulomb}$
estão separadas por $\SI{4}{\centi\meter}$ no vácuo. Determine o módulo e o tipo
(atração ou repulsão) da força entre elas.
\dados{$\ke = \SI{9e9}{\newton\meter\squared\per\coulomb\squared}$.}
\resposta{$F = \SI{33.75}{\newton}$, atrativa}
\resolucao{Com $d = \SI{4}{\centi\meter} = \SI{0.04}{\meter}$:
\[
F=\ke\frac{|q_1||q_2|}{d^2}
=9\cdot10^9\cdot\frac{(\pot{2}{-6})(\pot{3}{-6})}{(0{,}04)^2}
=\SI{33.75}{\newton}.
\]
Como as cargas têm sinais opostos, a força é \textbf{atrativa}.}
\end{questao}

\begin{questao}[1]{}
Duas cargas $q_1 = q_2 = \SI{+5}{\micro\coulomb}$ estão a
$\SI{2}{\centi\meter}$ no vácuo. Determine o módulo e o tipo da força.
\dados{$\ke = \SI{9e9}{\newton\meter\squared\per\coulomb\squared}$.}
\resposta{$F = \SI{562.5}{\newton}$, repulsiva}
\resolucao{$F = 9\cdot10^9\cdot\dfrac{(\pot{5}{-6})^2}{(0{,}02)^2}
=\SI{562.5}{\newton}$, \textbf{repulsiva} (sinais iguais).}
\end{questao}

\begin{questao}[1]{}
Duas cargas puntiformes se repelem com força de intensidade
$\pot{42}{-5}\,\si{\newton}$. Reduzindo-se a distância entre elas à
\textbf{metade}, a nova força de repulsão será:
\begin{alternativas}[2]
\item $\pot{10.5}{-5}\,\si{\newton}$
\item $\pot{21}{-5}\,\si{\newton}$
\item $\pot{84}{-5}\,\si{\newton}$
\item $\pot{168}{-5}\,\si{\newton}$
\item inalterada
\end{alternativas}
\resposta{(d)}
\resolucao{Como $F \propto \dfrac{1}{d^2}$, reduzir $d$ à metade \emph{multiplica}
a força por 4:
\[
F' = 4F = 4\cdot\pot{42}{-5} = \pot{168}{-5}\,\si{\newton}.
\]}
\end{questao}

\blocoaplicacao

\begin{questao}[2]{}
A força elétrica entre duas partículas de cargas $q$ e $q/2$, separadas por uma
distância $d$ no vácuo, é $F$. A força entre duas partículas de cargas $q$ e $2q$,
separadas por $d/2$, também no vácuo, é:
\begin{alternativas}[3]
\item $F$
\item $2F$
\item $4F$
\item $8F$
\item $16F$
\end{alternativas}
\resposta{(e)}
\resolucao{Situação inicial: $F = \ke\dfrac{q\cdot(q/2)}{d^2} = \dfrac{\ke q^2}{2d^2}$.\\
Situação nova: $F' = \ke\dfrac{q\cdot 2q}{(d/2)^2} = \dfrac{2\ke q^2}{d^2/4}
= \dfrac{8\ke q^2}{d^2}$.\\
Dividindo: $\dfrac{F'}{F} = \dfrac{8\ke q^2/d^2}{\ke q^2/(2d^2)} = 16$, logo
$F' = 16F$. O produto das cargas quadruplicou e a distância caiu à metade
(fator 4 na força): $4\times4 = 16$.}
\end{questao}

\begin{questao}[2]{}
Se a força entre dois objetos carregados se mantém \emph{constante} mesmo quando a
carga de cada um é reduzida à metade, então a distância entre eles:
\begin{alternativas}
\item foi quadruplicada.
\item foi duplicada.
\item foi reduzida à quarta parte.
\item foi reduzida à metade.
\item permaneceu constante.
\end{alternativas}
\resposta{(d)}
\resolucao{Reduzir cada carga à metade divide o produto $q_1 q_2$ por 4, o que
sozinho reduziria $F$ a $1/4$. Para $F$ permanecer constante, o fator $1/d^2$
precisa \emph{quadruplicar}, ou seja, $d^2$ deve cair a $1/4$, o que significa $d$
reduzido à \textbf{metade}:
\[
F=\ke\frac{(q_1/2)(q_2/2)}{d'^2}=\ke\frac{q_1q_2}{4d'^2}
\stackrel{!}{=}\ke\frac{q_1q_2}{d^2}
\ \Rightarrow\ d'^2=\frac{d^2}{4}\ \Rightarrow\ d'=\frac{d}{2}.
\]}
\end{questao}

\begin{questao}[2]{}
Três cargas iguais de $\SI{+1}{\micro\coulomb}$ ocupam os vértices de um triângulo
equilátero de lado $\SI{5}{\centi\meter}$. Determine o módulo da força resultante
sobre uma delas.

\circuitoq{%
\begin{tikzpicture}[scale=1,line width=0.8pt,>=Stealth]
 \coordinate (A) at (0,0);
 \coordinate (B) at (3,0);
 \coordinate (C) at (1.5,2.6);
 \foreach \p/\n in {A/A,B/B,C/C}{
   \shade[ball color=Vcol!25] (\p) circle (0.28);
   \draw[Vcol] (\p) circle (0.28);
   \node at (\p) {\footnotesize$+q$};}
 \draw[dashed,cinza] (A)--(B)--(C)--cycle;
 \node[below] at (1.5,-0.35) {$\SI{5}{\centi\meter}$};
\end{tikzpicture}}{Três cargas iguais nos vértices de um triângulo equilátero.}

\resposta{$F_R \approx \SI{6.24}{\newton}$}
\resolucao{Cada par de cargas exerce uma força de módulo
\[
F=\ke\frac{q^2}{d^2}=9\cdot10^9\cdot\frac{(\pot{1}{-6})^2}{(0{,}05)^2}
=\SI{3.6}{\newton}.
\]
Sobre a carga escolhida atuam \emph{duas} forças iguais de \SI{3.6}{\newton}, com
$\SI{60}{\degree}$ entre elas (geometria do equilátero). A resultante é
\[
F_R = \sqrt{F^2+F^2+2F^2\cos 60^\circ}=F\sqrt{3}
=3{,}6\sqrt{3}\approx\SI{6.24}{\newton},
\]
dirigida para fora do triângulo, ao longo da bissetriz do vértice.}
\end{questao}

\begin{questao}[2]{}
Quatro cargas positivas iguais de $\SI{1}{\micro\coulomb}$ ocupam os vértices de
um quadrado de lado $\SI{10}{\centi\meter}$. Determine o módulo da força
resultante sobre uma delas.
\dados{$\ke = \SI{9e9}{\newton\meter\squared\per\coulomb\squared}$; $\sqrt2\approx1{,}41$.}

\circuitoq{%
\begin{tikzpicture}[scale=1,line width=0.8pt,>=Stealth]
 \foreach \p in {(0,0),(2.4,0),(2.4,2.4),(0,2.4)}{
   \shade[ball color=Vcol!25] \p circle (0.22);}
 \node at (0,0){\footnotesize$+q$};\node at (2.4,0){\footnotesize$+q$};
 \node at (2.4,2.4){\footnotesize$+q$};\node at (0,2.4){\footnotesize$+q$};
 \draw[dashed,cinza] (0,0)--(2.4,0)--(2.4,2.4)--(0,2.4)--cycle;
 \draw[dashed,cinza] (0,0)--(2.4,2.4);
 \node[below] at (1.2,-0.3){$\SI{10}{\centi\meter}$};
\end{tikzpicture}}{Quatro cargas iguais nos vértices de um quadrado.}

\resposta{$F_R \approx \SI{1.72}{\newton}$}
\resolucao{Sobre a carga de um vértice atuam três forças repulsivas:
duas dos vértices \emph{adjacentes} (ao longo dos lados, perpendiculares entre si)
e uma do vértice \emph{oposto} (ao longo da diagonal).\\
\textbf{Força de um lado} ($d = \SI{10}{\centi\meter}$):
$F_L = \ke\dfrac{q^2}{d^2} = 9\cdot10^9\dfrac{(\pot{1}{-6})^2}{(0{,}1)^2}
= \SI{0.9}{\newton}$.\\
\textbf{Força da diagonal} ($d' = d\sqrt2$):
$F_D = \ke\dfrac{q^2}{(d\sqrt2)^2} = \dfrac{F_L}{2} = \SI{0.45}{\newton}$.\\
As duas forças de lado, perpendiculares, têm resultante
$F_L\sqrt2 \approx \SI{1.27}{\newton}$, na direção da diagonal --- \emph{mesma}
direção de $F_D$. Somando:
\[
F_R = F_L\sqrt2 + F_D \approx 1{,}27 + 0{,}45 = \SI{1.72}{\newton}.
\]}
\end{questao}

\begin{questao}[2]{}
Uma pequena esfera de massa $\SI{2}{\gram}$, eletrizada com carga $q$, permanece
em equilíbrio, suspensa, sob a ação combinada do peso e de um campo elétrico
uniforme vertical de intensidade $\SI{e5}{\newton\per\coulomb}$, dirigido para
cima. Determine a carga $q$.
\dados{$g = \SI{10}{\meter\per\second\squared}$.}
\resposta{$q = \SI{0.2}{\micro\coulomb}$}
\resolucao{No equilíbrio, a força elétrica (para cima) equilibra o peso (para
baixo): $qE = mg$, logo
\[
q=\frac{mg}{E}=\frac{\pot{2}{-3}\cdot10}{\pot{1}{5}}
=\pot{2}{-7}\,\si{\coulomb}=\SI{0.2}{\micro\coulomb}.
\]
Para a força ser para cima (mesmo sentido de $E$), a carga deve ser
\emph{positiva}. É o princípio do experimento de Millikan, que mediu a carga do
elétron equilibrando gotas de óleo.}
\end{questao}

\blocoaprofundamento

\begin{questao}[3]{}
Duas cargas $Q_1$ e $Q_2 = 4Q_1$ (ambas positivas) estão fixas nos pontos A e B,
distantes $\SI{30}{\centi\meter}$. Em que posição $x$, medida a partir de A sobre
a reta AB, deve-se colocar uma carga $Q_3$ para que ela fique em equilíbrio sob
ação apenas de forças elétricas?

\circuitoq{%
\begin{tikzpicture}[scale=1,line width=0.8pt]
 \draw[cinza] (0,0)--(6,0);
 \shade[ball color=Vcol!25] (0,0) circle (0.25);
 \shade[ball color=Vcol!25] (6,0) circle (0.25);
 \node at (0,0) {\footnotesize$Q_1$}; \node at (6,0) {\footnotesize$Q_2$};
 \node[below] at (0,-0.35) {A}; \node[below] at (6,-0.35) {B};
 \fill[laranja] (2,0) circle (0.09);
 \node[above,laranja] at (2,0.15) {$Q_3$};
 \draw[<->,cinza] (0,-0.9)--(2,-0.9) node[midway,below]{$x$};
 \draw[<->,cinza] (0,-1.5)--(6,-1.5) node[midway,below]{$\SI{30}{\centi\meter}$};
\end{tikzpicture}}{Posição de equilíbrio entre duas cargas.}

\resposta{$x = \SI{10}{\centi\meter}$ a partir de A}
\resolucao{No equilíbrio, as forças que $Q_1$ e $Q_2$ exercem sobre $Q_3$ têm
mesmo módulo e sentidos opostos (o que só é possível \emph{entre} as cargas, já
que ambas são positivas). Igualando os módulos, com $d = \SI{30}{\centi\meter}$:
\[
\ke\frac{Q_1 Q_3}{x^2}=\ke\frac{4Q_1\,Q_3}{(d-x)^2}
\ \Rightarrow\
\frac{1}{x^2}=\frac{4}{(d-x)^2}
\ \Rightarrow\
\frac{d-x}{x}=2.
\]
Logo $d - x = 2x \Rightarrow x = \dfrac{d}{3} = \SI{10}{\centi\meter}$. O ponto de
equilíbrio fica \emph{mais próximo} da carga menor ($Q_1$), como esperado.
Note que o valor de $Q_3$ e seu sinal não influenciam a posição.}
\end{questao}

\begin{questao}[3]{}
Quatro cargas ocupam os vértices de um quadrado: duas delas são $+Q$ e $-Q$ (mesmo
módulo), e as outras duas, $q_1$ e $q_2$, são desconhecidas. Coloca-se uma carga
de prova positiva no centro e observa-se que a força resultante sobre ela aponta
na direção da diagonal, \emph{do vértice $-Q$ para o vértice $+Q$}. Sobre $q_1$ e
$q_2$ (nos outros dois vértices), pode-se afirmar que, para que essa força
resultante fique exatamente sobre a diagonal $-Q\,$--$\,+Q$:
\begin{alternativas}
\item devem ser iguais em módulo e de mesmo sinal entre si.
\item devem ser iguais em módulo e de sinais opostos.
\item devem ser ambas nulas.
\item devem ser ambas positivas e diferentes.
\item nada se pode afirmar.
\end{alternativas}
\resposta{(a)}
\resolucao{A força devida ao par $+Q$/$-Q$ já está sobre a diagonal
$-Q\,$--$\,+Q$. Para que a resultante \emph{total} continue sobre essa mesma
diagonal, a contribuição do par $q_1$/$q_2$ (na outra diagonal) não pode ter
componente perpendicular a ela: as forças de $q_1$ e $q_2$ sobre a carga central
devem se \emph{cancelar} nessa direção transversal. Isso ocorre quando $q_1$ e
$q_2$ têm o mesmo módulo e o mesmo sinal (por simetria em relação à diagonal
principal), produzindo forças opostas sobre a carga de prova que se anulam.}
\end{questao}

\begin{questao}[3]{}
Duas cargas fixas, $q_1 = \SI{+9}{\micro\coulomb}$ em $x=0$ e
$q_2 = \SI{+4}{\micro\coulomb}$ em $x = \SI{10}{\centi\meter}$, estão sobre o eixo
$x$. Uma terceira carga $q_3$ é colocada sobre esse eixo.
\begin{itensq}
\item Determine a posição em que $q_3$ fica em equilíbrio.
\item Mostre que existe uma \emph{segunda} raiz matemática e explique por que ela
      deve ser descartada.
\item O equilíbrio encontrado é \emph{estável} ou \emph{instável}? Analise para
      $q_3$ positiva e negativa.
\end{itensq}
\resposta{(a) $x = \SI{6}{\centi\meter}$; (b) $x=\SI{30}{\centi\meter}$ (fora do
segmento); (c) estável para $q_3<0$, instável para $q_3>0$}
\resolucao{(a) Igualando os módulos das forças (a carga $q_3$ cancela-se):
\[
\frac{\ke\,9}{x^{2}}=\frac{\ke\,4}{(10-x)^{2}}
\;\Longrightarrow\;
\frac{3}{x}=\frac{2}{10-x}
\;\Longrightarrow\;
30-3x=2x
\;\Longrightarrow\;
x=\SI{6}{\centi\meter}.
\]
(b) Ao elevar ao quadrado, a equação admite também $\dfrac{3}{x}=-\dfrac{2}{10-x}$,
que fornece $x = \SI{30}{\centi\meter}$. Nesse ponto, \emph{fora} do segmento,
as duas forças apontam no \textbf{mesmo sentido} (ambas as cargas são positivas e
empurram $q_3$ para longe) --- elas se somam e nunca se cancelam. A raiz é
espúria, introduzida pela elevação ao quadrado.

(c) \textbf{Análise de estabilidade.} Suponha $q_3$ deslocada ligeiramente para a
direita (aproximando-se de $q_2$):
\begin{itemize}[leftmargin=1.6em,itemsep=1pt,topsep=2pt]
\item Se $q_3 > 0$: a repulsão de $q_2$ (agora mais próxima) supera a de $q_1$, e
      $q_3$ é empurrada \emph{ainda mais} para longe do ponto de equilíbrio ---
      \textbf{instável}.
\item Se $q_3 < 0$: a atração por $q_2$ cresce, puxando $q_3$ \emph{de volta}?
      Não --- a atração por $q_2$ a puxa para longe de $x=6$. Nesse eixo,
      portanto, também é instável.
\end{itemize}
\textbf{Análise completa:} para $q_3<0$, o equilíbrio é estável na direção
\emph{transversal} ao eixo (qualquer deslocamento lateral gera força restauradora),
mas instável ao longo do eixo. Para $q_3>0$, ocorre o oposto.\\
\textbf{Conclusão profunda --- teorema de Earnshaw:} \emph{nenhuma} configuração
de cargas estáticas produz equilíbrio estável em \textbf{todas} as direções. Sempre
há ao menos uma direção instável. É por isso que não se consegue levitar um objeto
apenas com cargas fixas --- são necessários campos variáveis (armadilhas de Paul)
ou diamagnetismo.}
\end{questao}

\begin{questao}[3]{}
Compare a força elétrica e a força gravitacional entre dois elétrons.
\begin{itensq}
\item Mostre que a razão $F_e/F_g$ \emph{independe} da distância.
\item Calcule seu valor numérico.
\item Que conclusão isso permite sobre a estrutura da matéria e sobre a
      astronomia?
\end{itensq}
\dados{$\ke=\SI{9e9}{}$; $G=\SI{6.67e-11}{}$; $e=\SI{1.6e-19}{\coulomb}$;
$m_e=\SI{9.11e-31}{\kilogram}$ (SI).}
\resposta{(a) demonstração; (b) $\approx\pot{4.2}{42}$; (c) ver comentário}
\resolucao{(a) Ambas as forças seguem a lei do inverso do quadrado:
\[
F_e=\ke\frac{e^{2}}{d^{2}},
\qquad
F_g=G\frac{m_e^{2}}{d^{2}}
\;\Longrightarrow\;
\frac{F_e}{F_g}=\frac{\ke\,e^{2}}{G\,m_e^{2}}.
\]
O fator $d^{2}$ \textbf{cancela-se}: a razão é uma constante universal.\\
(b) Numericamente:
\[
\frac{F_e}{F_g}=\frac{(9\cdot10^{9})(\pot{1.6}{-19})^{2}}
{(\pot{6.67}{-11})(\pot{9.11}{-31})^{2}}\approx\pot{4.2}{42}.
\]
(c) \textbf{Duas conclusões profundas:}
\begin{itemize}[leftmargin=1.6em,itemsep=2pt,topsep=3pt]
\item \textbf{Na escala atômica, a gravidade é irrelevante.} A força elétrica é
      $10^{42}$ vezes maior --- por isso a estrutura de átomos, moléculas,
      cristais e de toda a química é governada \emph{exclusivamente} por forças
      eletromagnéticas. Nenhum modelo atômico precisa incluir gravidade.
\item \textbf{Na escala astronômica, ocorre o oposto.} Corpos celestes são
      eletricamente \emph{neutros} (as cargas positivas e negativas se cancelam
      quase perfeitamente), então a força elétrica resultante é praticamente
      nula. Já a massa só se \emph{acumula} --- não existe "massa negativa". Por
      isso a gravidade, apesar de absurdamente mais fraca, domina o cosmos.
\end{itemize}
\emph{A natureza usa a força fraca para construir galáxias e a força forte para
construir átomos --- precisamente porque uma se cancela e a outra não.}}
\end{questao}

\begin{questao}[3]{}
Duas pequenas esferas idênticas, de massa $\SI{0.5}{\gram}$ cada, são suspensas
do mesmo ponto por fios isolantes de $\SI{20}{\centi\meter}$. Ao receberem cargas
iguais $q$, elas se afastam até que cada fio forme $\SI{10}{\degree}$ com a
vertical.

\circuitoq{%
\begin{tikzpicture}[scale=1,line width=0.8pt,>=Stealth]
 \draw[cinza!60,line width=2pt] (-1.6,3)--(1.6,3);
 \coordinate (O) at (0,3);
 \coordinate (P1) at ({-2.4*sin(10)},{3-2.4*cos(10)});
 \coordinate (P2) at ({ 2.4*sin(10)},{3-2.4*cos(10)});
 \draw[cinza,thick] (O)--(P1); \draw[cinza,thick] (O)--(P2);
 \draw[dashed,cinza!70] (O)--(0,0.3);
 \shade[ball color=Vcol!30] (P1) circle (0.2);
 \shade[ball color=Vcol!30] (P2) circle (0.2);
 \node[font=\footnotesize] at (0.34,2.55) {$10^\circ$};
 \node[left,font=\footnotesize] at ($(P1)+(-0.25,0)$) {$q$};
 \node[right,font=\footnotesize] at ($(P2)+(0.25,0)$) {$q$};
 \draw[<->,cinza] ($(P1)+(0,-0.45)$)--($(P2)+(0,-0.45)$)
       node[midway,below,font=\footnotesize]{$d$};
 \node[right,font=\footnotesize,cinza] at (0.9,1.9) {$L=\SI{20}{\centi\meter}$};
\end{tikzpicture}}{Pêndulos eletrostáticos idênticos.}

\begin{itensq}
\item Determine a separação $d$ entre as esferas.
\item Calcule a força elétrica entre elas.
\item Determine a carga $q$ de cada esfera.
\end{itensq}
\dados{$g=\SI{10}{\meter\per\second\squared}$; $\sin10^\circ=0{,}174$;
$\tan10^\circ=0{,}176$; $\ke=\SI{9e9}{}$.}
\resposta{(a) $d\approx\SI{6.9}{\centi\meter}$; (b) $F\approx\SI{0.88}{\milli\newton}$;
(c) $q\approx\SI{22}{\nano\coulomb}$}
\resolucao{(a) Cada esfera se afasta $L\sin\theta$ da vertical, logo
\[
d=2L\sin\theta=2(0{,}20)(0{,}174)\approx\SI{0.069}{\meter}=\SI{6.9}{\centi\meter}.
\]
(b) No equilíbrio de cada esfera (tração, peso e força elétrica),
a decomposição fornece $\tan\theta = \dfrac{F_e}{mg}$:
\[
F_e = mg\tan\theta=(\pot{0.5}{-3})(10)(0{,}176)
\approx\pot{8.8}{-4}\,\si{\newton}=\SI{0.88}{\milli\newton}.
\]
(c) Da lei de Coulomb, com as duas cargas iguais:
\[
F_e=\ke\frac{q^{2}}{d^{2}}
\;\Longrightarrow\;
q=d\sqrt{\frac{F_e}{\ke}}
=0{,}069\sqrt{\frac{\pot{8.8}{-4}}{9\cdot10^{9}}}
\approx\pot{2.2}{-8}\,\si{\coulomb}=\SI{22}{\nano\coulomb}.
\]
\textbf{Ordem de grandeza:} apenas $\sim\pot{1.4}{11}$ elétrons em excesso ---
uma quantidade ínfima --- já produz um afastamento visível. \emph{Isso mostra
quão intensa é a força elétrica em escala macroscópica.}}
\end{questao}

\begin{questao}[3]{}
Quatro cargas de mesmo módulo $q=\SI{1}{\micro\coulomb}$ ocupam os vértices de um
quadrado de lado $a=\SI{10}{\centi\meter}$, \textbf{alternando os sinais}
($+,-,+,-$ ao percorrer o perímetro). Determine o módulo e a direção da força
resultante sobre uma das cargas positivas.
\dados{$\ke=\SI{9e9}{}$; $\sqrt2\approx1{,}41$.}
\resposta{$F_R \approx \SI{0.82}{\newton}$, dirigida ao centro do quadrado}
\resolucao{\textbf{Forças envolvidas.} Sobre a carga positiva escolhida atuam:
\begin{itemize}[leftmargin=1.6em,itemsep=1pt,topsep=2pt]
\item duas forças dos vizinhos \emph{adjacentes} (que são \textbf{negativos}):
      são de \textbf{atração}, ao longo dos lados;
\item uma força da carga da \emph{diagonal} (que é \textbf{positiva}):
      \textbf{repulsão}, ao longo da diagonal.
\end{itemize}
\textbf{Módulos:}
\[
F_L=\ke\frac{q^{2}}{a^{2}}=9\cdot10^{9}\frac{(\pot{1}{-6})^{2}}{(0{,}1)^{2}}
=\SI{0.9}{\newton},
\qquad
F_D=\ke\frac{q^{2}}{(a\sqrt2)^{2}}=\frac{F_L}{2}=\SI{0.45}{\newton}.
\]
\textbf{Composição.} As duas atrações, perpendiculares entre si, têm resultante
$F_L\sqrt2 \approx \SI{1.27}{\newton}$ apontando \emph{para o centro} do quadrado
(ao longo da diagonal). A repulsão da diagonal aponta \emph{para fora}, na mesma
reta. Logo elas se \textbf{subtraem}:
\[
F_R=F_L\sqrt2-F_D = 1{,}27-0{,}45\approx\SI{0.82}{\newton},
\]
dirigida \textbf{para o centro} do quadrado.\\
\textbf{Contraste com o caso de cargas iguais} (visto anteriormente): lá as três
forças eram repulsivas e se \emph{somavam} ($F_R=F_L\sqrt2+F_D=\SI{1.72}{\newton}$),
apontando para fora. Aqui, alternar os sinais inverte o sentido e reduz o módulo
para menos da metade. \emph{A geometria é a mesma; a configuração de sinais muda
tudo.}}
\end{questao}

\begin{questao}[3]{}
Duas cargas fixas, $q_1=+q$ e $q_2=+4q$, estão separadas por $d$. Deseja-se
colocar uma terceira carga $q_3$ de modo que \textbf{as três} fiquem em
equilíbrio simultaneamente.
\begin{itensq}
\item Determine a posição de $q_3$.
\item Determine o valor e o sinal de $q_3$.
\item Esse equilíbrio é estável?
\end{itensq}
\resposta{(a) a $d/3$ de $q_1$; (b) $q_3=-\dfrac{4q}{9}$; (c) instável}
\resolucao{(a) Para $q_3$ estar em equilíbrio, ela deve ficar onde os campos de
$q_1$ e $q_2$ se cancelam (calculado anteriormente):
\[
\frac{q}{x^{2}}=\frac{4q}{(d-x)^{2}}
\;\Rightarrow\;
\frac{d-x}{x}=2
\;\Rightarrow\;
x=\frac{d}{3}.
\]
(b) Mas isso garante apenas o equilíbrio de $q_3$. Para que \emph{$q_1$} também
esteja equilibrada, a repulsão de $q_2$ sobre ela deve ser cancelada pela atração
de $q_3$:
\[
\underbrace{\ke\frac{q\cdot 4q}{d^{2}}}_{\text{repulsão de }q_2}
=\underbrace{\ke\frac{q\,|q_3|}{(d/3)^{2}}}_{\text{atração de }q_3}
\;\Longrightarrow\;
\frac{4q}{d^{2}}=\frac{9|q_3|}{d^{2}}
\;\Longrightarrow\;
|q_3|=\frac{4q}{9}.
\]
O sinal deve ser \textbf{negativo} (para atrair), logo
$q_3 = -\dfrac{4q}{9}$.\\
\textbf{Verificação em $q_2$:} a repulsão de $q_1$ sobre $q_2$ é
$\ke\dfrac{4q^2}{d^2}$; a atração de $q_3$ (à distância $2d/3$) é
$\ke\dfrac{4q\cdot(4q/9)}{(2d/3)^{2}}=\ke\dfrac{(16q^2/9)}{(4d^2/9)}
=\ke\dfrac{4q^{2}}{d^{2}}$. \checkmark\ São iguais --- as três estão em
equilíbrio.\\
(c) \textbf{Instável}, pelo teorema de Earnshaw: qualquer deslocamento infinitesimal
de qualquer das cargas destrói o equilíbrio, e as forças resultantes \emph{afastam}
ainda mais o sistema da configuração original. Não existe arranjo eletrostático
puro que seja estável.}
\end{questao}

\begin{questao}[3]{}
Uma pequena esfera de massa $m = \SI{10}{\gram}$, eletrizada com carga $q$, é
suspensa por um fio isolante de massa desprezível em uma região de campo elétrico
\emph{horizontal} e uniforme, de intensidade $E = \SI{e4}{\newton\per\coulomb}$.
No equilíbrio, o fio forma $\SI{30}{\degree}$ com a vertical.

\circuitoq{%
\begin{tikzpicture}[scale=1,line width=0.8pt,>=Stealth]
 \draw[cinza!60,line width=2pt] (-2,2.6)--(2,2.6);
 \coordinate (O) at (0,2.6);
 \coordinate (P) at ({2.2*sin(30)},{2.6-2.2*cos(30)});
 \draw[cinza,thick] (O)--(P);
 \draw[dashed,cinza!70] (O)--(0,0.2);
 \shade[ball color=Vcol!30] (P) circle (0.22);
 \node[right,font=\footnotesize] at ($(P)+(0.28,0)$) {$m,\,q$};
 \draw[->,Vcol,thick] (P)--($(P)+(1.0,0)$) node[right,font=\footnotesize]{$\vec F_e$};
 \draw[->,cinza,thick] (P)--($(P)+(0,-0.9)$) node[below,font=\footnotesize]{$\vec P$};
 \node[font=\footnotesize] at (0.42,1.9) {$30^\circ$};
 \foreach \y in {0.4,1.0,1.6} \draw[->,Vcol!45] (-1.9,\y)--(-1.2,\y);
 \node[Vcol,font=\footnotesize] at (-1.55,2.05) {$\vec E$};
\end{tikzpicture}}{Pêndulo eletrostático em campo horizontal.}

\begin{itensq}
\item Determine a carga $q$ da esfera.
\item Calcule a tração no fio.
\item Se o campo dobrasse, o ângulo dobraria? Justifique quantitativamente.
\end{itensq}
\dados{$g = \SI{10}{\meter\per\second\squared}$; $\tan30^\circ\approx0{,}577$.}
\resposta{(a) $q \approx \SI{5.8}{\micro\coulomb}$; (b) $T\approx\SI{0.115}{\newton}$;
(c) não --- o ângulo cresce menos que proporcionalmente}
\resolucao{\textbf{(a)} No equilíbrio, decompondo a tração:
\[
T\sin\theta = F_e = qE,
\qquad
T\cos\theta = P = mg
\quad\Longrightarrow\quad
\tan\theta=\frac{qE}{mg}.
\]
Isolando $q$:
\[
q=\frac{mg\tan\theta}{E}
=\frac{(\pot{10}{-3})(10)(0{,}577)}{\pot{1}{4}}
\approx\pot{5.8}{-6}\,\si{\coulomb}=\SI{5.8}{\micro\coulomb}.
\]
\textbf{(b)} $T = \dfrac{mg}{\cos\theta} = \dfrac{0{,}1}{0{,}866}
\approx\SI{0.115}{\newton}$.\\
\textbf{(c)} \textbf{Não.} A relação é $\tan\theta \propto E$, e \emph{não}
$\theta \propto E$. Dobrando o campo:
\[
\tan\theta' = 2\tan30^\circ = 1{,}155
\;\Rightarrow\;
\theta' = \arctan(1{,}155)\approx 49{,}1^\circ,
\]
que é \emph{menos} que $2\times30^\circ = 60^\circ$. A função tangente cresce mais
rápido que o ângulo, então o ângulo responde de forma \emph{sublinear}. No limite
$E\to\infty$, $\theta\to90^\circ$ mas nunca o atinge --- o fio jamais fica
horizontal.}
\end{questao}

\begin{questao}[3]{}
Três cargas puntiformes estão sobre uma reta: $+Q$ na origem, $-2Q$ em $x=a$ e
$+Q$ em $x=2a$ (um \emph{quadrupolo linear}).
\begin{itensq}
\item Determine a força resultante sobre a carga central.
\item Determine a força resultante sobre a carga da extremidade em $x=0$.
\item Explique por que o sistema, embora tenha carga total nula, ainda exerce
      força sobre uma carga externa distante.
\end{itensq}
\resposta{(a) nula, por simetria; (b) $F=\dfrac{7\ke Q^{2}}{4a^{2}}$ dirigida para
$+x$; (c) ver comentário}
\resolucao{(a) A carga central $-2Q$ está \emph{equidistante} das duas cargas
$+Q$, que são iguais. As duas atrações têm mesmo módulo e sentidos opostos: a
resultante é \textbf{nula} por simetria.\\
(b) Sobre a carga em $x=0$ atuam:
\begin{itemize}[leftmargin=1.6em,itemsep=1pt,topsep=2pt]
\item atração de $-2Q$ (distância $a$):
      $F_1=\ke\dfrac{Q\cdot2Q}{a^{2}}=\dfrac{2\ke Q^{2}}{a^{2}}$, para $+x$;
\item repulsão de $+Q$ (distância $2a$):
      $F_2=\ke\dfrac{Q^{2}}{4a^{2}}$, para $-x$.
\end{itemize}
Resultante:
\[
F=F_1-F_2=\frac{2\ke Q^{2}}{a^{2}}-\frac{\ke Q^{2}}{4a^{2}}
=\frac{\ke Q^{2}}{a^{2}}\left(2-\frac14\right)
=\frac{7\ke Q^{2}}{4a^{2}},
\]
dirigida para $+x$ (para dentro do conjunto).\\
(c) \textbf{Por que a carga total nula não implica campo nulo.} A carga total é
$Q-2Q+Q=0$, então \emph{a longas distâncias} o campo decai muito mais rápido que
o de uma carga isolada. Mas ele não é nulo: as cargas ocupam \emph{posições
diferentes}, e essa separação espacial cria momentos de ordem superior.\\
\textbf{Hierarquia dos decaimentos:}
\begin{center}\small\sffamily
\begin{tabular}{@{}l c c@{}}
\toprule
\textbf{Configuração} & \textbf{Carga total} & \textbf{Campo distante}\\
\midrule
Monopolo (carga isolada) & $\neq0$ & $\propto 1/r^{2}$\\
Dipolo ($+q,-q$)         & $0$     & $\propto 1/r^{3}$\\
Quadrupolo (este caso)   & $0$     & $\propto 1/r^{4}$\\
\bottomrule
\end{tabular}
\end{center}
\emph{Cada cancelamento sucessivo acelera o decaimento --- é o princípio da
expansão multipolar, base da análise de antenas e de moléculas polares.}}
\end{questao}

\blocodesafio

\begin{questao}[4]{}
Uma carga $Q$ é dividida em duas partes, $q$ e $(Q-q)$, que são então colocadas a
uma distância fixa $d$ uma da outra.
\begin{itensq}
\item Obtenha a expressão da força entre elas em função de $q$.
\item Determine o valor de $q$ que \textbf{maximiza} essa força.
\item Qual é a força máxima?
\end{itensq}
\resposta{(a) $F=\dfrac{\ke\,q(Q-q)}{d^{2}}$; (b) $q=Q/2$;
(c) $F_{\max}=\dfrac{\ke Q^{2}}{4d^{2}}$}
\resolucao{(a) Pela lei de Coulomb,
\[
F(q)=\ke\,\frac{q\,(Q-q)}{d^{2}}=\frac{\ke}{d^{2}}\left(Qq-q^{2}\right).
\]
(b) É uma \textbf{parábola} com concavidade para baixo em função de $q$. O máximo
ocorre no vértice. Derivando:
\[
\frac{\mathrm{d}F}{\mathrm{d}q}=\frac{\ke}{d^{2}}\left(Q-2q\right)=0
\;\Longrightarrow\;
\boxed{\,q=\frac{Q}{2}\,}
\]
(Sem cálculo: o vértice de $Qq-q^2$ está em $q = -\dfrac{b}{2a}=\dfrac{Q}{2}$.)\\
(c) Substituindo:
\[
F_{\max}=\frac{\ke}{d^{2}}\cdot\frac{Q}{2}\cdot\frac{Q}{2}
=\frac{\ke\,Q^{2}}{4d^{2}}.
\]
\textbf{Interpretação.} A força é máxima quando a carga é dividida
\emph{igualmente}. Isso decorre da desigualdade entre médias: para soma fixa
($q+(Q-q)=Q$), o produto é máximo quando as parcelas são iguais.\\
\textbf{Casos-limite:} se $q\to0$ ou $q\to Q$, uma das partes fica sem carga e
$F\to0$ --- coerente. \emph{Este é um problema de otimização disfarçado de
eletrostática: reconhecer a estrutura matemática vale mais que decorar a
fórmula.}}
\end{questao}

% BEGIN BANCO COMPLEMENTAR Lei de Coulomb
% =====================================================================
%  Banco complementar --- Lei de Coulomb  (REVISADO)
%  Substitui a versao gerada por template: 32 questoes, todas com
%  enunciado distinto. Nivel 4 reclassificado conforme o criterio da
%  obra (sintese, demonstracao, otimizacao ou modelagem nao rotineira).
%  Distribuicao mantida: 7 N1 + 10 N2 + 6 N3 + 9 N4 = 32
% =====================================================================

\subsubsection*{Banco complementar --- Lei de Coulomb}

\begin{atencao}[Organização do banco]
As questões abaixo complementam o tópico sem repetir enunciados da seção
principal. Cada nível explora uma \emph{estrutura de raciocínio} diferente:
N1 aplica a fórmula em suas quatro formas isoladas ($F$, $q$, $r$, meio);
N2 exige composição ou reinterpretação; N3 combina vetores, equilíbrio e
redistribuição de carga; N4 pede demonstração, integração, otimização ou
modelagem.
\end{atencao}

\blocofixacao

\begin{questao}[1]{}
Duas cargas puntiformes $q_1=\SI{20}{\nano\coulomb}$ e
$q_2=\SI{30}{\nano\coulomb}$ estão separadas por $\SI{3}{\centi\meter}$ no
vácuo. Determine o módulo da força elétrica e classifique-a.
\dados{$\ke=\SI{9e9}{\newton\meter\squared\per\coulomb\squared}$}
\resposta{$F=\SI{6}{\milli\newton}$, repulsiva}
\resolucao{Convertendo antes de substituir --- este é o erro mais comum:
$q_1=\pot{2}{-8}\,\si{\coulomb}$, $q_2=\pot{3}{-8}\,\si{\coulomb}$,
$r=\pot{3}{-2}\,\si{\meter}$.
\[
F=\ke\frac{q_1q_2}{r^{2}}
 =\frac{\pot{9}{9}\cdot\pot{2}{-8}\cdot\pot{3}{-8}}{(\pot{3}{-2})^{2}}
 =\frac{\pot{5.4}{-6}}{\pot{9}{-4}}=\pot{6}{-3}\,\si{\newton}.
\]
Ambas positivas $\Rightarrow$ repulsão.}
\end{questao}

\begin{questao}[1]{}
Duas esferas idênticas recebem cargas iguais $q$ e se repelem com força de
$\SI{0.90}{\newton}$ quando separadas por $\SI{20}{\centi\meter}$.
\begin{itensq}
\item Determine $q$.
\item Quantos elétrons foram \textbf{retirados} de cada esfera?
\end{itensq}
\dados{$e=\pot{1.6}{-19}\,\si{\coulomb}$}
\resposta{(a) $q=\SI{2}{\micro\coulomb}$; (b) $n=\pot{1.25}{13}$ elétrons}
\resolucao{(a) De $F=\ke q^{2}/r^{2}$,
\[
q=\sqrt{\frac{Fr^{2}}{\ke}}=\sqrt{\frac{0{,}90\cdot(0{,}20)^{2}}{\pot{9}{9}}}
 =\sqrt{\pot{4}{-12}}=\pot{2}{-6}\,\si{\coulomb}.
\]
(b) Repulsão com cargas iguais e sinal positivo indica \emph{falta} de
elétrons: $n=q/e=\pot{2}{-6}/\pot{1.6}{-19}=\pot{1.25}{13}$.}
\end{questao}

\begin{questao}[1]{}
Duas cargas de $\SI{5}{\micro\coulomb}$ cada devem exercer uma sobre a outra
força de $\SI{2.5}{\newton}$ no vácuo. A que distância devem ser colocadas?
\resposta{$r=\SI{0.30}{\meter}$}
\resolucao{Isolando $r$ na lei de Coulomb:
\[
r=\sqrt{\frac{\ke q_1q_2}{F}}
 =\sqrt{\frac{\pot{9}{9}\cdot(\pot{5}{-6})^{2}}{2{,}5}}
 =\sqrt{\frac{0{,}225}{2{,}5}}=\sqrt{0{,}09}=\SI{0.30}{\meter}.
\]}
\end{questao}

\begin{questao}[1]{}
Duas cargas puntiformes se atraem com força $F$. Mantendo as cargas e
\textbf{dobrando} a distância entre elas, a nova força será:
\begin{alternativas}[3]
\item $F/4$
\item $F/2$
\item $F$
\item $2F$
\item $4F$
\end{alternativas}
\resposta{(a)}
\resolucao{$F\propto 1/r^{2}$. Dobrar $r$ divide a força por $2^{2}=4$:
$F'=F/4$. Note que a natureza (atrativa) não muda --- ela depende apenas dos
sinais, nunca da distância.}
\end{questao}

\begin{questao}[1]{}
Duas cargas separadas por uma distância fixa exercem força de
$\SI{0.80}{\newton}$ no vácuo. Todo o conjunto é então imerso em um líquido
isolante de permissividade relativa $\varepsilon_r=4$. Determine a nova força.
\resposta{$F'=\SI{0.20}{\newton}$}
\resolucao{Em um meio dielétrico, $F_{\text{meio}}=F_{\text{vácuo}}/\varepsilon_r$,
pois $\ke$ é substituída por $\ke/\varepsilon_r$:
\[
F'=\frac{0{,}80}{4}=\SI{0.20}{\newton}.
\]
O dielétrico se polariza e blinda parcialmente as cargas --- por isso a força
\emph{sempre} diminui ($\varepsilon_r\ge 1$).}
\end{questao}

\begin{questao}[1]{}
Um bastão carregado positivamente é aproximado --- sem tocar --- de uma esfera
metálica \textbf{neutra} e isolada. A força entre eles é:
\begin{alternativas}[2]
\item nula, pois a esfera é neutra
\item atrativa
\item repulsiva
\item atrativa ou repulsiva, dependendo da distância
\end{alternativas}
\resposta{(b)}
\resolucao{A carga \emph{líquida} da esfera continua nula, mas o bastão
positivo atrai elétrons livres para a face próxima e repele carga positiva
para a face distante (\textbf{indução eletrostática}). A face próxima, negativa,
está \emph{mais perto} do bastão que a face distante, positiva. Como
$F\propto 1/r^{2}$, a atração vence: resultante \textbf{atrativa}. É o mesmo
princípio pelo qual um canudo atritado atrai pedaços de papel neutros.}
\end{questao}

\begin{questao}[1]{}
Estime a força entre duas cargas de $\SI{1}{\coulomb}$ separadas por
$\SI{1}{\kilo\meter}$ e comente o resultado.
\resposta{$F=\SI{9}{\kilo\newton}$ (o peso de cerca de \SI{900}{\kilogram})}
\resolucao{
\[
F=\frac{\pot{9}{9}\cdot 1\cdot 1}{(10^{3})^{2}}=\pot{9}{3}\,\si{\newton}.
\]
Nove mil newtons --- o peso de quase uma tonelada --- a \emph{um quilômetro}
de distância. Isso mostra por que o coulomb é uma unidade enorme para a
eletrostática: cargas realmente separáveis por atrito ficam na faixa de
$\si{\nano\coulomb}$ a $\si{\micro\coulomb}$.}
\end{questao}

\blocoaplicacao

\begin{questao}[2]{}
Três cargas estão fixas sobre o eixo $x$: $q_1=\SI{+3}{\micro\coulomb}$ em
$x=0$, $q_2=\SI{+2}{\micro\coulomb}$ em $x=\SI{20}{\centi\meter}$ e
$q_3=\SI{-5}{\micro\coulomb}$ em $x=\SI{50}{\centi\meter}$. Determine a força
resultante sobre $q_2$ (módulo e sentido).
\resposta{$F=\SI{2.35}{\newton}$, no sentido $+x$}
\resolucao{Trate cada par separadamente e some \emph{com sinal}.
\[
F_{12}=\frac{\pot{9}{9}\cdot\pot{3}{-6}\cdot\pot{2}{-6}}{(0{,}20)^{2}}
=\SI{1.35}{\newton}\quad(\text{repulsão}\Rightarrow +x)
\]
\[
F_{32}=\frac{\pot{9}{9}\cdot\pot{2}{-6}\cdot\pot{5}{-6}}{(0{,}30)^{2}}
=\SI{1.00}{\newton}\quad(\text{atração por } q_3 \Rightarrow +x)
\]
As duas apontam no mesmo sentido: $F=1{,}35+1{,}00=\SI{2.35}{\newton}$ em $+x$.
\textbf{Cuidado:} não some as cargas antes; a lei de Coulomb vale par a par
(princípio da superposição).}
\end{questao}

\begin{questao}[2]{}
Duas esferas metálicas idênticas com $q_1=\SI{+8}{\micro\coulomb}$ e
$q_2=\SI{-2}{\micro\coulomb}$ estão a $\SI{30}{\centi\meter}$. Elas são postas
em contato e recolocadas na mesma posição.
\begin{itensq}
\item Calcule a força antes do contato.
\item Calcule a força depois.
\end{itensq}
\resposta{(a) $\SI{1.6}{\newton}$, atrativa; (b) $\SI{0.9}{\newton}$, repulsiva}
\resolucao{(a) $F=\dfrac{\pot{9}{9}\cdot\pot{8}{-6}\cdot\pot{2}{-6}}{(0{,}30)^{2}}
=\dfrac{0{,}144}{0{,}09}=\SI{1.6}{\newton}$, atrativa (sinais opostos).\\
(b) Esferas \emph{idênticas} em contato dividem a carga total igualmente:
\[
q'=\frac{q_1+q_2}{2}=\frac{+8-2}{2}=\SI{+3}{\micro\coulomb}\ \text{cada}.
\]
\[
F'=\frac{\pot{9}{9}\cdot(\pot{3}{-6})^{2}}{0{,}09}=\SI{0.9}{\newton},
\ \text{agora repulsiva}.
\]
A força \emph{e} a natureza mudaram: o contato conserva a carga total, não a
individual.}
\end{questao}

\begin{questao}[2]{}
Duas cargas fixas exercem entre si $\SI{0.60}{\newton}$ no vácuo. O conjunto é
mergulhado em óleo com $\varepsilon_r=2{,}5$.
\begin{itensq}
\item Qual é a nova força, mantida a distância?
\item Que fração da distância original restauraria a força de
      $\SI{0.60}{\newton}$?
\end{itensq}
\resposta{(a) $\SI{0.24}{\newton}$; (b) $r'=r/\sqrt{2{,}5}\approx 0{,}63\,r$}
\resolucao{(a) $F'=F/\varepsilon_r=0{,}60/2{,}5=\SI{0.24}{\newton}$.\\
(b) Queremos $\dfrac{\ke q_1q_2}{\varepsilon_r r'^{2}}=\dfrac{\ke q_1q_2}{r^{2}}$,
ou seja $\varepsilon_r r'^{2}=r^{2}$:
\[
r'=\frac{r}{\sqrt{\varepsilon_r}}=\frac{r}{\sqrt{2{,}5}}=0{,}632\,r.
\]
Aproximar as cargas em cerca de $37\%$ compensa exatamente o efeito do meio.}
\end{questao}

\begin{questao}[2]{}
A força entre duas cargas $q_1$ e $q_2$ separadas por $d$ vale $F$. Se $q_1$ for
\textbf{triplicada} e a distância for reduzida à \textbf{metade}, a nova força
será:
\begin{alternativas}[3]
\item $1{,}5F$
\item $3F$
\item $6F$
\item $12F$
\item $24F$
\end{alternativas}
\resposta{(d)}
\resolucao{Os efeitos são multiplicativos e independentes:
\[
F'=\ke\frac{(3q_1)q_2}{(d/2)^{2}}=3\cdot 4\cdot \ke\frac{q_1q_2}{d^{2}}=12F.
\]
Fator $3$ da carga $\times$ fator $4$ da distância $=12$.}
\end{questao}

\begin{questao}[2]{}
Em uma impressora a jato eletrostático, cada gota de tinta tem massa
$\pot{1.0}{-11}\,\si{\kilogram}$ e é carregada pela remoção de
$\pot{2.0}{5}$ elétrons. Duas gotas vizinhas ficam a $\SI{1.0}{\milli\meter}$.
Compare a força elétrica entre elas com o peso de uma gota.
\dados{$e=\pot{1.6}{-19}\,\si{\coulomb}$; $g=\SI{9.8}{\meter\per\second\squared}$}
\resposta{$F_e=\pot{9.2}{-12}\,\si{\newton}$; $P=\pot{9.8}{-11}\,\si{\newton}$;
$F_e/P\approx 0{,}094$}
\resolucao{Carga da gota: $q=ne=\pot{2.0}{5}\cdot\pot{1.6}{-19}
=\pot{3.2}{-14}\,\si{\coulomb}$.
\[
F_e=\frac{\pot{9}{9}\cdot(\pot{3.2}{-14})^{2}}{(\pot{1.0}{-3})^{2}}
=\pot{9.2}{-12}\,\si{\newton}.
\]
Peso: $P=mg=\pot{1.0}{-11}\cdot 9{,}8=\pot{9.8}{-11}\,\si{\newton}$.
Razão $\approx 9{,}4\%$: a repulsão entre gotas é pequena diante do peso, mas
não desprezível --- é ela que degrada a nitidez quando as gotas são emitidas
muito próximas.}
\end{questao}

\begin{questao}[2]{}
Duas cargas positivas, $q_1=\SI{4}{\micro\coulomb}$ em $x=0$ e
$q_2=\SI{9}{\micro\coulomb}$ em $x=\SI{50}{\centi\meter}$, estão fixas.
Determine o ponto onde uma terceira carga ficaria em equilíbrio.
\resposta{$x=\SI{20}{\centi\meter}$}
\resolucao{Como as duas são positivas, o ponto de força nula está
\emph{entre} elas. Seja $x$ a distância a $q_1$:
\[
\frac{\ke q_1 q_3}{x^{2}}=\frac{\ke q_2 q_3}{(0{,}50-x)^{2}}
\;\Longrightarrow\;
\frac{4}{x^{2}}=\frac{9}{(0{,}50-x)^{2}}.
\]
Extraindo a raiz (ambos os lados positivos):
$\dfrac{2}{x}=\dfrac{3}{0{,}50-x}\Rightarrow 2(0{,}50-x)=3x
\Rightarrow x=\SI{0.20}{\meter}$.
O ponto fica \emph{mais próximo da carga menor}, como esperado. Note que $q_3$
se cancela: a posição não depende dela.}
\end{questao}

\begin{questao}[2]{}
Uma carga $q$ e outra $100q$ interagem. Sobre os módulos das forças que uma
exerce sobre a outra, é correto afirmar:
\begin{alternativas}[2]
\item a força sobre $q$ é 100 vezes maior
\item a força sobre $100q$ é 100 vezes maior
\item as forças têm o mesmo módulo
\item depende das massas
\end{alternativas}
\resposta{(c)}
\resolucao{A lei de Coulomb é simétrica no produto $q_1q_2$: ambas valem
$\ke\,q\cdot 100q/r^{2}$. Formam um par ação--reação (3\textsuperscript{a} lei
de Newton): mesmo módulo, mesma direção, sentidos opostos. O que \emph{difere}
é a \textbf{aceleração}, se as massas forem diferentes --- confundir força com
aceleração é o erro clássico aqui.}
\end{questao}

\begin{questao}[2]{}
No modelo de Bohr para o hidrogênio, o elétron orbita o próton a
$\pot{5.3}{-11}\,\si{\meter}$. Determine a força elétrica e a aceleração
centrípeta do elétron.
\dados{$e=\pot{1.6}{-19}\,\si{\coulomb}$; $m_e=\pot{9.11}{-31}\,\si{\kilogram}$}
\resposta{$F=\pot{8.2}{-8}\,\si{\newton}$;
$a=\pot{9.0}{22}\,\si{\meter\per\second\squared}$}
\resolucao{
\[
F=\frac{\pot{9}{9}\cdot(\pot{1.6}{-19})^{2}}{(\pot{5.3}{-11})^{2}}
=\frac{\pot{2.304}{-28}}{\pot{2.81}{-21}}=\pot{8.2}{-8}\,\si{\newton}.
\]
\[
a=\frac{F}{m_e}=\frac{\pot{8.2}{-8}}{\pot{9.11}{-31}}
=\pot{9.0}{22}\,\si{\meter\per\second\squared}.
\]
Uma força minúscula produz aceleração de $10^{22}\,\si{\meter\per\second\squared}$
porque a massa do elétron é ainda mais minúscula.}
\end{questao}

\begin{questao}[2]{}
Duas esferas metálicas idênticas com cargas $+Q$ e $+3Q$, separadas por $d$,
repelem-se com força $F$. Elas são postas em contato e depois afastadas até
$2d$. Determine a nova força em função de $F$.
\resposta{$F'=F/3$}
\resolucao{Antes: $F=\ke\dfrac{Q\cdot 3Q}{d^{2}}=\dfrac{3\ke Q^{2}}{d^{2}}$.\\
Após o contato, cada esfera fica com $\dfrac{Q+3Q}{2}=2Q$. A $2d$:
\[
F'=\ke\frac{(2Q)(2Q)}{(2d)^{2}}=\ke\frac{4Q^{2}}{4d^{2}}=\frac{\ke Q^{2}}{d^{2}}
=\frac{F}{3}.
\]
O produto das cargas subiu (de $3Q^2$ para $4Q^2$), mas o fator $4$ da
distância dominou.}
\end{questao}

\begin{questao}[2]{}
Compare a força entre duas cargas de $\SI{1.0}{\micro\coulomb}$ separadas por
$\SI{1.0}{\centi\meter}$ com o peso de um corpo de $\SI{1.0}{\kilogram}$, e
interprete.
\dados{$g=\SI{9.8}{\meter\per\second\squared}$}
\resposta{$F=\SI{90}{\newton}$ contra $P=\SI{9.8}{\newton}$; razão $\approx 9{,}2$}
\resolucao{
\[
F=\frac{\pot{9}{9}\cdot(\pot{1}{-6})^{2}}{(\pot{1}{-2})^{2}}
=\frac{\pot{9}{-3}}{\pot{1}{-4}}=\SI{90}{\newton};
\qquad P=1{,}0\cdot 9{,}8=\SI{9.8}{\newton}.
\]
Um milionésimo de coulomb, a um centímetro, produz nove vezes o peso de um
quilograma. Essa desproporção entre a força elétrica e a gravitacional é a razão
pela qual é a matéria macroscópica é praticamente neutra: qualquer desequilíbrio
apreciável de carga seria violentamente desfeito.}
\end{questao}

\blocoaprofundamento

\begin{questao}[3]{}
Duas cargas fixas sobre o eixo $x$: $q_1=\SI{+4}{\micro\coulomb}$ em $x=0$ e
$q_2=\SI{-1}{\micro\coulomb}$ em $x=\SI{30}{\centi\meter}$. Determine todos os
pontos do eixo onde a força sobre uma carga de prova seria nula.
\resposta{Um único ponto: $x=\SI{60}{\centi\meter}$}
\resolucao{Com \textbf{sinais opostos}, o ponto não pode estar entre as cargas
(ali as duas forças apontam para o mesmo lado). Deve estar \emph{fora} do
segmento e, mais especificamente, do lado da carga de \emph{menor} módulo ---
só assim a maior proximidade compensa o menor valor. Testando $x>0{,}30$:
\[
\frac{\ke\cdot 4\mu}{x^{2}}=\frac{\ke\cdot 1\mu}{(x-0{,}30)^{2}}
\;\Longrightarrow\;
\frac{2}{x}=\frac{1}{x-0{,}30}
\;\Longrightarrow\;
2x-0{,}60=x \Rightarrow x=\SI{0.60}{\meter}.
\]
\textbf{Verificação:}
$\ke\cdot4\mu/(0{,}60)^{2}=\ke\cdot1\mu/(0{,}30)^{2}$ --- confere. A raiz
$x=0{,}20$ da equação quadrática completa é \emph{espúria}: corresponde ao
ponto interno, onde as forças se somam em vez de se cancelar.}
\end{questao}

\begin{questao}[3]{}
Seis cargas iguais $Q=\SI{2}{\micro\coulomb}$ ocupam os vértices de um hexágono
regular de lado $a=\SI{10}{\centi\meter}$, com uma carga
$q=\SI{1}{\micro\coulomb}$ no centro.
\begin{itensq}
\item Qual é a força resultante sobre $q$?
\item Uma das seis cargas é removida. Qual passa a ser a força sobre $q$
      (módulo e direção)?
\end{itensq}
\resposta{(a) nula; (b) $\SI{1.8}{\newton}$, apontando \emph{para} o vértice vazio}
\resolucao{(a) Em um hexágono regular, o centro dista $a$ de todos os vértices
e as cargas ocupam posições diametralmente opostas duas a duas. Cada par
contribui com forças de mesmo módulo e sentidos opostos: resultante nula, por
simetria.\\
(b) Use o argumento da superposição \emph{ao contrário}: 
$\vec{F}_{6}= \vec{F}_{5}+\vec{F}_{\text{removida}}=\vec{0}$, logo
$\vec{F}_{5}=-\vec{F}_{\text{removida}}$.
\[
|\vec{F}_{5}|=\frac{\ke Qq}{a^{2}}
=\frac{\pot{9}{9}\cdot\pot{2}{-6}\cdot\pot{1}{-6}}{(0{,}10)^{2}}
=\SI{1.8}{\newton}.
\]
A carga removida \emph{repelia} $q$ para longe do vértice; sem ela, a resultante
aponta \textbf{para} o vértice vazio. Este truque de simetria evita somar cinco
vetores.}
\end{questao}

\begin{questao}[3]{}
Três cargas ocupam os vértices de um triângulo retângulo:
$q_A=\SI{+2}{\micro\coulomb}$ na origem, $q_B=\SI{+3}{\micro\coulomb}$ em
$(\SI{30}{\centi\meter},0)$ e $q_C=\SI{+4}{\micro\coulomb}$ em
$(0,\SI{40}{\centi\meter})$. Determine a força resultante sobre $q_A$: módulo e
ângulo com o eixo $x$.
\resposta{$F=\SI{0.75}{\newton}$, a $\SI{216.9}{\degree}$ (isto é, $\SI{36.9}{\degree}$ abaixo
do semieixo $-x$)}
\resolucao{As duas forças sobre $q_A$ são perpendiculares --- o que torna a
composição imediata.
\[
F_{AB}=\frac{\pot{9}{9}\cdot\pot{2}{-6}\cdot\pot{3}{-6}}{(0{,}30)^{2}}
=\SI{0.60}{\newton}\ \ (\text{repulsão}\Rightarrow -x)
\]
\[
F_{AC}=\frac{\pot{9}{9}\cdot\pot{2}{-6}\cdot\pot{4}{-6}}{(0{,}40)^{2}}
=\SI{0.45}{\newton}\ \ (\text{repulsão}\Rightarrow -y)
\]
\[
F=\sqrt{0{,}60^{2}+0{,}45^{2}}=\sqrt{0{,}5625}=\SI{0.75}{\newton},
\qquad
\tan\theta=\frac{0{,}45}{0{,}60}=0{,}75 \Rightarrow \theta=\SI{36.87}{\degree}.
\]
Um clássico triângulo $3$--$4$--$5$ escondido. \textbf{Atenção:} a hipotenusa
$BC=\SI{50}{\centi\meter}$ é irrelevante --- ela governa a força \emph{entre}
$q_B$ e $q_C$, não a força sobre $q_A$.}
\end{questao}

\begin{questao}[3]{}
Duas pequenas esferas metálicas idênticas, separadas por $\SI{10}{\centi\meter}$,
\textbf{atraem-se} com força de $\SI{0.108}{\newton}$. Postas em contato e
recolocadas na mesma distância, passam a \textbf{repelir-se} com
$\SI{0.036}{\newton}$. Determine as cargas iniciais.
\resposta{$q_1=\SI{+0.6}{\micro\coulomb}$ e $q_2=\SI{-0.2}{\micro\coulomb}$
(ou a permuta)}
\resolucao{\textbf{Etapa 1 --- produto.} A atração indica sinais opostos:
\[
|q_1q_2|=\frac{F_1r^{2}}{\ke}=\frac{0{,}108\cdot 0{,}01}{\pot{9}{9}}
=\pot{1.2}{-13}\ \Rightarrow\ q_1q_2=-\pot{1.2}{-13}\,\si{\coulomb\squared}.
\]
\textbf{Etapa 2 --- soma.} Após o contato cada esfera fica com
$q'=(q_1+q_2)/2$:
\[
q'=\sqrt{\frac{F_2r^{2}}{\ke}}=\sqrt{\frac{0{,}036\cdot0{,}01}{\pot{9}{9}}}
=\pot{2}{-7}\,\si{\coulomb}
\ \Rightarrow\ q_1+q_2=\pot{4}{-7}\,\si{\coulomb}.
\]
\textbf{Etapa 3 --- soma e produto.} As cargas são raízes de
$x^{2}-Sx+P=0$ com $S=\pot{4}{-7}$ e $P=-\pot{1.2}{-13}$:
\[
x=\frac{\pot{4}{-7}\pm\sqrt{\pot{1.6}{-13}+\pot{4.8}{-13}}}{2}
 =\frac{\pot{4}{-7}\pm\pot{8}{-7}}{2}
 \Rightarrow x=\pot{6}{-7}\ \text{ou}\ -\pot{2}{-7}.
\]
A repulsão final confirma que a soma é não nula. Se as cargas fossem simétricas
($q$ e $-q$), o contato as neutralizaria e a força final seria zero.}
\end{questao}

\begin{questao}[3]{}
Um pequeno bloco de massa $m=\SI{20}{\gram}$, com carga $q=\SI{1}{\micro\coulomb}$,
repousa sobre um plano inclinado \textbf{liso} de $\SI{30}{\degree}$. Uma carga fixa
$Q=\SI{1}{\micro\coulomb}$ está presa na base do plano, e a força elétrica atua
ao longo da superfície. Determine a distância de equilíbrio $d$ medida sobre o
plano.
\dados{$g=\SI{9.8}{\meter\per\second\squared}$}
\resposta{$d\approx\SI{30.3}{\centi\meter}$}
\resolucao{Sobre o plano liso, o equilíbrio na direção da superfície envolve
apenas a componente do peso e a repulsão elétrica (a normal é perpendicular):
\[
\frac{\ke Qq}{d^{2}}=mg\sin\theta.
\]
\[
mg\sin\SI{30}{\degree}=0{,}020\cdot 9{,}8\cdot 0{,}5=\SI{0.098}{\newton};
\qquad
d=\sqrt{\frac{\pot{9}{9}\cdot(\pot{1}{-6})^{2}}{0{,}098}}
=\sqrt{0{,}0918}=\SI{0.303}{\meter}.
\]
\textbf{Estabilidade:} se o bloco descer, $d$ diminui e a repulsão cresce como
$1/d^{2}$, empurrando-o de volta; se subir, a repulsão cai e o peso o traz de
volta. O equilíbrio é \emph{estável}.}
\end{questao}

\begin{questao}[3]{}
Um dipolo é formado por $+q$ e $-q$, com $q=\SI{1}{\micro\coulomb}$, fixos em
$(\pm\SI{5}{\centi\meter},0)$. Uma carga $Q=\SI{2}{\micro\coulomb}$ é colocada
em $(0,\SI{12}{\centi\meter})$, sobre a mediatriz.
\begin{itensq}
\item Mostre que a resultante sobre $Q$ é paralela ao eixo do dipolo.
\item Calcule seu módulo.
\end{itensq}
\resposta{(a) as componentes verticais se cancelam; (b) $F=\SI{0.82}{\newton}$,
no sentido $-x$ (de $+q$ para $-q$)}
\resolucao{(a) Cada carga dista
$r=\sqrt{0{,}05^{2}+0{,}12^{2}}=\SI{0.13}{\meter}$ de $Q$ (triângulo
$5$--$12$--$13$). A carga $+q$ \emph{repele} $Q$ e $-q$ a \emph{atrai}: as duas
forças têm o mesmo módulo e as componentes em $y$ apontam em sentidos opostos,
cancelando-se. Restam as componentes em $x$, que se \textbf{somam}.\\
(b) Cada força vale $F_1=\ke qQ/r^{2}$ e sua componente horizontal é
$F_1\cdot(a/r)$, com $a=\SI{0.05}{\meter}$:
\[
F=2\,\frac{\ke qQ}{r^{2}}\cdot\frac{a}{r}=\frac{2\ke qQ\,a}{r^{3}}
=\frac{2\cdot\pot{9}{9}\cdot\pot{1}{-6}\cdot\pot{2}{-6}\cdot 0{,}05}
       {(0{,}13)^{3}}=\SI{0.82}{\newton}.
\]
Note a dependência $1/r^{3}$: longe do dipolo a força cai \emph{mais rápido}
que a de uma carga isolada, porque as duas cargas quase se cancelam.}
\end{questao}

\blocodesafio

\begin{questao}[4]{}
\textbf{(Modelagem experimental.)} Em uma reprodução da balança de torção, duas
esferas com cargas iguais $q=\SI{0.10}{\micro\coulomb}$ foram mantidas a várias
distâncias e a força foi medida:

\begin{center}
\begin{tabular}{cccccc}
\hline
$r$ (m) & 0,020 & 0,030 & 0,040 & 0,050 & 0,060\\
$F$ (N) & 0,223 & 0,101 & 0,0558 & 0,0362 & 0,0249\\
\hline
\end{tabular}
\end{center}

\begin{itensq}
\item Proponha uma linearização que permita testar a hipótese $F=C\,r^{n}$.
\item Estime $n$ e discuta se os dados sustentam a lei do inverso do quadrado.
\item Obtenha o valor experimental de $\ke$ e compare com o tabelado.
\end{itensq}
\resposta{(a) $\log F=\log C+n\log r$; (b) $n\approx-2{,}00$; (c)
$\ke\approx\pot{9.1}{9}\,\si{\newton\meter\squared\per\coulomb\squared}$
(erro $\approx 1\%$)}
\resolucao{(a) Tomando logaritmos, uma lei de potência vira uma reta:
$\log F = \log C + n\log r$. O coeficiente angular do gráfico
$\log F \times \log r$ é o expoente procurado.\\
(b) Ajustando os cinco pontos por mínimos quadrados obtém-se
$n=-1{,}997$ e $C=\pot{9.09}{-5}$. O expoente reproduz $-2$ com desvio de
$0{,}15\%$ --- dentro do que se espera de medidas de força dessa ordem. Um teste
rápido sem regressão: dobrar $r$ de $0{,}020$ para $0{,}040$ deveria dividir $F$
por $4$; de fato $0{,}223/0{,}0558=4{,}00$.\\
(c) Como $C=\ke q^{2}$ e $q^{2}=\pot{1}{-14}\,\si{\coulomb\squared}$:
\[
\ke=\frac{C}{q^{2}}=\frac{\pot{9.09}{-5}}{\pot{1}{-14}}
=\pot{9.09}{9}\,\si{\newton\meter\squared\per\coulomb\squared},
\]
contra $\pot{8.99}{9}$ tabelado --- erro relativo de $1{,}1\%$.\\
\textbf{Observação de método:} a linearização logarítmica é preferível a ajustar
$F\times 1/r^{2}$ diretamente, porque \emph{não pressupõe} o expoente --- ela o
mede. Esse é exatamente o argumento que Coulomb precisou construir em 1785.}
\end{questao}

\begin{questao}[4]{}
\textbf{(Otimização com cálculo.)} Um anel fino de raio $R$, com carga total $Q$
uniformemente distribuída, está no plano $yz$. Uma carga puntiforme $q$ é
colocada sobre o eixo do anel, a uma distância $x$ do centro.
\begin{itensq}
\item Mostre que $F(x)=\dfrac{\ke Q q\,x}{(x^{2}+R^{2})^{3/2}}$.
\item Determine o valor de $x$ que \textbf{maximiza} a força.
\item Calcule $F_{\max}$ para $Q=\SI{10}{\micro\coulomb}$,
      $q=\SI{2}{\micro\coulomb}$ e $R=\SI{10}{\centi\meter}$.
\end{itensq}
\resposta{(b) $x=R/\sqrt{2}\approx\SI{7.07}{\centi\meter}$;
(c) $F_{\max}=\dfrac{2\ke Qq}{3\sqrt{3}\,R^{2}}=\SI{6.93}{\newton}$}
\resolucao{(a) Divida o anel em elementos $\mathrm{d} Q$. Cada um dista
$s=\sqrt{x^{2}+R^{2}}$ de $q$ e contribui com $\ke q\,\mathrm{d} Q/s^{2}$. Por
simetria, as componentes perpendiculares ao eixo se cancelam aos pares; sobra a
projeção axial, fator $\cos\alpha=x/s$:
\[
F=\int \frac{\ke q\,\mathrm{d} Q}{s^{2}}\cdot\frac{x}{s}
 =\frac{\ke q x}{(x^{2}+R^{2})^{3/2}}\int \mathrm{d} Q
 =\frac{\ke Q q\,x}{(x^{2}+R^{2})^{3/2}}.
\]
(b) Derivando e igualando a zero:
\[
\frac{\mathrm{d} F}{\mathrm{d} x}\propto (x^{2}+R^{2})^{3/2}-x\cdot 3x(x^{2}+R^{2})^{1/2}=0
\;\Rightarrow\; x^{2}+R^{2}-3x^{2}=0 \;\Rightarrow\; x=\frac{R}{\sqrt{2}}.
\]
(c) Substituindo $x=R/\sqrt2$:
\[
F_{\max}=\frac{2\,\ke Qq}{3\sqrt{3}\,R^{2}}
=\frac{2\cdot\pot{9}{9}\cdot\pot{1}{-5}\cdot\pot{2}{-6}}{3\sqrt3\cdot 0{,}01}
=\frac{0{,}36}{0{,}0520}=\SI{6.93}{\newton}.
\]
\textbf{Casos-limite:} em $x=0$ a força é nula (simetria perfeita); para
$x\gg R$, $F\to\ke Qq/x^{2}$ --- o anel se comporta como carga puntiforme.
Entre esses dois zeros deve existir um máximo, e ele está em $R/\sqrt2$.}
\end{questao}

\begin{questao}[4]{}
\textbf{(Distribuição contínua.)} Um bastão isolante retilíneo de comprimento
$L$ possui carga $Q$ uniformemente distribuída. Uma carga puntiforme $q$ está
sobre o prolongamento do bastão, a uma distância $d$ da extremidade mais
próxima.
\begin{itensq}
\item Deduza $F=\dfrac{\ke Qq}{d\,(d+L)}$.
\item Verifique o comportamento para $L\ll d$.
\item Calcule $F$ para $Q=\SI{4}{\micro\coulomb}$, $L=\SI{20}{\centi\meter}$,
      $q=\SI{1}{\micro\coulomb}$ e $d=\SI{10}{\centi\meter}$.
\end{itensq}
\resposta{(b) $F\to\ke Qq/d^{2}$; (c) $F=\SI{1.2}{\newton}$}
\resolucao{(a) Densidade linear $\lambda=Q/L$. Um elemento $\mathrm{d} x$ a uma
distância $x$ da carga carrega $\mathrm{d} Q=\lambda\,\mathrm{d} x$ e todas as contribuições
são \emph{colineares} --- por isso a integral é escalar:
\[
F=\int_{d}^{d+L}\frac{\ke q\lambda\,\mathrm{d} x}{x^{2}}
 =\ke q\lambda\left[-\frac{1}{x}\right]_{d}^{d+L}
 =\ke q\lambda\left(\frac{1}{d}-\frac{1}{d+L}\right)
 =\frac{\ke q\lambda L}{d(d+L)}=\frac{\ke Qq}{d(d+L)}.
\]
(b) Se $L\ll d$, então $d+L\approx d$ e $F\to\ke Qq/d^{2}$: recuperamos a lei de
Coulomb puntiforme, como deve ser.\\
(c) $F=\dfrac{\pot{9}{9}\cdot\pot{4}{-6}\cdot\pot{1}{-6}}{0{,}10\cdot 0{,}30}
=\dfrac{0{,}036}{0{,}03}=\SI{1.2}{\newton}$.\\
\textbf{Compare:} se toda a carga estivesse concentrada no centro do bastão
($\SI{20}{\centi\meter}$), teríamos $\SI{0.9}{\newton}$. A força real é maior
porque as porções mais próximas pesam com $1/x^{2}$ --- concentrar carga no
centro geométrico \emph{subestima} a força.}
\end{questao}

\begin{questao}[4]{}
\textbf{(Demonstração --- estabilidade.)} Duas cargas $+Q$ estão fixas em
$x=-d$ e $x=+d$. Uma carga $q$ é colocada na origem.
\begin{itensq}
\item Mostre que a origem é posição de equilíbrio para qualquer $q$.
\item Para $q>0$, analise a estabilidade a um pequeno deslocamento
      \emph{longitudinal} ($\pm x$).
\item Repita para um deslocamento \emph{transversal} ($\pm y$) e conclua.
\end{itensq}
\resposta{(a) por simetria; (b) estável, com $F\approx-4\ke Qq\,x/d^{3}$;
(c) instável, com $F\approx+2\ke Qq\,y/d^{3}$ --- não há equilíbrio estável em
todas as direções}
\resolucao{(a) As duas forças têm o mesmo módulo e sentidos opostos: resultante
nula, independentemente do sinal e do valor de $q$.\\
(b) Deslocando $q>0$ para $x$ pequeno:
\[
F_x=\ke Qq\left[\frac{1}{(d+x)^{2}}-\frac{1}{(d-x)^{2}}\right].
\]
Com $(d\pm x)^{-2}\approx d^{-2}(1\mp 2x/d)$:
\[
F_x\approx\frac{\ke Qq}{d^{2}}\left(-\frac{4x}{d}\right)=-\frac{4\ke Qq}{d^{3}}\,x.
\]
Sinal negativo $\Rightarrow$ força \textbf{restauradora}: estável nessa direção
(e o movimento é harmônico simples).\\
(c) Deslocando para $y$ pequeno, cada carga dista $\sqrt{d^{2}+y^{2}}\approx d$ e
repele $q$ com componente transversal $\ke Qq\,y/d^{3}$; as duas se somam:
\[
F_y\approx +\frac{2\ke Qq}{d^{3}}\,y,
\]
sinal positivo $\Rightarrow$ força que \textbf{afasta}: instável.\\
\textbf{Conclusão.} O equilíbrio é do tipo \emph{ponto de sela}: estável em uma
direção, instável na outra. Este é um caso particular do \textbf{teorema de
Earnshaw}: nenhuma configuração de cargas fixas pode manter uma carga puntiforme
em equilíbrio estável apenas por forças eletrostáticas. É por isso que modelos
puramente eletrostáticos do átomo não se sustentam.}
\end{questao}

\begin{questao}[4]{}
\textbf{(Série infinita.)} Cargas puntiformes de mesmo módulo $Q$ e
\textbf{sinais alternados} ocupam as posições $x=a,\,2a,\,3a,\dots$ sobre um
eixo: $-Q$ em $a$, $+Q$ em $2a$, $-Q$ em $3a$ e assim por diante,
indefinidamente. Uma carga $+q$ está na origem.
\begin{itensq}
\item Escreva a força resultante sobre $q$ como uma série.
\item Some a série e obtenha a expressão fechada.
\end{itensq}
\dados{$\displaystyle\sum_{n=1}^{\infty}\frac{(-1)^{n+1}}{n^{2}}=\frac{\pi^{2}}{12}$}
\resposta{$F=\dfrac{\pi^{2}}{12}\cdot\dfrac{\ke Qq}{a^{2}}
\approx 0{,}822\,\dfrac{\ke Qq}{a^{2}}$, atrativa (sentido $+x$)}
\resolucao{(a) A $n$-ésima carga está em $x=na$ e tem sinal $(-1)^{n}Q$. Uma
carga negativa \emph{atrai} $q$ (puxa no sentido $+x$); uma positiva
\emph{repele} (sentido $-x$). Tomando $+x$ como positivo:
\[
F=\frac{\ke Qq}{a^{2}}\left(\frac{1}{1^{2}}-\frac{1}{2^{2}}+\frac{1}{3^{2}}
-\frac{1}{4^{2}}+\cdots\right)
=\frac{\ke Qq}{a^{2}}\sum_{n=1}^{\infty}\frac{(-1)^{n+1}}{n^{2}}.
\]
(b) Usando o resultado dado:
\[
F=\frac{\pi^{2}}{12}\cdot\frac{\ke Qq}{a^{2}}\approx 0{,}822\,\frac{\ke Qq}{a^{2}}.
\]
\textbf{Interpretação.} Apesar de haver infinitas cargas, a força é
\emph{finita} e vale menos que a da primeira carga sozinha: a alternância de
sinais produz cancelamento progressivo, e o decaimento $1/n^{2}$ garante
convergência absoluta. Se todas as cargas fossem negativas, a soma seria
$\pi^2/6\approx1{,}645$ --- exatamente o dobro.}
\end{questao}

\begin{questao}[4]{}
\textbf{(Oscilações.)} Uma partícula de massa $m=\SI{1.0}{\gram}$ e carga
$-q=\SI{-2}{\micro\coulomb}$ é liberada em repouso sobre o eixo de um anel de
raio $R=\SI{10}{\centi\meter}$ e carga $Q=\SI{5}{\micro\coulomb}$, a uma
distância $x_0\ll R$ do centro.
\begin{itensq}
\item Mostre que, para $x\ll R$, o movimento é harmônico simples.
\item Obtenha $\omega$ e o período $T$.
\end{itensq}
\resposta{(b) $\omega=\sqrt{\ke Qq/(mR^{3})}=\SI{300}{\radian\per\second}$;
$T\approx\SI{20.9}{\milli\second}$}
\resolucao{(a) A força axial sobre a carga negativa é atrativa, dirigida ao
centro, de módulo $F=\ke Qq\,x/(x^{2}+R^{2})^{3/2}$. Para $x\ll R$, o
denominador tende a $R^{3}$:
\[
F_x\approx-\frac{\ke Qq}{R^{3}}\,x.
\]
Força proporcional ao deslocamento e de sinal oposto --- exatamente a forma
$F=-kx$ da lei de Hooke, com constante efetiva $k_{\text{ef}}=\ke Qq/R^{3}$.
Logo o MHS.\\
(b)
\[
\omega=\sqrt{\frac{k_{\text{ef}}}{m}}=\sqrt{\frac{\ke Qq}{mR^{3}}}
=\sqrt{\frac{\pot{9}{9}\cdot\pot{5}{-6}\cdot\pot{2}{-6}}
{\pot{1}{-3}\cdot(0{,}10)^{3}}}
=\sqrt{\frac{0{,}09}{\pot{1}{-6}}}=\SI{300}{\radian\per\second},
\]
\[
T=\frac{2\pi}{\omega}=\SI{20.9}{\milli\second}.
\]
\textbf{Observe:} o período não depende de $x_0$ --- é a assinatura do MHS. Fora
da aproximação $x\ll R$, a força deixa de ser linear e o movimento continua
periódico, mas com período dependente da amplitude.}
\end{questao}

\begin{questao}[4]{}
\textbf{(Modelagem com taxas.)} Duas esferas idênticas de massa
$m=\SI{2.0}{\gram}$, cada uma com carga $q$, pendem de fios isolantes de
comprimento $L=\SI{60}{\centi\meter}$ presos ao mesmo ponto. Elas se afastam até
uma separação $x\ll L$. Por fuga de carga para o ar úmido, as esferas se
aproximam lentamente.
\begin{itensq}
\item Mostre que $x^{3}=\dfrac{2\ke q^{2}L}{mg}$.
\item Demonstre que
      $\dfrac{\mathrm{d} x}{\mathrm{d} t}=\dfrac{2x}{3q}\dfrac{\mathrm{d} q}{\mathrm{d} t}$.
\item Se em certo instante $x=\SI{6.0}{\centi\meter}$ e as esferas se aproximam
      a $\SI{1.0}{\milli\meter\per\second}$, calcule $\mathrm{d} q/\mathrm{d} t$.
\end{itensq}
\dados{$g=\SI{9.8}{\meter\per\second\squared}$}
\resposta{(c) $q\approx\SI{19.8}{\nano\coulomb}$ e
$\mathrm{d} q/\mathrm{d} t\approx\SI{-0.49}{\nano\coulomb\per\second}$}
\resolucao{(a) Para cada esfera: tensão, peso e repulsão. Do triângulo de
forças, $\tan\theta=F/(mg)$. Com $x\ll L$,
$\tan\theta\approx\sin\theta=(x/2)/L$. Igualando:
\[
\frac{x}{2L}=\frac{\ke q^{2}}{x^{2}mg}
\;\Longrightarrow\;
x^{3}=\frac{2\ke q^{2}L}{mg}.
\]
(b) Derivando implicitamente em relação ao tempo:
\[
3x^{2}\frac{\mathrm{d} x}{\mathrm{d} t}=\frac{4\ke L}{mg}\,q\,\frac{\mathrm{d} q}{\mathrm{d} t}.
\]
De (a), $\dfrac{2\ke L}{mg}=\dfrac{x^{3}}{q^{2}}$; substituindo:
\[
3x^{2}\frac{\mathrm{d} x}{\mathrm{d} t}=\frac{2x^{3}}{q}\frac{\mathrm{d} q}{\mathrm{d} t}
\;\Longrightarrow\;
\frac{\mathrm{d} x}{\mathrm{d} t}=\frac{2x}{3q}\frac{\mathrm{d} q}{\mathrm{d} t}.
\]
(c) Primeiro $q$, por (a):
\[
q=\sqrt{\frac{x^{3}mg}{2\ke L}}
=\sqrt{\frac{(0{,}060)^{3}\cdot 0{,}0020\cdot 9{,}8}{2\cdot\pot{9}{9}\cdot0{,}60}}
=\pot{1.98}{-8}\,\si{\coulomb}.
\]
Com $\mathrm{d} x/\mathrm{d} t=-\pot{1.0}{-3}\,\si{\meter\per\second}$ (aproximação):
\[
\frac{\mathrm{d} q}{\mathrm{d} t}=\frac{3q}{2x}\frac{\mathrm{d} x}{\mathrm{d} t}
=\frac{3\cdot\pot{1.98}{-8}}{2\cdot 0{,}060}\cdot(-\pot{1}{-3})
=-\pot{4.9}{-10}\,\si{\coulomb\per\second}.
\]
\textbf{Leitura física:} medir uma \emph{velocidade} converte-se em medir uma
\emph{corrente de fuga} de meio nanoampère --- pequena demais para muitos
amperímetros, mas visível a olho nu pelo movimento das esferas.}
\end{questao}

\begin{questao}[4]{}
\textbf{(Simetria e integração.)} Um fio semicircular de raio
$R=\SI{5.0}{\centi\meter}$ possui carga $Q=\SI{3}{\micro\coulomb}$ distribuída
uniformemente. Uma carga $q=\SI{1}{\micro\coulomb}$ ocupa o centro do
semicírculo.
\begin{itensq}
\item Justifique por que a resultante é dirigida ao longo do eixo de simetria.
\item Mostre que $F=\dfrac{2\ke Qq}{\pi R^{2}}$ e calcule seu valor.
\end{itensq}
\resposta{$F=\SI{6.9}{\newton}$, ao longo do eixo de simetria}
\resolucao{(a) Para cada elemento do fio existe outro simétrico em relação ao
eixo vertical. Suas contribuições têm componentes perpendiculares ao eixo iguais
e opostas --- que se cancelam --- e componentes ao longo do eixo que se somam.\\
(b) Densidade linear $\lambda=Q/(\pi R)$. O elemento em $\theta$ tem
$\mathrm{d} Q=\lambda R\,\mathrm{d}\theta$, dista $R$ do centro e contribui com
$\ke q\,\mathrm{d} Q/R^{2}$; sua projeção no eixo vale $\sin\theta$:
\[
F=\int_{0}^{\pi}\frac{\ke q\lambda R\,\mathrm{d}\theta}{R^{2}}\sin\theta
 =\frac{\ke q\lambda}{R}\underbrace{\int_{0}^{\pi}\sin\theta\,\mathrm{d}\theta}_{=2}
 =\frac{2\ke q\lambda}{R}=\frac{2\ke Qq}{\pi R^{2}}.
\]
\[
F=\frac{2\cdot\pot{9}{9}\cdot\pot{3}{-6}\cdot\pot{1}{-6}}
{\pi\,(0{,}050)^{2}}=\frac{0{,}054}{\pot{7.85}{-3}}=\SI{6.9}{\newton}.
\]
\textbf{Compare:} se a mesma carga $Q$ estivesse concentrada em um ponto a $R$,
teríamos $\ke Qq/R^{2}=\SI{10.8}{\newton}$. O fator $2/\pi\approx 0{,}64$ mede
exatamente a perda por cancelamento vetorial ao longo do arco. Um anel
\emph{completo} levaria esse fator a zero.}
\end{questao}

\begin{questao}[4]{}
\textbf{(Síntese --- equilíbrio de um sistema.)} Quatro cargas iguais
$Q=\SI{2}{\micro\coulomb}$ ocupam os vértices de um quadrado de lado $a$. Uma
quinta carga $q$ é colocada no centro.
\begin{itensq}
\item Determine $q$ (sinal e módulo, em função de $Q$) para que \textbf{cada}
      carga dos vértices fique em equilíbrio.
\item Calcule $q$ numericamente.
\item O sistema completo está em equilíbrio estável? Justifique.
\end{itensq}
\resposta{(a) $q=-\dfrac{Q(2\sqrt{2}+1)}{4}$; (b) $q\approx\SI{-1.91}{\micro\coulomb}$;
(c) não --- o equilíbrio é instável}
\resolucao{(a) Por simetria, basta impor o equilíbrio de \emph{um} vértice; a
resultante das outras três cargas aponta ao longo da diagonal, para fora.
\begin{itemize}
\item Duas vizinhas, a distância $a$: cada uma contribui $\ke Q^{2}/a^{2}$, e
suas projeções sobre a diagonal valem $\cos\SI{45}{\degree}=\sqrt2/2$ cada, somando
$\sqrt{2}\,\ke Q^{2}/a^{2}$.
\item A oposta, a distância $a\sqrt{2}$: contribui
$\ke Q^{2}/(2a^{2})$, já ao longo da diagonal.
\end{itemize}
\[
F_{\text{vértices}}=\frac{\ke Q^{2}}{a^{2}}\left(\sqrt{2}+\frac{1}{2}\right).
\]
A carga central dista $a/\sqrt{2}$ do vértice, logo atua com
$F_{c}=\ke|q|Q/(a^{2}/2)=2\ke|q|Q/a^{2}$ e deve ser \textbf{atrativa} ---
portanto $q<0$. Igualando:
\[
2|q|=Q\left(\sqrt2+\frac12\right)
\;\Longrightarrow\;
|q|=\frac{Q\,(2\sqrt{2}+1)}{4}\approx 0{,}957\,Q.
\]
(b) $q\approx-0{,}957\cdot 2=\SI{-1.91}{\micro\coulomb}$.\\
(c) \textbf{Não.} A condição imposta zera a resultante sobre os vértices, mas
não sobre a carga central --- que já está em equilíbrio por simetria, embora
instável. Mais fundamentalmente, o teorema de Earnshaw proíbe equilíbrio
eletrostático estável: qualquer perturbação de $q$ ao longo de uma diagonal a
faz cair sobre um vértice. O equilíbrio existe, mas não sobrevive a uma
perturbação infinitesimal.}
\end{questao}

% END BANCO COMPLEMENTAR

\gabarito[1.2 Lei de Coulomb]

% =====================================================================
\subsection{Campo elétrico}
% =====================================================================

\begin{conceito}[Antes de começar]
\textbf{Campo de uma carga puntiforme:}
\[
E = \ke\,\frac{|q|}{d^2}\quad(\si{\newton\per\coulomb}\ \text{ou}\ \si{\volt\per\meter}).
\]
O campo \emph{aponta para fora} de cargas positivas e \emph{para dentro} de cargas
negativas.\par
\textbf{Força sobre uma carga no campo:} $F = qE$ (mesmo sentido de $E$ se $q>0$;
oposto se $q<0$).\par
\textbf{Superposição:} o campo resultante é a \emph{soma vetorial} dos campos
individuais.\par
\textbf{Ponto de campo nulo} entre duas cargas de mesmo sinal: onde os módulos se
igualam, $\dfrac{|q_1|}{d_1^2} = \dfrac{|q_2|}{d_2^2}$.
\end{conceito}

\legendaniveis

\blocofixacao

\begin{questao}[1]{}
Qual a intensidade do campo elétrico gerado por uma carga de
$\SI{-15}{\micro\coulomb}$, a $\SI{10}{\milli\meter}$ dela, em um meio de constante
$\ke = \SI{8e9}{\newton\meter\squared\per\coulomb\squared}$?
\resposta{$E = \pot{1.2}{9}\,\si{\newton\per\coulomb}$}
\resolucao{Com $d = \SI{10}{\milli\meter} = \SI{0.01}{\meter}$:
\[
E=\ke\frac{|q|}{d^2}
=8\cdot10^9\cdot\frac{\pot{15}{-6}}{(0{,}01)^2}=\pot{1.2}{9}\,\si{\newton\per\coulomb}.
\]
O campo aponta \emph{para} a carga, por ela ser negativa.}
\end{questao}

\begin{questao}[1]{}
O módulo do campo elétrico de uma carga puntiforme $q$, em um ponto P a uma dada
distância, vale $E$. Afastando-se a carga de modo que a distância a P
\textbf{dobre}, o campo em P passa a ser:
\begin{alternativas}[3]
\item $E$
\item $E/2$
\item $E/4$
\item $2E$
\item $4E$
\end{alternativas}
\resposta{(c)}
\resolucao{Como $E \propto \dfrac{1}{d^2}$, dobrar a distância divide o campo por
$2^2 = 4$: o novo valor é $E/4$.}
\end{questao}

\blocoaplicacao

\begin{questao}[2]{}
Uma carga de prova de $\SI{e-9}{\coulomb}$, colocada em um ponto P de um campo
elétrico, fica sujeita a uma força de $\SI{e-2}{\newton}$, vertical e para baixo.
Determine:
\begin{itensq}
\item a intensidade, direção e sentido do campo em P;
\item a força que atuaria sobre uma carga de $\SI{3}{\milli\coulomb}$ colocada em P.
\end{itensq}
\resposta{(a) $E=\pot{1.0}{7}\,\si{\newton\per\coulomb}$, vertical para baixo;
(b) $F = \pot{3}{4}\,\si{\newton}$}
\resolucao{(a) $E = \dfrac{F}{q} = \dfrac{\pot{1}{-2}}{\pot{1}{-9}}
= \pot{1.0}{7}\,\si{\newton\per\coulomb}$. Como a carga de prova é positiva, o
campo tem o \emph{mesmo} sentido da força: vertical, para baixo.\\
(b) $F = qE = \pot{3}{-3}\cdot\pot{1}{7} = \pot{3}{4}\,\si{\newton}$,
também vertical para baixo (carga positiva).}
\end{questao}

\begin{questao}[2]{}
Uma carga positiva $Q = \SI{4.0}{\micro\coulomb}$ está no ar. Determine a
intensidade, a direção e o sentido do campo elétrico em um ponto M, a
$\SI{20}{\centi\meter}$ da carga.
\dados{$\ke = \SI{9e9}{\newton\meter\squared\per\coulomb\squared}$.}
\resposta{$E = \pot{9.0}{5}\,\si{\newton\per\coulomb}$, radial, afastando-se de $Q$}
\resolucao{$E = \ke\dfrac{Q}{d^2}
= 9\cdot10^9\cdot\dfrac{\pot{4}{-6}}{(0{,}2)^2} = \pot{9.0}{5}\,\si{\newton\per\coulomb}$.
A direção é a da reta que une $Q$ a M, e o sentido é \emph{para fora} de $Q$
(carga positiva).}
\end{questao}

\begin{questao}[2]{}
Duas cargas puntiformes, $Q$ e $4Q$ (de mesmo sinal), estão fixas no vácuo. O
módulo do campo elétrico resultante é \textbf{nulo} em um ponto sobre a reta que
as une. Esse ponto está tal que sua distância à carga menor, $d_1$, e à carga
maior, $d_2$, satisfazem:
\begin{alternativas}[2]
\item $d_2 = d_1$
\item $d_2 = 2d_1$
\item $d_2 = 4d_1$
\item $d_1 = 2d_2$
\item $d_1 = 4d_2$
\end{alternativas}
\resposta{(b)}
\resolucao{O campo se anula \emph{entre} as cargas, onde os módulos se igualam:
\[
\ke\frac{Q}{d_1^2}=\ke\frac{4Q}{d_2^2}
\ \Rightarrow\
\frac{d_2^2}{d_1^2}=4
\ \Rightarrow\
d_2 = 2d_1.
\]
O ponto de campo nulo fica \emph{mais próximo} da carga menor.}
\end{questao}

\begin{questao}[2]{}
Duas cargas de mesmo módulo $Q$ estão fixas, separadas por $\SI{20}{\centi\meter}$.
Analise o campo elétrico resultante no ponto médio M do segmento que as une, em
dois casos:

\circuitoq{%
\begin{tikzpicture}[scale=0.9,line width=0.8pt,>=Stealth]
 % caso 1: +Q e -Q
 \begin{scope}
   \shade[ball color=Vcol!25] (-1.6,0) circle (0.25); \node at (-1.6,0){\footnotesize$+$};
   \shade[ball color=Icol!25] (1.6,0) circle (0.25);   \node at (1.6,0){\footnotesize$-$};
   \fill[laranja] (0,0) circle (0.06); \node[above,font=\footnotesize] at (0,0.12){M};
   \draw[Vcol,->,thick] (0,-0.5)--(1.1,-0.5); \node[Vcol,font=\footnotesize] at (0.55,-0.8){$\vec E_1$};
   \draw[Vcol,->,thick] (0,-1.1)--(1.1,-1.1); \node[Vcol,font=\footnotesize] at (0.55,-1.4){$\vec E_2$};
   \node[font=\footnotesize] at (0,1.0){Caso I: $+Q$ e $-Q$};
 \end{scope}
\end{tikzpicture}}{Campo no ponto médio entre duas cargas.}

No \textbf{Caso I} ($+Q$ e $-Q$) e no \textbf{Caso II} (ambas $+Q$), o campo
resultante em M é, respectivamente:
\begin{alternativas}
\item nulo em I e nulo em II.
\item não-nulo em I e nulo em II.
\item nulo em I e não-nulo em II.
\item não-nulo em ambos.
\item nulo em ambos.
\end{alternativas}
\resposta{(b)}
\resolucao{Em M, cada carga produz um campo de mesmo módulo (mesma distância,
mesmo $|Q|$).\\
\textbf{Caso I ($+Q$, $-Q$):} os dois campos apontam no \emph{mesmo} sentido (para
longe da positiva e em direção à negativa), então se \textbf{somam}: campo
resultante \emph{não-nulo}, $E = 2\cdot\dfrac{\ke Q}{d^2}$.\\
\textbf{Caso II ($+Q$, $+Q$):} os campos apontam em sentidos \emph{opostos} e se
\textbf{cancelam}: campo resultante \emph{nulo}.\\
Cuidado: para o \emph{potencial} a conclusão se inverte (nulo no caso I, não-nulo
no caso II), porque potencial é escalar e campo é vetorial.}
\end{questao}

\blocoaprofundamento

\begin{questao}[3]{}
As cargas $q_1 = \SI{20}{\coulomb}$ e $q_2 = \SI{64}{\coulomb}$ estão fixas no
vácuo nos pontos A e B, separados por $\SI{9}{\meter}$. Determine a posição, sobre
o segmento AB, onde o campo elétrico resultante é nulo.
\dados{$\ke = \SI{9e9}{\newton\meter\squared\per\coulomb\squared}$.}
\resposta{A $\SI{4}{\meter}$ de A (e $\SI{5}{\meter}$ de B)}
\resolucao{Como as cargas têm o mesmo sinal, o campo se anula entre elas. Seja $x$
a distância a partir de A; a distância a B é $9 - x$. Igualando os módulos:
\[
\ke\frac{q_1}{x^2}=\ke\frac{q_2}{(9-x)^2}
\ \Rightarrow\
\frac{20}{x^2}=\frac{64}{(9-x)^2}.
\]
Extraindo a raiz (grandezas positivas):
$\dfrac{9-x}{x}=\sqrt{\dfrac{64}{20}}=\dfrac{8}{\sqrt{20}}=\dfrac{8}{2\sqrt5}
=\dfrac{4}{\sqrt5}$. Resolvendo,
$9 - x = \dfrac{4}{\sqrt5}x \Rightarrow x = \dfrac{9}{1+4/\sqrt5}\approx\SI{4}{\meter}$.
O ponto neutro fica a cerca de \SI{4}{\meter} de A e \SI{5}{\meter} de B --- mais
perto da carga menor, como sempre ocorre.}
\end{questao}

\begin{questao}[3]{}
A intensidade do campo elétrico de uma carga puntiforme positiva, em função da
distância $d$, obedece ao gráfico $E \times d$ abaixo (curva do tipo $E = c/d^2$).
Sabendo que a $\SI{3}{\meter}$ o campo vale $\SI{80}{\newton\per\coulomb}$,
determine o campo a $\SI{6}{\meter}$.

\circuitoqq{%
\begin{tikzpicture}
 \begin{axis}[width=9cm,height=5cm,xlabel={$d$ (m)},ylabel={$E$ (N/C)},
   xmin=0,xmax=8,ymin=0,ymax=90,xtick={0,2,3,4,6,8},ytick={0,20,40,60,80},
   grid=both,axis lines=left,domain=1.6:8,samples=100]
   \addplot[Vcol,very thick] {720/(x^2)};
   \addplot[only marks,mark=*,Vcol] coordinates {(3,80)(6,20)};
   \node[Vcol,anchor=south west,font=\footnotesize] at (axis cs:3,80) {$(3,\,80)$};
   \node[Vcol,anchor=south west,font=\footnotesize] at (axis cs:6,20) {$(6,\,20)$};
 \end{axis}
\end{tikzpicture}}{Campo elétrico em função da distância ($E \propto 1/d^2$).}

\resposta{$E = \SI{20}{\newton\per\coulomb}$}
\resolucao{Como $E \propto \dfrac{1}{d^2}$, dobrar a distância (de \SI{3}{} para
\SI{6}{\meter}) divide o campo por 4:
\[
E(6)=\frac{E(3)}{4}=\frac{80}{4}=\SI{20}{\newton\per\coulomb}.
\]
Formalmente, $E(3)\cdot3^2 = E(6)\cdot6^2$, ou seja,
$80\cdot9 = E(6)\cdot36 \Rightarrow E(6) = \SI{20}{\newton\per\coulomb}$.}
\end{questao}

\begin{questao}[3]{}
Duas cargas $q_1 = \SI{+9}{\micro\coulomb}$ e $q_2 = \SI{+4}{\micro\coulomb}$
estão fixas e separadas por $\SI{10}{\centi\meter}$. Determine a que distância da
carga $q_1$, sobre a reta que as une, o campo elétrico resultante é nulo.
\resposta{A $\SI{6}{\centi\meter}$ de $q_1$}
\resolucao{Cargas de mesmo sinal: o campo se anula \emph{entre} elas. Seja $x$ a
distância a $q_1$; a distância a $q_2$ é $(10 - x)$. Igualando os módulos dos
campos:
\[
\frac{\ke q_1}{x^2}=\frac{\ke q_2}{(10-x)^2}
\ \Rightarrow\
\frac{9}{x^2}=\frac{4}{(10-x)^2}.
\]
Extraindo a raiz: $\dfrac{3}{x}=\dfrac{2}{10-x} \Rightarrow 3(10-x)=2x
\Rightarrow 30 = 5x \Rightarrow x = \SI{6}{\centi\meter}$.
O ponto neutro fica mais perto da carga \emph{menor} ($q_2$), a
\SI{4}{\centi\meter} dela.}
\end{questao}

\begin{questao}[3]{}
Duas cargas puntiformes, $+4Q$ e $-Q$, estão fixas no vácuo, separadas por uma
distância $d$.
\begin{itensq}
\item Mostre que existe um ponto sobre a reta que as une onde o campo elétrico
      resultante é \textbf{nulo}, e determine sua posição.
\item Explique por que esse ponto está \emph{fora} do segmento que une as cargas,
      e do lado da carga de \emph{menor} módulo.
\item Nesse mesmo ponto, o potencial elétrico é nulo? Justifique.
\end{itensq}
\resposta{(a) a $2d$ da carga $+4Q$ (isto é, a $d$ além de $-Q$);
(b) e (c) ver comentário}
\resolucao{\textbf{(a)} Seja $x$ a distância medida a partir de $+4Q$, com o ponto
além de $-Q$ (logo a distância a $-Q$ é $x-d$). Igualando os módulos:
\[
\frac{\ke\,4Q}{x^{2}}=\frac{\ke\,Q}{(x-d)^{2}}
\;\Longrightarrow\;
4(x-d)^2 = x^2
\;\Longrightarrow\;
2(x-d)=\pm x.
\]
A raiz fisicamente válida é $2x-2d = x \Rightarrow \boxed{x = 2d}$, ou seja, a
$2d$ da carga maior e a $d$ além da carga menor.
(A outra raiz, $x = \tfrac{2d}{3}$, cai \emph{entre} as cargas, onde os campos
apontam no \emph{mesmo} sentido e portanto não podem se cancelar --- deve ser
descartada.)

\textbf{(b)} Entre cargas de sinais \emph{opostos}, os dois vetores campo apontam
no mesmo sentido (saindo da positiva, entrando na negativa) e se \emph{somam};
nunca se anulam ali. Fora do segmento, eles têm sentidos contrários, e o
cancelamento só é possível onde o campo da carga \emph{menor} --- que está mais
próxima --- consegue igualar o da maior. Por isso o ponto neutro fica sempre do
lado da carga de menor módulo.

\textbf{(c) Não.} Potencial é grandeza \emph{escalar}:
\[
V = \ke\frac{4Q}{2d}+\ke\frac{(-Q)}{d}
= \frac{2\ke Q}{d}-\frac{\ke Q}{d}
= \frac{\ke Q}{d}\neq 0.
\]
\textbf{Conclusão conceitual:} campo nulo e potencial nulo são condições
\emph{independentes}. Neste caso, $\vec E = 0$ mas $V \neq 0$. No ponto médio de um
dipolo simétrico ocorre o oposto: $V = 0$ mas $\vec E \neq 0$.}
\end{questao}

\begin{questao}[3]{}
Um elétron entra com velocidade horizontal $v_0=\SI{2e7}{\meter\per\second}$ na
região entre duas placas paralelas de comprimento $L=\SI{5}{\centi\meter}$, onde
existe campo elétrico uniforme $E=\SI{e4}{\volt\per\meter}$, perpendicular à
velocidade.

\circuitoq{%
\begin{tikzpicture}[scale=1,line width=0.8pt,>=Stealth]
 \draw[cinza!70,line width=2.2pt] (0,1.3)--(4.2,1.3);
 \draw[cinza!70,line width=2.2pt] (0,-1.3)--(4.2,-1.3);
 \node[above,font=\footnotesize] at (2.1,1.35) {$+$\ placa};
 \node[below,font=\footnotesize] at (2.1,-1.35) {$-$\ placa};
 \foreach \x in {0.6,1.5,2.4,3.3}
   \draw[->,Vcol!70] (\x,1.2)--(\x,-1.2);
 \node[Vcol,font=\footnotesize] at (3.85,0.55) {$\vec E$};
 \draw[->,laranja,line width=1.2pt] (-1.4,0)--(0,0)
       node[midway,above,font=\footnotesize]{$v_0$};
 \draw[laranja,line width=1.2pt,domain=0:4.2,samples=40,smooth]
       plot (\x,{0.055*\x*\x});
 \draw[->,laranja,line width=1.2pt] (4.2,{0.055*4.2*4.2})--(5.3,{0.055*4.2*4.2+0.55});
 \draw[<->,cinza] (0,-1.75)--(4.2,-1.75) node[midway,below,font=\footnotesize]{$L$};
 \draw[<->,cinza] (4.5,0)--(4.5,{0.055*4.2*4.2}) node[midway,right,font=\footnotesize]{$y$};
\end{tikzpicture}}{Deflexão de um elétron em campo elétrico uniforme.}

\begin{itensq}
\item Calcule a aceleração do elétron.
\item Determine o tempo de permanência entre as placas e o desvio vertical $y$.
\item Calcule o ângulo de saída em relação à direção inicial.
\end{itensq}
\dados{$e=\SI{1.6e-19}{\coulomb}$; $m_e=\SI{9.11e-31}{\kilogram}$.}
\resposta{(a) $a\approx\SI{1.76e15}{\meter\per\second\squared}$;
(b) $t=\SI{2.5}{\nano\second}$, $y\approx\SI{5.5}{\milli\meter}$;
(c) $\theta\approx\SI{12.4}{\degree}$}
\resolucao{Este é um \textbf{lançamento de projétil} eletrostático: movimento
uniforme na horizontal e uniformemente acelerado na vertical.\\
(a) $a=\dfrac{eE}{m_e}=\dfrac{(\pot{1.6}{-19})(\pot{1}{4})}{\pot{9.11}{-31}}
\approx\pot{1.76}{15}\,\si{\meter\per\second\squared}$
--- cerca de $10^{14}$ vezes a gravidade, o que justifica desprezar o peso.\\
(b) Na horizontal (MU):
$t=\dfrac{L}{v_0}=\dfrac{0{,}05}{\pot{2}{7}}=\pot{2.5}{-9}\,\si{\second}
=\SI{2.5}{\nano\second}$.\\
Na vertical (MUV, partindo do repouso):
\[
y=\tfrac12 a t^{2}=\tfrac12(\pot{1.76}{15})(\pot{2.5}{-9})^{2}
\approx\pot{5.5}{-3}\,\si{\meter}=\SI{5.5}{\milli\meter}.
\]
(c) A velocidade vertical na saída é $v_y = a\,t \approx
\pot{4.4}{6}\,\si{\meter\per\second}$, logo
\[
\tan\theta=\frac{v_y}{v_0}=\frac{\pot{4.4}{6}}{\pot{2}{7}}=0{,}22
\;\Longrightarrow\;
\theta\approx\SI{12.4}{\degree}.
\]
\textbf{Aplicação:} é exatamente esse o mecanismo de deflexão dos antigos tubos de
raios catódicos (TV e osciloscópios) e das impressoras a jato de tinta, que
defletem gotículas carregadas para formar caracteres.}
\end{questao}

\begin{questao}[3]{}
Quatro cargas de módulo $q=\SI{1}{\micro\coulomb}$ ocupam os vértices de um
quadrado de lado $a=\SI{10}{\centi\meter}$. Compare o \textbf{campo elétrico no
centro} em três configurações:
\begin{itensq}
\item todas positivas;
\item duas positivas e duas negativas, alternadas ($+,-,+,-$);
\item duas positivas e duas negativas, adjacentes ($+,+,-,-$).
\end{itensq}
\dados{$\ke=\SI{9e9}{}$; $\sqrt2\approx1{,}41$.}
\resposta{(a) $\vec E=0$; (b) $\vec E=0$; (c) $E\approx\pot{5.1}{6}\,\si{\newton\per\coulomb}$}
\resolucao{Cada carga dista $\dfrac{a\sqrt2}{2}=\SI{0.0707}{\meter}$ do centro,
produzindo individualmente
\[
E_1=\frac{\ke q}{(a\sqrt2/2)^{2}}=\frac{9\cdot10^{9}\cdot\pot{1}{-6}}{0{,}005}
=\pot{1.8}{6}\,\si{\newton\per\coulomb}.
\]
\textbf{(a) Todas positivas.} Os quatro vetores apontam radialmente para fora,
aos pares diametralmente opostos: cancelam-se dois a dois. $\vec E = 0$.\\
\textbf{(b) Alternadas ($+,-,+,-$).} As cargas diametralmente opostas têm o
\emph{mesmo} sinal, e seus campos continuam se cancelando aos pares.
$\vec E = 0$ novamente.\\
\textbf{(c) Adjacentes ($+,+,-,-$).} Agora cada par diametralmente oposto tem
sinais \emph{contrários}: os campos se \textbf{somam} em vez de cancelar. Os dois
pares contribuem na mesma direção (perpendicular à linha que separa as positivas
das negativas):
\[
E=2\,(E_1\sqrt2)=2\sqrt2\,E_1\approx 2{,}83\cdot\pot{1.8}{6}
\approx\pot{5.1}{6}\,\si{\newton\per\coulomb}.
\]
\textbf{Lição central.} Em (a) e (b), configurações \emph{muito} diferentes
produzem o mesmo resultado nulo. Em (c), com as mesmas quatro cargas, apenas
reposicionadas, o campo é máximo. \emph{O campo resultante depende da geometria
tanto quanto dos valores das cargas} --- e a simetria é a ferramenta que revela
isso sem nenhum cálculo pesado.}
\end{questao}

\begin{questao}[3]{}
Uma carga $-Q$ está na origem e uma carga $+4Q$ está em $x=d$.
\begin{itensq}
\item Determine o ponto sobre o eixo $x$ onde o campo resultante é nulo.
\item Explique por que ele fica \emph{fora} do segmento e do lado da carga de
      menor módulo.
\item Nesse ponto, quanto vale o potencial elétrico?
\end{itensq}
\resposta{(a) $x=-d$ (a uma distância $d$ à esquerda de $-Q$);
(b) e (c) ver comentário; $V=\dfrac{\ke Q}{d}$}
\resolucao{(a) Como as cargas têm sinais \emph{opostos}, entre elas os campos
apontam no mesmo sentido e se somam --- o cancelamento só pode ocorrer fora do
segmento. Seja $x<0$ a posição procurada; a distância a $-Q$ é $|x|$ e a $+4Q$ é
$|x|+d$:
\[
\frac{\ke Q}{|x|^{2}}=\frac{\ke\,4Q}{(|x|+d)^{2}}
\;\Longrightarrow\;
\frac{|x|+d}{|x|}=2
\;\Longrightarrow\;
|x|=d.
\]
O ponto está em $x=-d$.\\
(b) Do lado da carga \emph{maior}, seu campo sempre domina (está mais perto
\emph{e} é mais intensa) --- nunca poderia ser equilibrado. Do lado da carga
menor, a proximidade compensa a menor intensidade, e existe exatamente um ponto
de equilíbrio.\\
(c) O potencial é escalar:
\[
V=\ke\frac{(-Q)}{d}+\ke\frac{4Q}{2d}
=-\frac{\ke Q}{d}+\frac{2\ke Q}{d}
=\frac{\ke Q}{d}\neq0.
\]
\textbf{Mais um exemplo de que $\vec E=0$ não implica $V=0$} --- e vice-versa. O
campo mede a \emph{taxa de variação} do potencial, não o seu valor.}
\end{questao}

\begin{questao}[3]{}
Um anel isolante de raio $R$ está uniformemente carregado com carga total $Q$.
Considere um ponto P sobre o eixo do anel, a uma distância $x$ do centro.
\begin{itensq}
\item Explique, por simetria, por que o campo em P é paralelo ao eixo.
\item Sabendo que $E(x)=\dfrac{\ke\,Q\,x}{\left(x^{2}+R^{2}\right)^{3/2}}$,
      determine $E$ no centro do anel e no limite $x\gg R$.
\item Mostre que o campo é \emph{máximo} em $x = \dfrac{R}{\sqrt2}$.
\end{itensq}
\resposta{(b) $E(0)=0$ e $E\to\ke Q/x^{2}$; (c) demonstração}
\resolucao{(a) \textbf{Argumento de simetria.} Para cada elemento de carga do
anel existe outro diametralmente oposto. As componentes de campo
\emph{perpendiculares} ao eixo desses dois elementos são iguais e opostas,
cancelando-se. Restam apenas as componentes \emph{paralelas} ao eixo, que se
somam. Logo $\vec E$ é axial.

(b) \textbf{No centro} ($x=0$): substituindo,
\[
E(0)=\frac{\ke Q\cdot 0}{R^{3}}=0.
\]
Faz sentido: no centro, as contribuições de elementos opostos se cancelam
\emph{integralmente}.\\
\textbf{Longe} ($x\gg R$): então $\left(x^{2}+R^{2}\right)^{3/2}\approx x^{3}$ e
\[
E\approx\frac{\ke Q x}{x^{3}}=\frac{\ke Q}{x^{2}},
\]
o campo de uma carga puntiforme --- como esperado, pois de longe o anel "vira"
um ponto.

(c) \textbf{Máximo.} Derivando $E(x)$ e igualando a zero:
\[
\frac{\mathrm{d}E}{\mathrm{d}x}
=\ke Q\,\frac{\left(x^{2}+R^{2}\right)^{3/2}-x\cdot\tfrac32\left(x^{2}+R^{2}\right)^{1/2}(2x)}
{\left(x^{2}+R^{2}\right)^{3}}=0.
\]
O numerador, dividido por $\left(x^{2}+R^{2}\right)^{1/2}$, dá
\[
\left(x^{2}+R^{2}\right)-3x^{2}=0
\;\Longrightarrow\;
R^{2}=2x^{2}
\;\Longrightarrow\;
\boxed{\,x=\frac{R}{\sqrt2}\,}
\]
\textbf{Interpretação física do máximo.} Perto do centro, as contribuições
axiais são pequenas (os vetores são quase perpendiculares ao eixo). Muito longe,
a distância cresce e o campo decai. Entre os dois regimes existe um ponto ótimo,
a $x=R/\sqrt2\approx0{,}71R$.\\
\textbf{Onde isso é usado:} essa geometria é a base das \emph{bobinas de
Helmholtz} (versão magnética do mesmo cálculo) e dos anéis focalizadores em
aceleradores de partículas e microscópios eletrônicos.}
\end{questao}

\blocodesafio

\begin{questao}[4]{}
Uma haste isolante fina de comprimento $L$ está uniformemente carregada com carga
total $Q$ (densidade linear $\lambda=Q/L$). Considere um ponto P sobre o
prolongamento da haste, a uma distância $a$ da extremidade mais próxima.
\begin{itensq}
\item Monte a integral que fornece o campo elétrico em P.
\item Mostre que $E=\dfrac{\ke\,Q}{a\,(a+L)}$.
\item Verifique o comportamento no limite $a \gg L$ e interprete.
\end{itensq}
\resposta{(b) demonstração; (c) $E\to \ke Q/a^{2}$ (carga puntiforme)}
\resolucao{(a) Tome um elemento de comprimento $\mathrm{d}x$ à distância $x$ de P
(com $x$ variando de $a$ até $a+L$). Sua carga é
$\mathrm{d}q=\lambda\,\mathrm{d}x$ e contribui com
\[
\mathrm{d}E=\frac{\ke\,\mathrm{d}q}{x^{2}}=\frac{\ke\lambda\,\mathrm{d}x}{x^{2}}.
\]
Como todos os elementos estão alinhados com P, as contribuições têm a
\emph{mesma direção} e somam-se escalarmente:
\[
E=\int_{a}^{a+L}\frac{\ke\lambda}{x^{2}}\,\mathrm{d}x.
\]
(b) Integrando:
\[
E=\ke\lambda\left[-\frac{1}{x}\right]_{a}^{a+L}
=\ke\lambda\left(\frac{1}{a}-\frac{1}{a+L}\right)
=\ke\lambda\,\frac{L}{a(a+L)}.
\]
Como $\lambda L = Q$:
\[
\boxed{\;E=\frac{\ke\,Q}{a\,(a+L)}\;}
\]
(c) \textbf{Limite $a\gg L$:} então $a+L\approx a$ e
\[
E\approx\frac{\ke Q}{a^{2}},
\]
que é exatamente o campo de uma \textbf{carga puntiforme} $Q$.\\
\textbf{Interpretação:} de longe, qualquer distribuição finita de carga
"parece" puntiforme. Esse teste de consistência --- verificar se a expressão
obtida reproduz o caso conhecido em algum limite --- é uma das melhores formas de
validar um resultado. \emph{Se o limite não fechasse, haveria erro na
integração.}\\
\textbf{Limite oposto ($a\ll L$):} $E\approx\dfrac{\ke Q}{aL}=\dfrac{\ke\lambda}{a}$
--- campo de um fio "infinito" visto de perto, decaindo com $1/a$ em vez de
$1/a^{2}$.}
\end{questao}

\begin{questao}[4]{}
Um dipolo elétrico é formado por cargas $+q$ e $-q$ separadas por uma pequena
distância $2\ell$. Define-se o \textbf{momento de dipolo} $p=q\,(2\ell)$.
\begin{itensq}
\item Mostre que, sobre a mediatriz do dipolo, a uma distância $r\gg\ell$ do
      centro, o campo vale $E\approx\dfrac{\ke\,p}{r^{3}}$.
\item Por que o campo do dipolo decai mais rápido que o de uma carga isolada?
\item Um dipolo colocado em campo uniforme sofre força resultante? E torque?
\end{itensq}
\resposta{(a) demonstração; (b) e (c) ver comentário}
\resolucao{(a) Sobre a mediatriz, cada carga dista
$d=\sqrt{r^{2}+\ell^{2}}$ do ponto e produz campo de módulo
$E_1=\dfrac{\ke q}{r^{2}+\ell^{2}}$. As componentes \emph{paralelas} à mediatriz
se cancelam (simetria); restam as componentes ao longo do eixo do dipolo, cada uma
valendo $E_1\cos\alpha$, com $\cos\alpha=\dfrac{\ell}{\sqrt{r^{2}+\ell^{2}}}$:
\[
E=2E_1\cos\alpha
=\frac{2\ke q\,\ell}{\left(r^{2}+\ell^{2}\right)^{3/2}}.
\]
Para $r\gg\ell$, $\left(r^{2}+\ell^{2}\right)^{3/2}\approx r^{3}$, e como
$p=2q\ell$:
\[
\boxed{\;E\approx\frac{\ke\,p}{r^{3}}\;}
\]
(b) Porque a carga \emph{total} do dipolo é nula. De longe, os campos das duas
cargas quase se cancelam --- sobra apenas o efeito da pequena separação entre
elas, que é de ordem superior. Daí o decaimento $1/r^{3}$ em vez de $1/r^{2}$.\\
(c) Em campo \textbf{uniforme}:
\begin{itemize}[leftmargin=1.6em,itemsep=1pt,topsep=2pt]
\item \textbf{Força resultante nula}: as forças $+qE$ e $-qE$ sobre as duas cargas
      têm mesmo módulo e sentidos opostos.
\item \textbf{Torque não nulo}: como elas atuam em \emph{pontos diferentes},
      formam um binário de momento
      \[
      \tau = p\,E\sin\theta,
      \]
      que tende a \emph{alinhar} o dipolo com o campo.
\end{itemize}
\textbf{Aplicações:} é esse torque que alinha moléculas polares (como a água) em
um campo --- princípio do forno de micro-ondas, onde o campo oscilante faz as
moléculas girarem e o atrito molecular gera calor. Em campo \emph{não} uniforme,
surge também força resultante --- razão pela qual um objeto neutro é atraído por
um corpo eletrizado.}
\end{questao}

% BEGIN BANCO COMPLEMENTAR Campo elétrico
% =====================================================================
%  Banco complementar --- Campo Eletrico  (REVISADO)
%  Substitui a versao gerada por template (8 clones em N1, 11 em N2,
%  7 em N3 e 8 questoes conceituais marcadas como N4).
%  O cap01b e denso e ja cobre: E de carga puntiforme, proporcionalidade
%  com a distancia, carga de prova com forca dada, direcao e sentido do
%  campo, PONTO DE CAMPO NULO (tres vezes: Q e 4Q; 9 e 4 uC; -Q e +4Q),
%  duas cargas de mesmo modulo a 20 cm, grafico E x d, ELETRON ENTRE
%  PLACAS PARALELAS com deflexao, quadrado de quatro cargas, anel
%  isolante no eixo, haste finita por integracao (N4) e momento dipolar
%  (N4).
%  Este banco evita esses casos e explora: E=0 com V nao nulo, linhas de
%  campo e por que nao se cruzam, PROJETO de placas paralelas, cavidade
%  em condutor, tracado numerico de linhas, esfera isolante macica por
%  Gauss, par de placas com +-sigma e o limite imposto pela rigidez
%  dieletrica.
%  Distribuicao mantida: 8 N1 + 11 N2 + 7 N3 + 8 N4 = 34
% =====================================================================

\subsubsection*{Banco complementar --- Campo elétrico}

\begin{atencao}[Organização do banco]
O campo é \textbf{vetorial}: contribuições somam-se por componentes, e podem
cancelar-se. Ele é definido por $\vec E=\vec F/q_0$, mas \emph{não depende} da
carga de prova --- é propriedade do ponto. Duas expressões dominam o tópico:
$E=\ke Q/r^{2}$ para carga puntiforme e $E=\sigma/\varepsilon_0$ para plano
carregado, esta última \emph{independente da distância}. O N4 trata de projeto,
condutores, distribuições contínuas e limites físicos.
\end{atencao}

\blocofixacao

\begin{questao}[1]{}
Determine o módulo do campo elétrico produzido por $Q=+\SI{3}{\micro\coulomb}$
a $\SI{30}{\centi\meter}$ de distância.
\dados{$\ke=\SI{9e9}{\newton\meter\squared\per\coulomb\squared}$}
\resposta{$E=\SI{300}{\kilo\newton\per\coulomb}$}
\resolucao{
\[
E=\frac{\ke Q}{r^{2}}=\frac{\pot{9}{9}\times\pot{3}{-6}}{(0{,}30)^{2}}
=\frac{\pot{2.7}{4}}{0{,}09}=\pot{3}{5}\,\si{\newton\per\coulomb}.
\]
O campo aponta \textbf{para longe} da carga, por ela ser positiva.\\
\textbf{Atenção à potência:} o campo cai com $1/r^{2}$, o potencial com
$1/r$. Dobrar a distância divide $E$ por quatro e $V$ por dois.}
\end{questao}

\begin{questao}[1]{}
Determine o campo produzido por $Q=-\SI{5}{\micro\coulomb}$ a
$\SI{50}{\centi\meter}$, indicando o sentido.
\resposta{$E=\SI{180}{\kilo\newton\per\coulomb}$, apontando \emph{para} a carga}
\resolucao{
\[
E=\frac{\pot{9}{9}\times\pot{5}{-6}}{0{,}25}
=\pot{1.8}{5}\,\si{\newton\per\coulomb}.
\]
\textbf{Sentido:} cargas negativas \emph{atraem} cargas de prova positivas, logo
o campo aponta \textbf{para} a carga fonte.\\
\textbf{Regra mnemônica:} o campo \emph{sai} do positivo e \emph{entra} no
negativo --- é a mesma regra que orienta as linhas de campo.}
\end{questao}

\begin{questao}[1]{}
Uma carga $q=+\SI{2}{\micro\coulomb}$ é colocada em um ponto onde
$E=\SI{5000}{\newton\per\coulomb}$. Determine a força sobre ela.
\resposta{$F=\SI{10}{\milli\newton}$, no sentido do campo}
\resolucao{
\[
F=qE=\pot{2}{-6}\times5000=\pot{1}{-2}\,\si{\newton}=\SI{10}{\milli\newton}.
\]
Como a carga é positiva, a força tem o \emph{mesmo} sentido de $\vec E$. Se
fosse negativa, teria sentido oposto --- e apenas o sentido mudaria, não o
módulo.}
\end{questao}

\begin{questao}[1]{}
Determine a carga que produz $E=\SI{2000}{\newton\per\coulomb}$ a
$\SI{30}{\centi\meter}$.
\resposta{$Q=\SI{20}{\nano\coulomb}$}
\resolucao{
\[
Q=\frac{Er^{2}}{\ke}=\frac{2000\times0{,}09}{\pot{9}{9}}
=\pot{2}{-8}\,\si{\coulomb}=\SI{20}{\nano\coulomb}.
\]
Uma carga de nanocoulombs --- ordem de grandeza típica da eletrização por
atrito --- já produz milhares de newtons por coulomb a trinta centímetros.}
\end{questao}

\begin{questao}[1]{}
A que distância de $Q=+\SI{8}{\micro\coulomb}$ o campo vale
$\SI{200}{\kilo\newton\per\coulomb}$?
\resposta{$r=\SI{0.6}{\meter}$}
\resolucao{
\[
r=\sqrt{\frac{\ke Q}{E}}=\sqrt{\frac{\pot{9}{9}\times\pot{8}{-6}}{\pot{2}{5}}}
=\sqrt{0{,}36}=\SI{0.6}{\meter}.
\]
\textbf{Superfície de campo constante:} todos os pontos a
$\SI{0.6}{\meter}$ da carga têm o mesmo \emph{módulo} de campo --- mas
\emph{direções} diferentes, pois o campo é radial. Diferentemente das
equipotenciais, essas superfícies não têm significado especial.}
\end{questao}

\begin{questao}[1]{}
O campo elétrico em um ponto, medido com carga de prova $q_0$, vale $E$. Se a
carga de prova for \textbf{dobrada}, o valor medido:
\begin{alternativas}[2]
\item dobra
\item cai à metade
\item não se altera
\item depende do sinal de $q_0$
\end{alternativas}
\resposta{(c)}
\resolucao{Por definição, $\vec E=\vec F/q_0$. Dobrar $q_0$ dobra a força
\emph{e} o denominador --- a razão permanece.\\
\textbf{Ponto conceitual:} o campo é uma propriedade \textbf{do ponto do
espaço}, criada pelas cargas-fonte. A carga de prova apenas o \emph{revela};
ela não o cria nem o modifica.\\
\textbf{Ressalva real:} se $q_0$ for grande demais, ela perturba as
cargas-fonte (por indução, em condutores) e a medida deixa de valer. Daí a
definição rigorosa usar o limite $q_0\to0$.}
\end{questao}

\begin{questao}[1]{}
Mostre que $\si{\newton\per\coulomb}$ e $\si{\volt\per\meter}$ são a mesma
unidade.
\resposta{$\si{\volt}=\si{\joule\per\coulomb}$ e
$\si{\joule}=\si{\newton\meter}$}
\resolucao{
\[
\frac{\si{\volt}}{\si{\meter}}
=\frac{\si{\joule}}{\si{\coulomb}\cdot\si{\meter}}
=\frac{\si{\newton}\cdot\si{\meter}}{\si{\coulomb}\cdot\si{\meter}}
=\frac{\si{\newton}}{\si{\coulomb}}.\ \checkmark
\]
\textbf{As duas leituras do campo.} $\si{\newton\per\coulomb}$ enfatiza a
\emph{força} por unidade de carga; $\si{\volt\per\meter}$ enfatiza a
\emph{taxa de queda do potencial} por unidade de comprimento. São descrições
equivalentes de $\vec E=-\nabla V$.\\
Na prática, usa-se $\si{\volt\per\meter}$ quando há tensões envolvidas (placas,
cabos) e $\si{\newton\per\coulomb}$ em problemas de dinâmica de partículas.}
\end{questao}

\begin{questao}[1]{}
Uma carga negativa é colocada em uma região de campo elétrico uniforme. Sobre
seu movimento inicial, é correto afirmar que ela:
\begin{alternativas}[2]
\item acelera no sentido do campo
\item acelera no sentido oposto ao campo
\item permanece em repouso
\item move-se perpendicularmente ao campo
\end{alternativas}
\resposta{(b)}
\resolucao{A força vale $\vec F=q\vec E$ com $q<0$, logo $\vec F$ é
\textbf{antiparalela} a $\vec E$.\\
\textbf{Consequência energética:} a carga negativa move-se espontaneamente no
sentido de $V$ \emph{crescente} --- o oposto de uma carga positiva. Em ambos os
casos, a energia potencial $U=qV$ \emph{diminui}, como deve ocorrer em qualquer
movimento espontâneo.}
\end{questao}

\blocoaplicacao

\begin{questao}[2]{}
Uma carga $q_0=-\SI{1}{\micro\coulomb}$ é colocada em um campo uniforme de
$E=\SI{2000}{\newton\per\coulomb}$.
\begin{itensq}
\item Determine a força.
\item Determine sua aceleração, se a massa for $\SI{4}{\gram}$.
\end{itensq}
\resposta{$F=\SI{2}{\milli\newton}$, oposta a $\vec E$;
$a=\SI{0.5}{\meter\per\second\squared}$}
\resolucao{(a) $F=|q|E=\pot{1}{-6}(2000)=\pot{2}{-3}\,\si{\newton}$, com sentido
\emph{oposto} ao campo.\\
(b)
\[
a=\frac{F}{m}=\frac{\pot{2}{-3}}{\pot{4}{-3}}=\SI{0.5}{\meter\per\second\squared}.
\]
\textbf{Comparação com a gravidade:} $a$ é apenas $5\%$ de $g$. Para objetos
macroscópicos, a força elétrica só compete com o peso se a carga for
relativamente grande ou a massa muito pequena --- daí os experimentos
eletrostáticos clássicos usarem pedaços de papel e gotas de óleo.}
\end{questao}

\begin{questao}[2]{}
Um próton é abandonado em um campo uniforme de
$E=\SI{1e4}{\newton\per\coulomb}$. Determine sua aceleração.
\dados{$e=\pot{1.6}{-19}\,\si{\coulomb}$; $m_p=\pot{1.67}{-27}\,\si{\kilogram}$}
\resposta{$a=\pot{9.6}{11}\,\si{\meter\per\second\squared}$}
\resolucao{
\[
a=\frac{eE}{m_p}=\frac{(\pot{1.6}{-19})(\pot{1}{4})}{\pot{1.67}{-27}}
=\pot{9.58}{11}\,\si{\meter\per\second\squared}.
\]
\textbf{Cerca de $10^{11}$ vezes $g$} --- a gravidade é completamente
desprezível na dinâmica de partículas carregadas.\\
\textbf{Para o elétron}, a aceleração seria $1836$ vezes maior
($\pot{1.76}{15}\,\si{\meter\per\second\squared}$), pela razão das massas. É por
isso que, em tubos e aceleradores, são os elétrons que respondem primeiro.}
\end{questao}

\begin{questao}[2]{}
Duas cargas estão sobre o eixo $x$: $+\SI{4}{\micro\coulomb}$ em $x=0$ e
$-\SI{2}{\micro\coulomb}$ em $x=\SI{0.30}{\meter}$. Determine o campo em
$x=\SI{0.60}{\meter}$.
\resposta{$E=\SI{100}{\kilo\newton\per\coulomb}$, no sentido $-x$}
\resolucao{
\[
E_1=\frac{\pot{9}{9}(\pot{4}{-6})}{(0{,}60)^{2}}
=\pot{1}{5}\,\si{\newton\per\coulomb}\ (\text{sentido }+x);
\]
\[
E_2=\frac{\pot{9}{9}(\pot{2}{-6})}{(0{,}30)^{2}}
=\pot{2}{5}\,\si{\newton\per\coulomb}\ (\text{sentido }-x).
\]
\[
E=2\times10^{5}-1\times10^{5}=\pot{1}{5}\,\si{\newton\per\coulomb}\ (-x).
\]
\textbf{Observe:} a carga \emph{menor} produziu o campo \emph{maior}, por estar
duas vezes mais próxima --- e $1/r^{2}$ vence o fator $2$ de carga.
Compare os fatores: carga $\times\frac12$, distância$^{2}$
$\times\frac14$.}
\end{questao}

\begin{questao}[2]{}
Duas cargas iguais $Q=+\SI{2}{\micro\coulomb}$ estão em
$(\pm\SI{0.30}{\meter},\,0)$. Determine o campo no ponto
$(0,\ \SI{0.40}{\meter})$.
\resposta{$E=\SI{115.2}{\kilo\newton\per\coulomb}$, ao longo de $+y$}
\resolucao{Cada carga dista $r=\sqrt{0{,}30^{2}+0{,}40^{2}}=\SI{0.5}{\meter}$
(triângulo $3$--$4$--$5$):
\[
E_1=E_2=\frac{\pot{9}{9}(\pot{2}{-6})}{0{,}25}
=\pot{7.2}{4}\,\si{\newton\per\coulomb}.
\]
\textbf{Por simetria}, as componentes horizontais se cancelam e as verticais se
somam, cada uma multiplicada por $\cos\theta=0{,}40/0{,}50=0{,}8$:
\[
E=2(\pot{7.2}{4})(0{,}8)=\pot{1.152}{5}\,\si{\newton\per\coulomb}.
\]
\textbf{Uso da simetria:} reconhecê-la elimina metade do trabalho. Sempre que a
configuração for simétrica em relação ao ponto de interesse, identifique quais
componentes se anulam \emph{antes} de calcular.}
\end{questao}

\begin{questao}[2]{}
Duas placas paralelas separadas por $\SI{2}{\milli\meter}$ estão submetidas a
$\SI{600}{\volt}$. Determine o campo entre elas.
\resposta{$E=\SI{300}{\kilo\volt\per\meter}$}
\resolucao{
\[
E=\frac{U}{d}=\frac{600}{\pot{2}{-3}}=\pot{3}{5}\,\si{\volt\per\meter}.
\]
\textbf{Campo uniforme:} entre placas extensas e próximas, $E$ é praticamente o
mesmo em todos os pontos --- e independe da posição. É a configuração de
referência para acelerar e defletir partículas.\\
\textbf{Verificação de segurança:} $\SI{300}{\kilo\volt\per\meter}$ é um décimo
da rigidez do ar ($\SI{3}{\mega\volt\per\meter}$) --- há folga confortável.
Reduzir a separação para $\SI{0.2}{\milli\meter}$ com a mesma tensão provocaria
ruptura.}
\end{questao}

\begin{questao}[2]{}
Uma placa condutora extensa tem densidade superficial de carga
$\sigma=\SI{2}{\micro\coulomb\per\meter\squared}$. Determine o campo próximo à
sua superfície.
\dados{$\varepsilon_0=\pot{8.85}{-12}\,\si{\farad\per\meter}$}
\resposta{$E=\SI{226}{\kilo\newton\per\coulomb}$}
\resolucao{Para um condutor em equilíbrio,
\[
E=\frac{\sigma}{\varepsilon_0}=\frac{\pot{2}{-6}}{\pot{8.85}{-12}}
=\pot{2.26}{5}\,\si{\newton\per\coulomb}.
\]
\textbf{Fato notável:} o campo \textbf{não depende da distância} à placa (desde
que ela seja extensa e o ponto próximo). Diferentemente da carga puntiforme, não
há queda com $1/r^{2}$.\\
\textbf{Cuidado com a fórmula:} para um \emph{plano isolante} com carga nas duas
faces, o campo vale $\sigma/2\varepsilon_0$ --- metade. A diferença é que no
condutor toda a carga fica em uma face, e o campo interno é nulo.}
\end{questao}

\begin{questao}[2]{}
Compare o campo produzido por $Q_1=+\SI{9}{\micro\coulomb}$ a
$\SI{30}{\centi\meter}$ com o de $Q_2=+\SI{4}{\micro\coulomb}$ a
$\SI{20}{\centi\meter}$.
\resposta{$E_1=E_2=\SI{900}{\kilo\newton\per\coulomb}$ --- iguais}
\resolucao{
\[
E_1=\frac{\pot{9}{9}(\pot{9}{-6})}{0{,}09}=\pot{9}{5};
\qquad
E_2=\frac{\pot{9}{9}(\pot{4}{-6})}{0{,}04}=\pot{9}{5}\,\si{\newton\per\coulomb}.
\]
\textbf{Coincidência não é acaso:} as razões $Q/r^{2}$ são iguais, pois
$9/9=4/4$. Sempre que $Q\propto r^{2}$, o campo se mantém.\\
\textbf{Consequência:} se essas duas cargas estivessem em lados opostos de um
ponto, o campo resultante ali seria \emph{nulo} --- e essa é exatamente a
condição do "ponto de campo nulo" entre cargas de mesmo sinal.}
\end{questao}

\begin{questao}[2]{}
Uma carga em um ponto sofre campos de $\SI{3}{\kilo\newton\per\coulomb}$ na
direção $x$ e $\SI{4}{\kilo\newton\per\coulomb}$ na direção $y$, produzidos por
duas fontes distintas. Determine o campo resultante.
\resposta{$E=\SI{5}{\kilo\newton\per\coulomb}$, a $\ang{53.1}$ do eixo $x$}
\resolucao{Sendo perpendiculares, compõem-se por Pitágoras:
\[
E=\sqrt{3^{2}+4^{2}}=\SI{5}{\kilo\newton\per\coulomb};
\qquad
\theta=\arctan\frac{4}{3}=\ang{53.1}.
\]
\textbf{Contraste com o potencial:} se esses mesmos pontos tivessem potenciais
de $3$ e $\SI{4}{\kilo\volt}$, o total seria simplesmente
$\SI{7}{\kilo\volt}$ --- soma algébrica, sem geometria.\\
É essa diferença que torna o potencial a ferramenta preferida para cálculos, e o
campo a ferramenta preferida para entender forças e movimento.}
\end{questao}

\begin{questao}[2]{}
Em uma representação por linhas de campo, o que indica uma região onde as
linhas estão mais próximas entre si?
\begin{alternativas}[2]
\item potencial maior
\item campo mais intenso
\item carga maior naquele ponto
\item ausência de campo
\end{alternativas}
\resposta{(b)}
\resolucao{A \textbf{densidade de linhas} é proporcional ao módulo do campo ---
essa é a convenção que dá utilidade quantitativa ao desenho.\\
\textbf{Três regras das linhas de campo:}
\begin{enumerate}
\item Nascem em cargas positivas e morrem em negativas (ou no infinito).
\item São tangentes a $\vec E$ em cada ponto.
\item \textbf{Nunca se cruzam} --- em um cruzamento, o campo teria duas
direções, o que é impossível.
\end{enumerate}
\textbf{Cuidado com (a):} potencial e campo são grandezas distintas. Linhas
próximas indicam campo intenso, o que corresponde a equipotenciais próximas ---
mas não diz nada sobre o \emph{valor} de $V$.}
\end{questao}

\begin{questao}[2]{}
Uma carga produz campo de $\SI{8}{\kilo\newton\per\coulomb}$ no vácuo. O
conjunto é imerso em um dielétrico de $\varepsilon_r=4$. Determine o novo campo.
\resposta{$E=\SI{2}{\kilo\newton\per\coulomb}$}
\resolucao{
\[
E_{meio}=\frac{E_{vácuo}}{\varepsilon_r}=\frac{8}{4}
=\SI{2}{\kilo\newton\per\coulomb}.
\]
\textbf{Mecanismo:} o dielétrico se \emph{polariza} --- suas moléculas se
orientam e criam um campo interno oposto, que cancela parcialmente o campo
original.\\
\textbf{Consequência tecnológica:} o dielétrico reduz o campo para a mesma
carga, o que permite armazenar \emph{mais} carga para a mesma tensão --- é
exatamente por isso que $C=\varepsilon_r C_0$ em capacitores.}
\end{questao}

\begin{questao}[2]{}
Uma gota de massa $\SI{2}{\milli\gram}$ e carga $\SI{1}{\nano\coulomb}$ deve
ficar suspensa em repouso. Determine o campo elétrico necessário.
\dados{$g=\SI{9.8}{\meter\per\second\squared}$}
\resposta{$E=\SI{19.6}{\kilo\newton\per\coulomb}$}
\resolucao{Equilíbrio entre peso e força elétrica:
\[
qE=mg
\;\Longrightarrow\;
E=\frac{mg}{q}=\frac{(\pot{2}{-6})(9{,}8)}{\pot{1}{-9}}
=\pot{1.96}{4}\,\si{\newton\per\coulomb}.
\]
\textbf{Orientação:} o campo deve produzir força \emph{para cima}. Se a gota for
positiva, $\vec E$ aponta para cima; se negativa, para baixo.\\
\textbf{Este é o princípio do experimento de Millikan}, que determinou a carga
elementar. Medindo $E$ e $m$ de gotas suspensas, ele encontrou que todas as
cargas eram múltiplas inteiras de $\pot{1.6}{-19}\,\si{\coulomb}$ --- a
primeira evidência direta da quantização.}
\end{questao}

\blocoaprofundamento

\begin{questao}[3]{}
Duas cargas iguais $Q=+\SI{1}{\micro\coulomb}$ ocupam
$(\pm\SI{0.20}{\meter},\,0)$.
\begin{itensq}
\item Determine o campo em $(0,\ \SI{0.20}{\meter})$.
\item Determine o campo no ponto médio e justifique.
\end{itensq}
\resposta{(a) $E=\SI{159}{\kilo\newton\per\coulomb}$ ao longo de $+y$;
(b) $\vec E=\vec0$}
\resolucao{(a) Cada carga dista
$r=\sqrt{2}\,(0{,}20)=\SI{0.283}{\meter}$:
\[
E_1=E_2=\frac{\pot{9}{9}(\pot{1}{-6})}{0{,}08}
=\pot{1.125}{5}\,\si{\newton\per\coulomb}.
\]
As componentes em $x$ cancelam; as em $y$ somam-se, cada uma multiplicada por
$\cos\ang{45}=0{,}707$:
\[
E=2(\pot{1.125}{5})(0{,}707)=\pot{1.59}{5}\,\si{\newton\per\coulomb}.
\]
(b) \textbf{No ponto médio}, as duas contribuições têm módulos iguais e sentidos
\emph{opostos} (cada carga empurra para longe de si):
\[
\vec E=\vec 0.
\]
\textbf{Mas o potencial ali não é nulo:}
$V=2\ke Q/0{,}20=\SI{90}{\kilo\volt}$. É o caso complementar ao do dipolo, em
que $V=0$ e $\vec E\ne\vec0$ --- e mostra novamente que as duas condições são
independentes.\\
\textbf{Estabilidade:} uma carga positiva no ponto médio está em equilíbrio, mas
\emph{instável} na direção perpendicular ao eixo --- conforme o teorema de
Earnshaw.}
\end{questao}

\begin{questao}[3]{}
Um dipolo é formado por $+q$ e $-q$ separadas por $2a$. Considere um ponto do
eixo perpendicular à linha que as une, a distância $y$ do centro.
\begin{itensq}
\item Deduza a expressão exata do campo.
\item Obtenha a forma assintótica para $y\gg a$.
\item Compare a queda com a de uma carga isolada.
\end{itensq}
\resposta{$E=\dfrac{2\ke qa}{(y^{2}+a^{2})^{3/2}}\to\dfrac{2\ke qa}{y^{3}}$}
\resolucao{(a) Cada carga dista $r=\sqrt{y^{2}+a^{2}}$ e produz
$\ke q/r^{2}$. Por simetria, as componentes ao longo de $y$ se cancelam (uma
aponta para cima, a outra para baixo) e restam as componentes \emph{paralelas ao
eixo do dipolo}, cada uma multiplicada por $a/r$:
\[
E=2\cdot\frac{\ke q}{y^{2}+a^{2}}\cdot\frac{a}{\sqrt{y^{2}+a^{2}}}
=\frac{2\ke qa}{(y^{2}+a^{2})^{3/2}}.
\]
(b) Para $y\gg a$, $(y^{2}+a^{2})^{3/2}\to y^{3}$:
\[
E\approx\frac{2\ke qa}{y^{3}}=\frac{\ke p}{y^{3}},
\]
com $p=2qa$ o momento dipolar.\\
(c) \textbf{O campo do dipolo cai com $1/y^{3}$}, muito mais rápido que o
$1/r^{2}$ de uma carga isolada.\\
\textbf{Razão física:} de longe, as duas cargas quase se cancelam --- o sistema
é \emph{globalmente neutro}. O que sobra é o efeito de sua \emph{separação},
uma correção de ordem superior.\\
\textbf{Padrão geral:} carga isolada $\sim1/r^{2}$; dipolo
$\sim1/r^{3}$; quadrupolo $\sim1/r^{4}$. Quanto mais neutra a distribuição, mais
rápido o campo se extingue --- e é por isso que a matéria comum, neutra, não
exerce força elétrica a distância.}
\end{questao}

\begin{questao}[3]{}
Duas placas condutoras extensas e paralelas têm densidades superficiais
$+\sigma$ e $-\sigma$, com
$\sigma=\SI{1}{\micro\coulomb\per\meter\squared}$.
\begin{itensq}
\item Determine o campo entre as placas.
\item Determine o campo fora delas.
\item Justifique fisicamente.
\end{itensq}
\resposta{Entre: $E=\SI{113}{\kilo\newton\per\coulomb}$; fora: $E=0$}
\resolucao{(a) Cada plano contribui com $\sigma/2\varepsilon_0$, e
\textbf{entre} as placas as duas contribuições têm o \emph{mesmo} sentido (do
positivo para o negativo):
\[
E=2\cdot\frac{\sigma}{2\varepsilon_0}=\frac{\sigma}{\varepsilon_0}
=\frac{\pot{1}{-6}}{\pot{8.85}{-12}}=\pot{1.13}{5}\,\si{\newton\per\coulomb}.
\]
(b) \textbf{Fora}, as duas contribuições têm sentidos \emph{opostos} e módulos
iguais:
\[
E=\frac{\sigma}{2\varepsilon_0}-\frac{\sigma}{2\varepsilon_0}=0.
\]
(c) \textbf{Justificativa:} o campo de um plano infinito é uniforme e aponta
para longe dele (se positivo) ou para ele (se negativo). Na região externa,
qualquer ponto está do mesmo lado das duas placas --- e os campos se opõem.
Entre elas, o ponto está "depois" da positiva e "antes" da negativa, e os campos
se somam.\\
\textbf{Consequência: blindagem natural.} O par de placas confina completamente
seu campo na região interna --- não há campo escapando para fora. É por isso que
um capacitor de placas paralelas não interfere no circuito vizinho, e por que a
energia armazenada está inteiramente entre as armaduras.}
\end{questao}

\begin{questao}[3]{}
Uma esfera condutora de raio $R=\SI{10}{\centi\meter}$ tem carga
$Q=\SI{5}{\micro\coulomb}$.
\begin{itensq}
\item Determine o campo no interior, na superfície e a $\SI{30}{\centi\meter}$
      do centro.
\item Avalie se a configuração é fisicamente sustentável no ar.
\end{itensq}
\resposta{$0$; $\SI{4.5}{\mega\volt\per\meter}$;
$\SI{500}{\kilo\volt\per\meter}$ --- \textbf{insustentável}}
\resolucao{(a) \textbf{Interior:} $E=0$, condição de equilíbrio de qualquer
condutor.\\
\textbf{Superfície:} toda a carga se comporta como concentrada no centro, para
pontos externos:
\[
E=\frac{\ke Q}{R^{2}}=\frac{\pot{9}{9}(\pot{5}{-6})}{0{,}01}
=\pot{4.5}{6}\,\si{\volt\per\meter}.
\]
\textbf{A $\SI{30}{\centi\meter}$:}
$E=\dfrac{\pot{4.5}{4}}{0{,}09}=\pot{5}{5}\,\si{\volt\per\meter}$.\\
(b) \textbf{Insustentável.} O campo na superfície,
$\SI{4.5}{\mega\volt\per\meter}$, excede a rigidez dielétrica do ar
($\SI{3}{\mega\volt\per\meter}$) em $50\%$. O ar ioniza e a carga escapa por
descarga corona.\\
\textbf{Carga máxima real:}
$Q_{\max}=E_{rup}R^{2}/\ke=\SI{3.33}{\micro\coulomb}$. A esfera simplesmente
\emph{não consegue} reter os $\SI{5}{\micro\coulomb}$ enunciados.\\
\textbf{Lição de método:} sempre verifique se o resultado é fisicamente
possível. Um problema pode ser aritmeticamente consistente e descrever uma
situação inexistente.}
\end{questao}

\begin{questao}[3]{}
Um elétron parte do repouso em um campo uniforme de
$E=\SI{1000}{\volt\per\meter}$ e percorre $\SI{2}{\centi\meter}$.
\begin{itensq}
\item Determine a aceleração e o tempo de percurso.
\item Determine a velocidade final e a energia adquirida.
\end{itensq}
\dados{$e=\pot{1.6}{-19}\,\si{\coulomb}$; $m_e=\pot{9.11}{-31}\,\si{\kilogram}$}
\resposta{$a=\pot{1.76}{14}\,\si{\meter\per\second\squared}$;
$t=\SI{15.1}{\nano\second}$; $v=\pot{2.65}{6}\,\si{\meter\per\second}$;
$\SI{20}{\electronvolt}$}
\resolucao{(a)
\[
a=\frac{eE}{m_e}=\frac{(\pot{1.6}{-19})(1000)}{\pot{9.11}{-31}}
=\pot{1.756}{14}\,\si{\meter\per\second\squared};
\]
\[
t=\sqrt{\frac{2d}{a}}=\sqrt{\frac{2(0{,}02)}{\pot{1.756}{14}}}
=\pot{1.51}{-8}\,\si{\second}=\SI{15.1}{\nano\second}.
\]
(b) $v=at=(\pot{1.756}{14})(\pot{1.51}{-8})
=\pot{2.65}{6}\,\si{\meter\per\second}$.\\
\textbf{Energia:} o mais rápido é usar a tensão percorrida,
$U=Ed=1000(0{,}02)=\SI{20}{\volt}$:
\[
W=eU=\SI{20}{\electronvolt}=\pot{3.2}{-18}\,\si{\joule}.
\]
\emph{(Conferindo: $\frac12m_ev^{2}=\frac12(\pot{9.11}{-31})(\pot{2.65}{6})^{2}
=\pot{3.2}{-18}\,\si{\joule}$ \checkmark.)}\\
\textbf{Verificação relativística:} $v/c=0{,}9\%$ --- plenamente no regime
clássico.}
\end{questao}

\begin{questao}[3]{}
Compare a queda do campo com a distância para três configurações: carga
puntiforme, fio infinito carregado e plano infinito carregado.
\begin{itensq}
\item Escreva a dependência de cada uma.
\item Explique a origem geométrica das diferenças.
\end{itensq}
\resposta{$1/r^{2}$, $1/r$ e constante}
\resolucao{(a)
\begin{center}
\begin{tabular}{lcc}
\hline
Distribuição & Campo & Queda\\
\hline
Carga puntiforme & $\dfrac{\ke Q}{r^{2}}$ & $1/r^{2}$\\[6pt]
Fio infinito ($\lambda$) & $\dfrac{\lambda}{2\pi\varepsilon_0 r}$ & $1/r$\\[6pt]
Plano infinito ($\sigma$) & $\dfrac{\sigma}{2\varepsilon_0}$ & constante\\
\hline
\end{tabular}
\end{center}
(b) \textbf{Origem geométrica.} Pela lei de Gauss, o fluxo total é fixo pela
carga envolvida, e o campo é o fluxo dividido pela área da superfície:
\begin{itemize}
\item \textbf{Ponto:} a superfície é uma \emph{esfera}, de área
$4\pi r^{2}$ --- o campo cai com $1/r^{2}$.
\item \textbf{Fio:} a superfície é um \emph{cilindro}, de área lateral
$2\pi r\ell$ --- cresce apenas com $r$, e o campo cai com $1/r$.
\item \textbf{Plano:} a superfície é um \emph{prisma} cuja área \emph{não
depende} de $r$ --- e o campo não cai.
\end{itemize}
\textbf{Regra unificadora:} o campo cai como $1/r^{n}$, onde $n$ é o número de
dimensões em que a fonte é \emph{limitada}. Ponto: limitado em três
($n=2$). Fio: em duas ($n=1$). Plano: em uma ($n=0$).\\
\textbf{Consequência prática:} perto de uma placa carregada extensa, afastar-se
não ajuda --- o campo permanece. É o oposto da intuição formada com cargas
puntiformes.}
\end{questao}

\begin{questao}[3]{}
Uma haste isolante fina é carregada de modo que produz, em certo ponto, campo
de $\SI{600}{\newton\per\coulomb}$. A haste é então partida ao meio, e apenas
uma das metades permanece.
\begin{itensq}
\item O campo cai à metade? Justifique.
\item Determine em que condição isso valeria.
\end{itensq}
\resposta{Não necessariamente --- só vale se a metade removida contribuísse com
exatamente metade do campo}
\resolucao{(a) \textbf{Não, em geral.} O campo é a soma \emph{vetorial} das
contribuições de todos os elementos da haste, e cada elemento contribui com
módulo e direção diferentes, dependendo de sua distância e posição angular em
relação ao ponto.\\
Se o ponto estiver \emph{sobre a mediatriz} da haste, as duas metades
contribuem igualmente e o campo de fato cai à metade --- mas a \emph{direção}
resultante também muda, pois as componentes que antes se cancelavam agora não
se cancelam mais.\\
Se o ponto estiver \emph{sobre o prolongamento} da haste, a metade mais próxima
contribui muito mais que a distante, e remover a distante reduz o campo bem
menos que $50\%$.\\
(b) \textbf{A queda é exatamente à metade quando:}
\begin{itemize}
\item as duas metades contribuem com vetores \emph{idênticos} --- o que exige
simetria completa em relação ao ponto; \emph{e}
\item essa condição só se realiza em pontos muito distantes, onde a haste inteira
se comporta como carga puntiforme e a divisão é simplesmente $Q\to Q/2$.
\end{itemize}
\textbf{Lição:} a superposição vale \emph{sempre}, mas ela é vetorial. Reduzir
a fonte pela metade não reduz o campo pela metade, exceto em casos de simetria
particular.}
\end{questao}

\blocodesafio

\begin{questao}[4]{}
\textbf{(Campo nulo com potencial não nulo.)} Construa uma configuração em que
$\vec E=\vec0$ em um ponto onde $V\ne0$.
\begin{itensq}
\item Proponha o arranjo mais simples e determine os valores.
\item Explique por que isso é possível.
\item Determine se o inverso ($V=0$ com $\vec E\ne\vec0$) também ocorre.
\end{itensq}
\resposta{Duas cargas iguais $+Q$; no ponto médio $\vec E=\vec0$ e
$V=2\ke Q/a$}
\resolucao{(a) \textbf{Arranjo:} duas cargas iguais $Q=+\SI{2}{\micro\coulomb}$
em $(\pm\SI{0.20}{\meter},\,0)$. No ponto médio:
\[
\vec E=\vec 0
\quad\text{(contribuições opostas)};
\]
\[
V=2\cdot\frac{\ke Q}{0{,}20}
=2\cdot\frac{\pot{9}{9}(\pot{2}{-6})}{0{,}20}
=\pot{1.8}{5}\,\si{\volt}=\SI{180}{\kilo\volt}.
\]
Campo nulo, potencial de cento e oitenta mil volts.\\
(b) \textbf{Por que é possível.} As duas grandezas se relacionam por
$\vec E=-\nabla V$: o campo é a \emph{taxa de variação} de $V$, não seu valor.
Uma função pode ter valor grande e derivada nula --- é exatamente o que ocorre
em um ponto \emph{estacionário}.\\
No ponto médio, $V$ tem um mínimo ao longo do eixo que une as cargas: o valor é
alto, mas a inclinação é zero.\\
(c) \textbf{Sim, o inverso também ocorre.} Substituindo uma das cargas por
$-Q$, o ponto médio passa a ter $V=0$ (contribuições opostas se cancelam
algebricamente) e $\vec E\ne\vec0$ (os vetores apontam para o mesmo lado e se
somam).\\
\textbf{Síntese:} $V=0$ e $\vec E=\vec0$ são condições \textbf{logicamente
independentes}. Nenhuma implica a outra, em nenhum dos dois sentidos.}
\end{questao}

\begin{questao}[4]{}
\textbf{(Linhas de campo.)} Analise as linhas de campo de um dipolo elétrico.
\begin{itensq}
\item Descreva seu traçado.
\item Demonstre que linhas de campo nunca se cruzam.
\item Explique o que ocorre em pontos de campo nulo.
\end{itensq}
\resposta{Saem de $+q$ e entram em $-q$; não se cruzam porque $\vec E$ tem
direção única em cada ponto}
\resolucao{(a) \textbf{Traçado.} As linhas saem radialmente de $+q$, curvam-se e
entram em $-q$. A linha que segue o eixo do dipolo é reta; as demais formam
arcos progressivamente mais abertos. Algumas partem para o infinito e retornam.\\
A densidade de linhas é máxima entre as cargas --- onde o campo é mais intenso
--- e decresce com $1/r^{3}$ ao se afastar, conforme demonstrado antes.\\
(b) \textbf{Demonstração.} Suponha, por absurdo, que duas linhas se cruzem em
um ponto $P$. Por definição, a linha de campo é \emph{tangente} a $\vec E$ em
cada ponto. Se duas linhas distintas passam por $P$, elas têm tangentes
distintas ali --- e portanto $\vec E(P)$ teria \textbf{duas direções}
simultaneamente.\\
Mas $\vec E$ é uma função vetorial \emph{bem definida}: a cada ponto do espaço
corresponde \emph{um} vetor. Contradição. Logo, as linhas não se cruzam.\\
\textbf{Corolário:} o mesmo argumento vale para linhas de corrente, linhas de
campo magnético e qualquer outro campo vetorial.\\
(c) \textbf{Pontos de campo nulo são a exceção.} Onde $\vec E=\vec0$, não há
direção definida --- e ali as linhas \emph{podem} se encontrar, sem contradição.
São os \textbf{pontos de sela} do potencial.\\
Em um dipolo ($+q$ e $-q$) não existe tal ponto no espaço finito. Mas em duas
cargas \emph{iguais}, o ponto médio é exatamente um ponto de campo nulo --- e
ali as linhas se aproximam de todos os lados sem se tocarem, definindo uma
separatriz.}
\end{questao}

\begin{questao}[4]{}
\textbf{(Projeto de acelerador de placas.)} Projete um arranjo de placas
paralelas que imprima aceleração de
$a=\pot{2}{12}\,\si{\meter\per\second\squared}$ a um elétron.
\begin{itensq}
\item Determine o campo necessário.
\item Determine a tensão para separação de $\SI{1}{\centi\meter}$.
\item Determine o tempo de trânsito e a velocidade de saída.
\end{itensq}
\dados{$e=\pot{1.6}{-19}\,\si{\coulomb}$; $m_e=\pot{9.11}{-31}\,\si{\kilogram}$}
\resposta{$E=\SI{11.4}{\volt\per\meter}$; $U=\SI{0.114}{\volt}$;
$t=\SI{100}{\nano\second}$; $v=\pot{2}{5}\,\si{\meter\per\second}$}
\resolucao{(a)
\[
E=\frac{m_ea}{e}=\frac{(\pot{9.11}{-31})(\pot{2}{12})}{\pot{1.6}{-19}}
=\SI{11.4}{\volt\per\meter}.
\]
(b) $U=Ed=11{,}4(0{,}01)=\SI{0.114}{\volt}$.\\
(c)
\[
t=\sqrt{\frac{2d}{a}}=\sqrt{\frac{0{,}02}{\pot{2}{12}}}
=\pot{1}{-7}\,\si{\second}=\SI{100}{\nano\second};
\]
\[
v=at=(\pot{2}{12})(\pot{1}{-7})=\pot{2}{5}\,\si{\meter\per\second}.
\]
\textbf{O resultado é surpreendente e instrutivo.} Uma aceleração de
$\pot{2}{12}\,\si{\meter\per\second\squared}$ --- duzentos bilhões de vezes
$g$ --- exige apenas $\SI{0.114}{\volt}$, tensão menor que a de uma pilha
descarregada.\\
\textbf{Razão:} a razão carga/massa do elétron é enorme
($\pot{1.76}{11}\,\si{\coulomb\per\kilogram}$). Campos elétricos modestos
produzem acelerações colossais em partículas leves.\\
\textbf{Ressalva prática:} com tensão tão baixa, o projeto seria dominado por
efeitos parasitas --- diferenças de função-trabalho entre os metais das placas
(décimos de volt), cargas superficiais e ruído térmico. Um projeto real usaria
tensões de dezenas de volts e aceleração proporcionalmente maior.}
\end{questao}

\begin{questao}[4]{}
\textbf{(Cavidade em condutor.)} Uma carga puntiforme $+q$ é colocada no
interior de uma cavidade vazia dentro de um condutor \emph{neutro} e isolado.
\begin{itensq}
\item Determine a carga induzida em cada superfície.
\item Determine o campo dentro do metal e fora do condutor.
\item Analise o que muda se o condutor for aterrado.
\end{itensq}
\resposta{$-q$ na parede da cavidade, $+q$ na superfície externa; campo nulo no
metal e não nulo fora}
\resolucao{(a) \textbf{Parede da cavidade:} aplicando a lei de Gauss a uma
superfície inteiramente dentro do metal (onde $\vec E=\vec0$), o fluxo é nulo e
a carga envolvida também. Como a superfície envolve $+q$ e a parede interna:
\[
q+q_{interna}=0
\;\Longrightarrow\;
q_{interna}=-q.
\]
\textbf{Superfície externa:} o condutor era neutro, então a carga total deve
permanecer nula:
\[
q_{externa}=+q.
\]
(b) \textbf{Dentro do metal:} $\vec E=\vec0$, por equilíbrio.\\
\textbf{Fora:} a carga $+q$ da superfície externa produz campo --- e ele é
\textbf{esfericamente simétrico} se a superfície externa for esférica,
\emph{independentemente} de onde a carga esteja dentro da cavidade.\\
\textbf{Consequência notável:} de fora, é impossível saber a posição da carga
interna, nem a forma da cavidade. O condutor "apaga" essa informação.\\
(c) \textbf{Se o condutor for aterrado:} a carga $+q$ da superfície externa
escoa para a terra. Resta apenas $-q$ na parede da cavidade, e
\[
\vec E=\vec 0\quad\text{em todo o exterior}.
\]
\textbf{A blindagem torna-se completa nos dois sentidos.} É exatamente esse o
princípio da gaiola de Faraday aterrada --- e a diferença entre blindar contra
campos externos (que dispensa aterramento) e impedir que fontes internas
irradiem (que o exige).}
\end{questao}

\begin{questao}[4]{}
\textbf{(Esfera isolante uniformemente carregada.)} Uma esfera isolante maciça
de raio $R$ tem carga $Q$ distribuída uniformemente em seu \emph{volume}.
\begin{itensq}
\item Determine o campo para $r<R$ e para $r>R$.
\item Identifique onde o campo é máximo.
\item Compare com a esfera \emph{condutora} de mesma carga.
\end{itensq}
\resposta{$E=\dfrac{\ke Qr}{R^{3}}$ dentro e $\dfrac{\ke Q}{r^{2}}$ fora;
máximo na superfície}
\resolucao{(a) \textbf{Fora ($r>R$):} pela lei de Gauss com superfície esférica,
toda a carga está envolvida:
\[
E=\frac{\ke Q}{r^{2}},
\]
idêntico ao de uma carga puntiforme no centro.\\
\textbf{Dentro ($r<R$):} apenas a fração de carga interna à superfície gaussiana
conta. Como a densidade é uniforme, essa fração é $(r/R)^{3}$:
\[
E=\frac{\ke\,Q(r/R)^{3}}{r^{2}}=\frac{\ke Q\,r}{R^{3}},
\]
\textbf{proporcional a $r$} --- cresce linearmente do centro à superfície.\\
(b) \textbf{Máximo na superfície} ($r=R$), onde as duas expressões coincidem em
$E=\ke Q/R^{2}$. O campo é contínuo ali, embora com derivadas diferentes de
cada lado.\\
\emph{No centro, $E=0$ por simetria.}\\
(c) \textbf{Comparação com a esfera condutora.}
\begin{center}
\begin{tabular}{lcc}
\hline
Região & Isolante & Condutora\\
\hline
$r<R$ & $\ke Qr/R^{3}$ & $0$\\
$r=R$ & $\ke Q/R^{2}$ & $\ke Q/R^{2}$\\
$r>R$ & $\ke Q/r^{2}$ & $\ke Q/r^{2}$\\
\hline
\end{tabular}
\end{center}
\textbf{São indistinguíveis do lado de fora} --- a lei de Gauss só "enxerga" a
carga total envolvida, não sua distribuição. A diferença está inteiramente no
interior: no condutor a carga migra para a superfície e o campo interno é nulo;
no isolante ela permanece onde foi colocada.}
\end{questao}

\begin{questao}[4]{}
\textbf{(Traçado numérico de linhas de campo.)} Descreva um procedimento
computacional para traçar linhas de campo a partir de uma expressão conhecida de
$\vec E(x,y)$.
\begin{itensq}
\item Formule o algoritmo.
\item Indique como escolher os pontos de partida.
\item Aponte as dificuldades numéricas e como contorná-las.
\end{itensq}
\resposta{Integrar $\dfrac{\mathrm{d}\vec r}{\mathrm{d}s}
=\dfrac{\vec E}{|\vec E|}$ a partir de pontos distribuídos em torno das cargas}
\resolucao{(a) \textbf{Algoritmo.} A linha de campo é, por definição, tangente a
$\vec E$. Parametrizando pelo comprimento de arco $s$:
\[
\frac{\mathrm{d}\vec r}{\mathrm{d}s}=\frac{\vec E(\vec r)}{|\vec E(\vec r)|}.
\]
Numericamente, com passo $h$:
\begin{enumerate}
\item Escolha um ponto inicial $\vec r_0$.
\item Calcule $\vec E(\vec r_n)$ e normalize.
\item Avance: $\vec r_{n+1}=\vec r_n+h\,\hat E$ (Euler) ou, melhor, por
Runge--Kutta de quarta ordem.
\item Repita até sair do domínio ou atingir uma carga negativa.
\end{enumerate}
\textbf{Normalizar é essencial:} sem isso, o passo seria proporcional a
$|\vec E|$ e a linha "voaria" perto das cargas e travaria longe delas.\\
(b) \textbf{Pontos de partida.} Distribua-os em pequenos círculos ao redor de
cada carga \emph{positiva}, com espaçamento angular uniforme. Para que a
densidade de linhas represente corretamente a intensidade, o número de linhas
por carga deve ser \emph{proporcional a $|q|$} --- oito linhas para $+q$,
dezesseis para $+2q$.\\
(c) \textbf{Dificuldades e soluções.}
\begin{itemize}
\item \textbf{Singularidades:} $|\vec E|\to\infty$ nas cargas. Solução: iniciar
a uma distância pequena mas finita, e interromper a integração ao entrar em um
raio mínimo em torno de uma carga negativa.
\item \textbf{Pontos de campo nulo:} ali $\hat E$ é indefinido e o algoritmo
trava ou oscila. Solução: detectar $|\vec E|<\epsilon$ e encerrar a linha.
\item \textbf{Acúmulo de erro:} o método de Euler desvia sistematicamente em
regiões de curvatura alta. Solução: Runge--Kutta e passo adaptativo, menor onde
o campo varia rápido.
\item \textbf{Linhas que escapam:} defina um domínio e encerre ao ultrapassá-lo.
\end{itemize}
\textbf{Verificação do resultado:} as linhas traçadas devem ser perpendiculares
às equipotenciais calculadas independentemente --- teste que detecta quase todo
erro de implementação.}
\end{questao}

\begin{questao}[4]{}
\textbf{(Limite de campo em capacitor.)} Um capacitor de placas paralelas com
$d=\SI{0.1}{\milli\meter}$ usa dielétrico de rigidez
$\SI{20}{\kilo\volt\per\milli\meter}$ e $\varepsilon_r=3$.
\begin{itensq}
\item Determine a tensão máxima e o campo correspondente.
\item Determine a densidade superficial de carga no limite.
\item Discuta o compromisso entre capacitância e tensão.
\end{itensq}
\dados{$\varepsilon_0=\pot{8.85}{-12}\,\si{\farad\per\meter}$}
\resposta{$U_{\max}=\SI{2}{\kilo\volt}$;
$E=\SI{2e7}{\volt\per\meter}$;
$\sigma=\SI{531}{\micro\coulomb\per\meter\squared}$}
\resolucao{(a)
\[
U_{\max}=E_{rup}\,d=20\times0{,}1=\SI{2}{\kilo\volt};
\qquad
E=\pot{2}{7}\,\si{\volt\per\meter}.
\]
(b) No dielétrico, $E=\dfrac{\sigma}{\varepsilon_r\varepsilon_0}$, logo
\[
\sigma=\varepsilon_r\varepsilon_0E
=3(\pot{8.85}{-12})(\pot{2}{7})
=\pot{5.31}{-4}\,\si{\coulomb\per\meter\squared}.
\]
(c) \textbf{O compromisso.} A capacitância por unidade de área é
$C/A=\varepsilon_r\varepsilon_0/d$: reduzir $d$ \emph{aumenta} $C$. Mas
$U_{\max}=E_{rup}d$: reduzir $d$ \emph{diminui} a tensão suportável.\\
O produto revela o que realmente importa:
\[
\frac{C\,U_{\max}}{A}=\varepsilon_r\varepsilon_0E_{rup},
\]
que \textbf{não depende de $d$}. E a energia por unidade de volume,
\[
u=\frac12\varepsilon_r\varepsilon_0E_{rup}^{2},
\]
também não. \textbf{A geometria não decide nada} --- só o material.\\
\textbf{Consequência:} não existe truque de projeto que aumente a energia
armazenada por volume. Todo avanço em capacitores vem de materiais com
$\varepsilon_r$ ou $E_{rup}$ maiores --- e esses dois parâmetros costumam ser
antagônicos, o que torna o problema difícil.}
\end{questao}

\begin{questao}[4]{}
\textbf{(Superposição vetorial em anel.)} Um anel de raio $R$ tem carga
$Q$ uniformemente distribuída.
\begin{itensq}
\item Explique por que o campo é nulo no centro.
\item Determine, sem integrar, o campo em um ponto muito distante no eixo.
\item Discuta o que ocorre se metade do anel for removida.
\end{itensq}
\resposta{Centro: $\vec E=\vec0$ por simetria; longe: $E\to\ke Q/x^{2}$;
com meio anel, campo não nulo no centro}
\resolucao{(a) \textbf{No centro,} para cada elemento de carga existe outro
diametralmente oposto, à mesma distância $R$. Suas contribuições têm módulos
iguais e sentidos opostos, cancelando-se aos pares:
\[
\vec E_{centro}=\vec 0.
\]
\textbf{Note:} o \emph{potencial} no centro não é nulo --- vale
$V=\ke Q/R$, pois potenciais somam-se sem cancelamento entre contribuições de
mesmo sinal. Mais um caso de $\vec E=\vec0$ com $V\ne0$.\\
(b) \textbf{Muito longe} ($x\gg R$), todos os elementos estão praticamente à
mesma distância $x$ e praticamente na mesma direção. O anel se comporta como uma
carga puntiforme:
\[
E\to\frac{\ke Q}{x^{2}}.
\]
\textbf{Este é o teste que toda expressão deve passar:} qualquer distribuição
finita, vista de longe, deve reproduzir a lei de Coulomb.\\
(c) \textbf{Com meio anel}, a simetria de cancelamento se perde. Cada elemento
não tem mais seu par oposto, e o campo no centro \textbf{deixa de ser nulo}.\\
Pelo argumento da superposição ao contrário: como o anel inteiro dá campo nulo,
o campo do semianel restante é o \emph{negativo} do campo do semianel removido
--- ou seja, os dois têm o mesmo módulo. Integrando (como no caso do fio
semicircular), obtém-se
\[
E=\frac{2\ke Q'}{\pi R^{2}},
\]
com $Q'=Q/2$ a carga de cada metade, dirigido ao longo do eixo de simetria.}
\end{questao}

% END BANCO COMPLEMENTAR

\gabarito[1.3 Campo elétrico]

% =====================================================================
%  CAPITULO 1 - continuacao
%  1.4 Potencial eletrico e superficies equipotenciais
% =====================================================================
\subsection{Potencial elétrico e superfícies equipotenciais}

\begin{conceito}[Antes de começar]
\textbf{Potencial de uma carga puntiforme:}
\[
V = \ke\,\frac{q}{d}\quad(\si{\volt}),
\]
com o \emph{sinal} da carga (positivo para $q>0$, negativo para $q<0$).\par
\textbf{Diferença de potencial (ddp)} e trabalho: o trabalho para levar uma carga
$q$ entre dois pontos é $W = q\,U_{AB} = q\,(V_A - V_B)$.\par
\textbf{Campo uniforme:} $U = E\,d$, onde $d$ é a distância medida \emph{ao longo}
da direção do campo.\par
\textbf{Superfície equipotencial:} conjunto de pontos com o mesmo potencial. O
campo elétrico é sempre \textbf{perpendicular} às equipotenciais e aponta no
sentido dos potenciais \emph{decrescentes}. Não há trabalho ao mover uma carga
sobre uma mesma equipotencial.
\end{conceito}

\legendaniveis

\blocofixacao

\begin{questao}[1]{}
Sobre superfícies equipotenciais, é correto afirmar que:
\begin{alternativas}
\item o campo elétrico é tangente a elas.
\item o campo elétrico é perpendicular a elas.
\item o potencial varia ao longo de uma mesma superfície.
\item é necessário realizar trabalho para mover uma carga sobre a superfície.
\item elas sempre coincidem com as linhas de campo.
\end{alternativas}
\resposta{(b)}
\resolucao{Por definição, o potencial é constante sobre uma equipotencial (isso
elimina c). Como não há variação de potencial ao longo dela, o trabalho para mover
uma carga sobre a superfície é nulo (elimina d). O campo, que aponta na direção de
maior variação do potencial, é sempre \emph{perpendicular} às equipotenciais ---
nunca tangente, o que elimina (a) e (e).}
\end{questao}

\begin{questao}[1]{}
Considere dois pontos A e B de um campo elétrico uniforme de intensidade
$E = \SI{e4}{\newton\per\coulomb}$, separados por $\SI{5}{\centi\meter}$ ao longo
da direção do campo. Calcule a ddp entre A e B.
\resposta{$U_{AB} = \SI{500}{\volt}$}
\resolucao{Em campo uniforme, $U = E\,d = \pot{1}{4}\cdot0{,}05 = \SI{500}{\volt}$
(com A no potencial maior, a favor do campo).}
\end{questao}

\blocoaplicacao

\begin{questao}[2]{}
A figura representa três superfícies equipotenciais de um campo elétrico
uniforme, com $V_A = \SI{60}{\volt}$, $V_B = \SI{40}{\volt}$ e
$V_C = \SI{20}{\volt}$, igualmente espaçadas de $\SI{2}{\centi\meter}$.

\circuitoq{%
\begin{tikzpicture}[scale=1,line width=0.8pt,>=Stealth]
 \foreach \x/\v in {0/A,2/B,4/C}{
   \draw[azulprof,thick] (\x,-1.3) -- (\x,1.3);
   \node[above,azulprof] at (\x,1.3) {$V_{\v}$};}
 \node[below] at (0,-1.3) {\footnotesize$\SI{60}{\volt}$};
 \node[below] at (2,-1.3) {\footnotesize$\SI{40}{\volt}$};
 \node[below] at (4,-1.3) {\footnotesize$\SI{20}{\volt}$};
 % campo (do maior para o menor potencial)
 \foreach \y in {-0.7,0,0.7}
   \draw[->,Vcol,thick] (0.3,\y) -- (3.7,\y);
 \node[Vcol] at (2,1.0) {$\vec{E}$};
 \draw[<->,cinza] (0,-1.9) -- (2,-1.9) node[midway,below]{\footnotesize$\SI{2}{\centi\meter}$};
\end{tikzpicture}}{Superfícies equipotenciais em campo uniforme.}

\begin{itensq}
\item Determine a intensidade do campo elétrico.
\item Qual é o trabalho para levar uma carga de $\SI{5}{\micro\coulomb}$ de A até C?
\end{itensq}
\resposta{(a) $E = \SI{1000}{\volt\per\meter}$; (b) $W = \SI{2e-4}{\joule}$}
\resolucao{(a) Entre superfícies adjacentes, $U = \SI{20}{\volt}$ em
$d = \SI{2}{\centi\meter}$:
\[
E=\frac{U}{d}=\frac{20}{0{,}02}=\SI{1000}{\volt\per\meter}.
\]
(b) $U_{AC} = V_A - V_C = 60 - 20 = \SI{40}{\volt}$:
\[
W = q\,U_{AC} = \pot{5}{-6}\cdot40 = \pot{2}{-4}\,\si{\joule}=\SI{2e-4}{\joule}.
\]
O trabalho é positivo porque a carga positiva se move no sentido do campo
(de maior para menor potencial).}
\end{questao}

\begin{questao}[2]{}
A figura mostra linhas equipotenciais no campo gerado por duas cargas de mesmo
módulo e sinais opostos. São dados $V_A = \SI{30}{\volt}$, $V_B = \SI{10}{\volt}$
e $V_C = \SI{-10}{\volt}$. Calcule a ddp entre A e B, e entre B e C.

\circuitoq{%
\begin{tikzpicture}[scale=0.85,line width=0.8pt]
 % duas cargas
 \shade[ball color=Vcol!25] (-3,0) circle (0.3); \node at (-3,0){\footnotesize$+$};
 \shade[ball color=Icol!25] (3,0) circle (0.3);  \node at (3,0){\footnotesize$-$};
 % equipotenciais (arcos estilizados)
 \foreach \r/\lab/\x in {1.2/A/-1.6, 2.0/B/0, 1.2/C/1.6}{}
 \draw[azulprof,thick] (-1.6,-1.6) .. controls (-2.2,0) .. (-1.6,1.6);
 \node[azulprof] at (-1.6,1.85) {$A$};
 \draw[azulprof,thick] (0,-1.9) -- (0,1.9);
 \node[azulprof] at (0,2.15) {$B$};
 \draw[azulprof,thick] (1.6,-1.6) .. controls (2.2,0) .. (1.6,1.6);
 \node[azulprof] at (1.6,1.85) {$C$};
 \node[below,font=\footnotesize] at (-1.6,-1.7) {$\SI{30}{\volt}$};
 \node[below,font=\footnotesize] at (0,-2.0) {$\SI{10}{\volt}$};
 \node[below,font=\footnotesize] at (1.6,-1.7) {$\SI{-10}{\volt}$};
\end{tikzpicture}}{Linhas equipotenciais de um dipolo elétrico.}

\resposta{$U_{AB} = \SI{20}{\volt}$; $U_{BC} = \SI{20}{\volt}$}
\resolucao{A ddp é sempre a diferença dos potenciais:
\[
U_{AB}=V_A-V_B=30-10=\SI{20}{\volt},
\qquad
U_{BC}=V_B-V_C=10-(-10)=\SI{20}{\volt}.
\]
Embora $U_{AB}=U_{BC}$, note que $V_C$ é \emph{negativo}: perto da carga negativa
o potencial fica abaixo de zero. O ponto médio B, equidistante das duas cargas de
módulos iguais, tem potencial intermediário.}
\end{questao}

\begin{questao}[2]{}
As linhas cheias da figura representam linhas de força de um campo eletrostático
e as tracejadas, linhas equipotenciais, com $V_A > V_B > V_C$. Uma partícula de
carga $q = \SI{2e-6}{\coulomb}$ é levada de A até C.

\circuitoq{%
\begin{tikzpicture}[scale=1,line width=0.8pt,>=Stealth]
 % linhas equipotenciais (verticais tracejadas)
 \foreach \x/\lab/\v in {0/A/{V_A},2.2/B/{V_B},4.4/C/{V_C}}{
   \draw[azulprof,dashed,thick] (\x,-1.4) -- (\x,1.4);
   \node[azulprof,above] at (\x,1.4) {$\v$};
   \fill[laranja] (\x,0) circle (0.07);
   \node[below,font=\footnotesize] at (\x,-0.15) {\lab};}
 % linhas de forca (horizontais com seta)
 \foreach \y in {-1,0,1}
   \draw[Vcol,->,thick] (-0.6,\y) -- (5.0,\y);
 \node[Vcol] at (5.2,0.7) {$\vec{E}$};
\end{tikzpicture}}{Linhas de força (cheias) e equipotenciais (tracejadas).}

Sobre o trabalho da força elétrica nesse deslocamento, é correto afirmar que:
\begin{alternativas}
\item é nulo, pois A e C estão na mesma linha de força.
\item é positivo, pois a carga se move no sentido do campo.
\item é negativo, pois a carga se move contra o campo.
\item depende do caminho seguido entre A e C.
\item é máximo se a carga for levada de A até B.
\end{alternativas}
\resposta{(b)}
\resolucao{A carga positiva se desloca de A (maior potencial) para C (menor
potencial), ou seja, \emph{a favor} do campo elétrico. O trabalho da força
elétrica é $W = q(V_A - V_C) > 0$. O trabalho \emph{não} depende do caminho (força
conservativa), o que elimina (d).}
\end{questao}

\begin{questao}[2]{}
A figura representa as linhas equipotenciais e as linhas de força do campo gerado
por uma carga puntiforme $Q$ fixa no vácuo. As equipotenciais são círculos
concêntricos.

\circuitoq{%
\begin{tikzpicture}[scale=0.9,line width=0.8pt,>=Stealth]
 \shade[ball color=Vcol!30] (0,0) circle (0.28);
 \node at (0,0) {\footnotesize$Q$};
 \foreach \r in {1,1.7,2.4}
   \draw[azulprof,dashed] (0,0) circle (\r);
 \foreach \a in {0,45,...,315}
   \draw[Vcol,->,thick] (\a:0.35) -- (\a:2.7);
 \fill[laranja] (0:1) circle (0.06) node[above right,font=\footnotesize]{A};
 \fill[laranja] (0:2.4) circle (0.06) node[above right,font=\footnotesize]{C};
\end{tikzpicture}}{Campo e equipotenciais de uma carga puntiforme positiva.}

Como as linhas de força \emph{apontam para fora} de $Q$, pode-se concluir que:
\begin{alternativas}
\item $Q$ é negativa e o potencial aumenta com a distância.
\item $Q$ é positiva e o potencial diminui com a distância.
\item $Q$ é positiva e o potencial aumenta com a distância.
\item o potencial é o mesmo em A e C.
\item o campo é uniforme.
\end{alternativas}
\resposta{(b)}
\resolucao{Linhas de força saindo indicam carga \textbf{positiva}. O potencial de
uma carga puntiforme, $V = \ke Q/d$, \emph{diminui} à medida que a distância
cresce: $V_A > V_C$. As equipotenciais, mais afastadas, têm menor potencial. O
campo não é uniforme --- é radial e enfraquece com $1/d^2$.}
\end{questao}

\begin{questao}[2]{}
A figura mostra a configuração das linhas de força do campo gerado por duas
partículas carregadas A e B. As linhas \emph{saem} de A e \emph{entram} em B.

\circuitoq{%
\begin{tikzpicture}[scale=0.85,line width=0.7pt,>=Stealth]
 \shade[ball color=Vcol!25] (-2,0) circle (0.3); \node at (-2,0){\footnotesize A};
 \shade[ball color=Icol!25] ( 2,0) circle (0.3); \node at (2,0){\footnotesize B};
 \foreach \a in {30,60,90,120,150,-30,-60,-90,-120,-150}
   \draw[Vcol,->] (-2,0) .. controls (\a:1.3) and ({180-\a}:1.3) .. (2,0);
 \draw[Vcol,->] (-2,0)--(-3.3,0);
 \draw[Vcol,->] (3.3,0)--(2,0);
\end{tikzpicture}}{Linhas de força de duas cargas A e B.}

Os sinais das cargas A e B são, respectivamente:
\begin{alternativas}[2]
\item positivo e positivo.
\item positivo e negativo.
\item negativo e positivo.
\item negativo e negativo.
\item ambos nulos.
\end{alternativas}
\resposta{(b)}
\resolucao{As linhas de força \emph{nascem} nas cargas positivas e \emph{morrem}
nas negativas. Como saem de A, esta é \textbf{positiva}; como entram em B, esta é
\textbf{negativa}. A configuração é a de um dipolo elétrico.}
\end{questao}

\begin{questao}[2]{}
Uma carga elétrica cria, em um ponto P situado a $\SI{20}{\centi\meter}$ dela, um
campo de intensidade $\SI{900}{\newton\per\coulomb}$. O módulo do potencial
elétrico em P é:
\begin{alternativas}[3]
\item \SI{45}{\volt}
\item \SI{90}{\volt}
\item \SI{180}{\volt}
\item \SI{450}{\volt}
\item \SI{900}{\volt}
\end{alternativas}
\resposta{(c)}
\resolucao{Para uma carga puntiforme, $E = \dfrac{\ke Q}{d^2}$ e
$V = \dfrac{\ke Q}{d}$. Dividindo, $\dfrac{V}{E} = d$, logo
\[
V = E\cdot d = 900\cdot0{,}20 = \SI{180}{\volt}.
\]
Essa relação $V = E\,d$ vale para carga puntiforme; não confundir com o campo
uniforme, onde também $U = E\,d$, mas por outro motivo.}
\end{questao}

\begin{questao}[2]{}
Próximo ao solo, em um dia limpo, a atmosfera apresenta um campo elétrico
aproximadamente uniforme, vertical, apontando para baixo, de intensidade
$\SI{120}{\volt\per\meter}$. Considere dois pontos M (mais alto) e N (mais baixo),
separados verticalmente por $\SI{10}{\meter}$.

\circuitoq{%
\begin{tikzpicture}[scale=0.9,line width=0.8pt,>=Stealth]
 \fill[cinza!25] (-2.5,-1.6) rectangle (2.5,-1.2);
 \node[cinza,font=\footnotesize] at (0,-1.4) {solo};
 \foreach \x in {-2,-1,0,1,2}
   \draw[Vcol,->,thick] (\x,1.4) -- (\x,-1.1);
 \node[Vcol] at (2.5,0.4) {$\vec{E}$};
 \fill[laranja] (-1,1.1) circle (0.08) node[left,font=\footnotesize]{M};
 \fill[laranja] (-1,-0.9) circle (0.08) node[left,font=\footnotesize]{N};
 \draw[<->,cinza] (0.6,1.1)--(0.6,-0.9) node[midway,right,font=\footnotesize]{$\SI{10}{\meter}$};
\end{tikzpicture}}{Campo elétrico atmosférico próximo ao solo.}

A diferença de potencial entre M e N vale:
\begin{alternativas}[2]
\item \SI{12}{\volt}
\item \SI{120}{\volt}
\item \SI{1200}{\volt}
\item \SI{12000}{\volt}
\item zero
\end{alternativas}
\resposta{(c)}
\resolucao{Em campo uniforme, $U = E\,d = 120\cdot10 = \SI{1200}{\volt}$. Como o
campo aponta para baixo (de M para N), M está a maior potencial:
$V_M - V_N = \SI{1200}{\volt}$.}
\end{questao}

\blocoaprofundamento

\begin{questao}[3]{}
Duas cargas $Q_1 = \SI{+2}{\micro\coulomb}$ e $Q_2 = \SI{-2}{\micro\coulomb}$
estão fixas, separadas por $\SI{20}{\centi\meter}$. Determine o potencial elétrico
no ponto médio M do segmento que as une.
\dados{$\ke = \SI{9e9}{\newton\meter\squared\per\coulomb\squared}$.}
\resposta{$V_M = 0$}
\resolucao{O potencial é uma grandeza \emph{escalar}: soma-se com sinal. No ponto
médio, cada carga está a $d = \SI{10}{\centi\meter}$:
\[
V_M = \ke\frac{Q_1}{d}+\ke\frac{Q_2}{d}
= \ke\frac{(+2)+(-2)}{d}\,\si{\micro\coulomb}=0.
\]
O potencial é nulo, mas o \emph{campo} em M \emph{não} é zero --- os vetores campo
das duas cargas apontam no mesmo sentido e se somam. Este é o ponto central: campo
nulo e potencial nulo são condições independentes.}
\end{questao}

\begin{questao}[3]{}
Uma carga puntiforme $Q = \SI{+5}{\micro\coulomb}$ está fixa no vácuo. Uma segunda
carga $q = \SI{+2}{\micro\coulomb}$ é trazida do infinito até um ponto a
$\SI{30}{\centi\meter}$ de $Q$.
\begin{itensq}
\item Qual é o potencial criado por $Q$ nesse ponto?
\item Qual é o trabalho realizado pela força externa para trazer $q$ (supondo-a
      trazida lentamente)?
\end{itensq}
\dados{$\ke = \SI{9e9}{\newton\meter\squared\per\coulomb\squared}$;
$V_\infty = 0$.}
\resposta{(a) $V = \SI{1.5e5}{\volt}$; (b) $W = \SI{0.3}{\joule}$}
\resolucao{(a) $V = \ke\dfrac{Q}{d}
= 9\cdot10^9\cdot\dfrac{\pot{5}{-6}}{0{,}3}=\pot{1.5}{5}\,\si{\volt}$.\\
(b) O trabalho da força externa iguala a energia potencial adquirida (partindo do
infinito, onde $V=0$):
\[
W = q\,(V - V_\infty) = \pot{2}{-6}\cdot\pot{1.5}{5} = \SI{0.3}{\joule}.
\]
Como ambas as cargas são positivas, elas se repelem, e o trabalho externo é
\emph{positivo} (precisamos empurrar $q$ contra a repulsão). Se as cargas tivessem
sinais opostos, o trabalho externo seria negativo.}
\end{questao}

\begin{questao}[3]{}
Um elétron ($q = \SI{-1.6e-19}{\coulomb}$, $m = \SI{9.1e-31}{\kilogram}$) parte do
repouso e é acelerado por uma ddp de $\SI{200}{\volt}$. Determine a energia
cinética adquirida e a velocidade final do elétron.
\resposta{$E_c = \SI{3.2e-17}{\joule}$; $v \approx \SI{8.4e6}{\meter\per\second}$}
\resolucao{Toda a energia potencial elétrica se converte em energia cinética:
\[
E_c = |q|\,U = \pot{1.6}{-19}\cdot200 = \pot{3.2}{-17}\,\si{\joule}.
\]
Da definição de energia cinética, $E_c = \tfrac12 m v^2$:
\[
v=\sqrt{\frac{2E_c}{m}}=\sqrt{\frac{2\cdot\pot{3.2}{-17}}{\pot{9.1}{-31}}}
\approx\SI{8.4e6}{\meter\per\second}.
\]
É o princípio de funcionamento dos antigos tubos de raios catódicos e dos
aceleradores de partículas: usar uma ddp para acelerar cargas.}
\end{questao}

\begin{questao}[3]{}
O diagrama representa o potencial elétrico $V$ criado por uma carga puntiforme em
função da distância $d$ até ela.

\circuitoqq{%
\begin{tikzpicture}
 \begin{axis}[width=9cm,height=5cm,xlabel={$d$ (m)},ylabel={$V$ (V)},
   xmin=0,xmax=6,ymin=0,ymax=90,xtick={0,1,2,3,4},ytick={0,20,40,80},
   grid=both,axis lines=left,domain=0.9:6,samples=100]
   \addplot[Vcol,very thick] {80/x};
   \addplot[only marks,mark=*,Vcol] coordinates {(1,80)(2,40)(4,20)};
   \node[Vcol,anchor=south west,font=\footnotesize] at (axis cs:2,40){$(2,\,40)$};
 \end{axis}
\end{tikzpicture}}{Potencial de uma carga puntiforme em função da distância.}

Sabendo que, a $\SI{2}{\meter}$, o potencial vale $\SI{40}{\volt}$, o potencial a
$\SI{4}{\meter}$ e o campo elétrico a $\SI{2}{\meter}$ valem, respectivamente:
\begin{alternativas}[2]
\item \SI{20}{\volt} e \SI{20}{\volt\per\meter}
\item \SI{20}{\volt} e \SI{40}{\volt\per\meter}
\item \SI{10}{\volt} e \SI{10}{\volt\per\meter}
\item \SI{40}{\volt} e \SI{40}{\volt\per\meter}
\item \SI{10}{\volt} e \SI{20}{\volt\per\meter}
\end{alternativas}
\resposta{(a)}
\resolucao{O potencial varia com $1/d$: dobrando a distância (de \SI{2}{} para
\SI{4}{\meter}), ele cai à metade, $V(4) = \SI{20}{\volt}$.\\
Para o campo, usamos $E = \dfrac{V}{d}$ (válido para carga puntiforme):
$E(2) = \dfrac{40}{2} = \SI{20}{\volt\per\meter}$.}
\end{questao}

\begin{questao}[3]{}
Na figura, uma partícula de carga $Q > 0$ está fixa em P. As superfícies A, B e C
são equipotenciais, e um ponto D situa-se sobre uma linha de força entre B e C.

\circuitoq{%
\begin{tikzpicture}[scale=0.9,line width=0.8pt,>=Stealth]
 \shade[ball color=Vcol!30] (0,0) circle (0.25); \node at (0,0){\footnotesize P};
 \foreach \r/\lab in {1.1/A,1.8/B,2.5/C}{
   \draw[azulprof,dashed] (0,0) circle (\r);
   \node[azulprof,font=\footnotesize] at (\r,0.28) {\lab};}
 \draw[Vcol,->,thick] (20:0.3)--(20:2.7);
 \fill[laranja] (20:2.15) circle (0.07) node[right,font=\footnotesize]{D};
\end{tikzpicture}}{Equipotenciais de uma carga puntiforme positiva.}

Sobre os potenciais, é correto afirmar que:
\begin{alternativas}
\item $V_A < V_B < V_C$.
\item $V_A > V_B > V_C$.
\item $V_A = V_B = V_C$.
\item $V_C$ é o maior, por estar mais longe.
\item nada se pode afirmar sem o valor de $Q$.
\end{alternativas}
\resposta{(b)}
\resolucao{Para carga positiva, $V = \ke Q/d$ \emph{decresce} com a distância. Como
A é a superfície mais interna (menor $d$) e C a mais externa (maior $d$):
$V_A > V_B > V_C$. O sinal de $Q$ determina a tendência; o \emph{valor} de $Q$ não
é necessário para ordenar os potenciais.}
\end{questao}

\begin{questao}[3]{}
Uma esfera condutora de raio $R = \SI{10}{\centi\meter}$ está carregada com
$Q = \SI{6}{\micro\coulomb}$, isolada no vácuo. Os gráficos mostram o campo
elétrico e o potencial em função da distância $d$ ao centro.

\circuitoqq{%
\begin{tikzpicture}
 \begin{axis}[width=6.2cm,height=4.2cm,xlabel={$d$ (m)},ylabel={$E$ (kN/C)},
   xmin=0,xmax=0.45,ymin=0,ymax=6200,xtick={0,0.1,0.2,0.3,0.4},
   ytick={0,2000,4000,5400},yticklabels={0,2000,4000,5400},
   grid=both,axis lines=left,domain=0.1:0.45,samples=80,
   scaled y ticks=false,y tick label style={/pgf/number format/fixed}]
   \addplot[Vcol,very thick] {9e9*6e-6/(x^2)/1000};
   \addplot[Vcol,very thick] coordinates {(0,0)(0.1,0)};
   \addplot[Vcol,dashed] coordinates {(0.1,0)(0.1,5400)};
 \end{axis}
\end{tikzpicture}\hspace{4mm}
\begin{tikzpicture}
 \begin{axis}[width=6.2cm,height=4.2cm,xlabel={$d$ (m)},ylabel={$V$ (kV)},
   xmin=0,xmax=0.45,ymin=0,ymax=620,xtick={0,0.1,0.2,0.3,0.4},
   ytick={0,180,540},grid=both,axis lines=left,domain=0.1:0.45,samples=80]
   \addplot[Icol,very thick] {9e9*6e-6/x/1000};
   \addplot[Icol,very thick] coordinates {(0,540)(0.1,540)};
 \end{axis}
\end{tikzpicture}}{Campo e potencial de uma esfera condutora carregada.}

\begin{itensq}
\item Calcule $E$ e $V$ na superfície da esfera.
\item Determine $E$ e $V$ a $\SI{30}{\centi\meter}$ do centro.
\item Explique por que, \emph{dentro} da esfera, $E = 0$ mas $V \neq 0$.
\end{itensq}
\dados{$\ke = \SI{9e9}{\newton\meter\squared\per\coulomb\squared}$.}
\resposta{(a) $E=\SI{5.4e6}{\newton\per\coulomb}$, $V=\SI{540}{\kilo\volt}$;
(b) $E=\SI{6e5}{\newton\per\coulomb}$, $V=\SI{180}{\kilo\volt}$; (c) ver comentário}
\resolucao{(a) Fora da esfera (inclusive na superfície), ela se comporta como uma
carga puntiforme concentrada no centro:
\[
E(R)=\frac{\ke Q}{R^{2}}=\frac{9\cdot10^9\cdot\pot{6}{-6}}{(0{,}1)^2}
=\pot{5.4}{6}\,\si{\newton\per\coulomb},
\]
\[
V(R)=\frac{\ke Q}{R}=\frac{9\cdot10^9\cdot\pot{6}{-6}}{0{,}1}
=\pot{5.4}{5}\,\si{\volt}=\SI{540}{\kilo\volt}.
\]
(b) A $\SI{30}{\centi\meter}$ (três vezes o raio): $E$ cai com $1/d^2$ (fator 9) e
$V$ com $1/d$ (fator 3):
\[
E=\frac{5{,}4\cdot10^6}{9}=\pot{6}{5}\,\si{\newton\per\coulomb},
\qquad
V=\frac{540}{3}=\SI{180}{\kilo\volt}.
\]
(c) \textbf{A distinção conceitual central.} Em um condutor em equilíbrio
eletrostático, toda a carga se distribui na \emph{superfície} e o campo interno é
nulo --- se não fosse, as cargas livres se moveriam e não haveria equilíbrio.\\
Mas o potencial \emph{não} é nulo: ele é a \emph{integral} do campo desde o
infinito até o ponto. Ao entrar na esfera, o campo passa a ser zero, então o
potencial \textbf{para de variar} --- ele permanece constante e igual ao valor da
superfície:
\[
V_{\text{interno}} = V(R) = \SI{540}{\kilo\volt}.
\]
Isso é coerente com $E = -\dfrac{\mathrm{d}V}{\mathrm{d}d}$: campo nulo significa
potencial \emph{constante}, não necessariamente nulo. Todo o volume do condutor é
uma única região equipotencial.\\
\textbf{Aplicação:} é o princípio da \textbf{gaiola de Faraday} --- dentro de um
condutor fechado o campo é nulo, protegendo o conteúdo, ainda que o conjunto esteja
a um potencial elevadíssimo. É por isso que se está seguro dentro de um carro
atingido por um raio.}
\end{questao}

\begin{questao}[3]{}
Três cargas puntiformes iguais, $q = \SI{1}{\micro\coulomb}$, são trazidas
\emph{uma a uma} desde o infinito até ocuparem os vértices de um triângulo
equilátero de lado $L = \SI{10}{\centi\meter}$.
\begin{itensq}
\item Calcule o trabalho externo necessário para trazer cada carga.
\item Determine a energia potencial total da configuração.
\item Se as cargas forem liberadas simultaneamente, qual é a energia cinética total
      quando estiverem infinitamente afastadas?
\end{itensq}
\dados{$\ke = \SI{9e9}{\newton\meter\squared\per\coulomb\squared}$.}
\resposta{(a) $0$, \SI{0.09}{\joule}, \SI{0.18}{\joule};
(b) $U = \SI{0.27}{\joule}$; (c) \SI{0.27}{\joule}}
\resolucao{\textbf{(a) Trabalho carga a carga.}
\begin{itemize}[leftmargin=1.6em,itemsep=1pt,topsep=2pt]
\item \textbf{1\textsuperscript{a} carga:} não há campo algum ainda, logo
      $W_1 = 0$.
\item \textbf{2\textsuperscript{a} carga:} vem contra o campo da primeira, a uma
      distância final $L$:
      \[
      W_2=\ke\frac{q^2}{L}=9\cdot10^9\frac{(\pot{1}{-6})^2}{0{,}1}=\SI{0.09}{\joule}.
      \]
\item \textbf{3\textsuperscript{a} carga:} sofre a ação das \emph{duas} anteriores,
      ambas à distância $L$:
      \[
      W_3=2\,\ke\frac{q^2}{L}=\SI{0.18}{\joule}.
      \]
\end{itemize}
\textbf{(b) Energia total da configuração.} É a soma dos trabalhos --- ou,
equivalentemente, a soma sobre os \emph{três pares} distintos:
\[
U = W_1+W_2+W_3 = 0+0{,}09+0{,}18 = \SI{0.27}{\joule}
= 3\,\ke\frac{q^2}{L}. \checkmark
\]
Note que a resposta \emph{não depende} da ordem em que as cargas foram trazidas
--- consequência de a força elétrica ser conservativa.

\textbf{(c)} Ao serem liberadas, toda a energia potencial armazenada se converte
em cinética (conservação de energia, sem dissipação):
\[
E_{c,\text{total}} = U = \SI{0.27}{\joule}.
\]
\textbf{Cuidado com o erro clássico:} a energia total \emph{não} é
$3\times W_3$ nem a soma de "cada carga vezes o potencial das outras" sem dividir
por 2 --- isso contaria cada par duas vezes. A fórmula geral é
$U=\tfrac12\sum_i q_i V_i$, onde $V_i$ é o potencial criado \emph{pelas demais}
na posição de $i$.}
\end{questao}

\begin{questao}[3]{}
Duas esferas condutoras, de raios $R$ e $3R$, estão distantes entre si e são
ligadas por um fio condutor longo e fino. A carga total do sistema é $Q$.
\begin{itensq}
\item Determine a carga final de cada esfera.
\item Compare os campos elétricos nas superfícies das duas esferas.
\item Explique o "poder das pontas" a partir desse resultado.
\end{itensq}
\resposta{(a) $Q_1=Q/4$, $Q_2=3Q/4$; (b) $E_1/E_2=3$; (c) ver comentário}
\resolucao{(a) Ligadas por um condutor, as esferas atingem o \emph{mesmo
potencial}:
\[
\frac{\ke Q_1}{R}=\frac{\ke Q_2}{3R}
\;\Longrightarrow\;
Q_2 = 3Q_1.
\]
Com $Q_1+Q_2=Q$: $Q_1=\dfrac{Q}{4}$ e $Q_2=\dfrac{3Q}{4}$.\\
(b) Campos nas superfícies:
\[
E_1=\frac{\ke Q_1}{R^{2}}=\frac{\ke Q}{4R^{2}},
\qquad
E_2=\frac{\ke Q_2}{(3R)^{2}}=\frac{3\ke Q}{4\cdot9R^{2}}=\frac{\ke Q}{12R^{2}}.
\]
\[
\frac{E_1}{E_2}=3.
\]
A esfera \textbf{menor}, embora tenha \emph{menos} carga, apresenta campo
\textbf{três vezes mais intenso} em sua superfície.\\
\textbf{Relação geral:} como $V=\ke Q/R$ e $E=\ke Q/R^{2}$, tem-se
$E = V/R$. Com $V$ comum às duas, \emph{o campo é inversamente proporcional ao
raio}.\\
(c) \textbf{Poder das pontas.} Uma "ponta" é uma região de raio de curvatura
muito pequeno. Sendo $E \propto 1/R$, o campo ali fica extremamente intenso ---
suficiente para ionizar o ar e provocar descarga, mesmo com o objeto a um
potencial moderado.\\
\textbf{Consequências práticas:}
\begin{itemize}[leftmargin=1.6em,itemsep=1pt,topsep=2pt]
\item \textbf{Para-raios} são pontiagudos \emph{de propósito}: concentram o campo
      e induzem a descarga em local controlado.
\item Equipamentos de alta tensão têm cantos \emph{arredondados} e cúpulas
      polidas, justamente para \emph{evitar} campos intensos e o efeito corona.
\end{itemize}}
\end{questao}

\blocodesafio

\begin{questao}[4]{}
Quatro cargas iguais $q=\SI{1}{\micro\coulomb}$ ocupam os vértices de um quadrado
de lado $a=\SI{10}{\centi\meter}$.
\begin{itensq}
\item Calcule a energia potencial elétrica total da configuração.
\item Quanto trabalho externo seria necessário para trazer uma quinta carga $q$
      desde o infinito até o centro do quadrado?
\item Se todas fossem liberadas, qual seria a energia cinética total no infinito?
\end{itensq}
\dados{$\ke=\SI{9e9}{}$; $\sqrt2\approx1{,}41$.}
\resposta{(a) $U\approx\SI{0.49}{\joule}$; (b) $W\approx\SI{0.51}{\joule}$;
(c) $\approx\SI{0.49}{\joule}$}
\resolucao{(a) Há $\binom{4}{2}=6$ pares: \textbf{4 lados} (distância $a$) e
\textbf{2 diagonais} (distância $a\sqrt2$):
\[
U=4\cdot\frac{\ke q^{2}}{a}+2\cdot\frac{\ke q^{2}}{a\sqrt2}
=\frac{\ke q^{2}}{a}\left(4+\frac{2}{\sqrt2}\right)
=\frac{\ke q^{2}}{a}\left(4+\sqrt2\right).
\]
Com $\dfrac{\ke q^{2}}{a}=\dfrac{9\cdot10^{9}(\pot{1}{-6})^{2}}{0{,}1}
=\SI{0.09}{\joule}$:
\[
U=0{,}09\,(4+1{,}41)\approx\SI{0.49}{\joule}.
\]
(b) A quinta carga, no centro, dista $\dfrac{a\sqrt2}{2}=\SI{0.0707}{\meter}$ de
\emph{cada} uma das quatro. O potencial no centro é
\[
V_c=4\cdot\frac{\ke q}{a\sqrt2/2}
=4\cdot\frac{9\cdot10^{9}\cdot\pot{1}{-6}}{0{,}0707}
\approx\pot{5.09}{5}\,\si{\volt},
\]
e o trabalho externo é
\[
W=q\,V_c=\pot{1}{-6}\cdot\pot{5.09}{5}\approx\SI{0.51}{\joule}.
\]
(c) Sem a quinta carga, ao liberar as quatro, toda a energia potencial da
configuração converte-se em cinética:
\[
E_{c}=U\approx\SI{0.49}{\joule}.
\]
\textbf{Atenção ao método.} Em (a) somam-se os \emph{pares} (cada par uma vez
só); em (b) usa-se $W=qV$ porque as quatro cargas já estavam posicionadas e
imóveis. Confundir os dois procedimentos --- contando pares duas vezes --- é o
erro mais comum neste tipo de problema.}
\end{questao}

\begin{questao}[4]{}
Uma esfera condutora oca de raio interno $a$ e raio externo $b$ envolve, no
centro, uma carga puntiforme $+Q$. A casca está inicialmente \emph{neutra} e
isolada.
\begin{itensq}
\item Determine as cargas induzidas nas superfícies interna e externa da casca.
\item Esboce o comportamento de $E(r)$ nas quatro regiões ($r<a$, $a<r<b$,
      $r>b$).
\item Se a casca for \textbf{aterrada}, o que muda?
\end{itensq}
\resposta{(a) $-Q$ na interna, $+Q$ na externa; (b) e (c) ver comentário}
\resolucao{(a) Dentro do \emph{material} do condutor ($a<r<b$) o campo deve ser
\textbf{nulo} em equilíbrio. Aplicando a lei de Gauss a uma superfície esférica
nessa região, a carga interna total deve ser zero:
\[
Q + q_{\text{int}} = 0 \;\Longrightarrow\; q_{\text{int}}=-Q.
\]
Como a casca é neutra no total, a superfície externa fica com
$q_{\text{ext}}=+Q$.\\
(b) \textbf{Comportamento do campo:}
\begin{center}\small\sffamily
\begin{tabular}{@{}l l l@{}}
\toprule
\textbf{Região} & \textbf{Carga envolvida} & \textbf{Campo}\\
\midrule
$r<a$      & $+Q$        & $E=\ke Q/r^{2}$ (radial, para fora)\\
$a<r<b$    & $Q-Q=0$     & $E=0$ (dentro do condutor)\\
$r>b$      & $Q-Q+Q=+Q$  & $E=\ke Q/r^{2}$\\
\bottomrule
\end{tabular}
\end{center}
Fora da casca, o campo é \emph{idêntico} ao que existiria sem ela --- a casca
neutra não blinda o exterior.\\
(c) \textbf{Aterrando a casca}, a superfície externa perde sua carga $+Q$ para a
terra (os elétrons sobem para neutralizá-la). Resultado:
\[
q_{\text{int}}=-Q,\qquad q_{\text{ext}}=0
\;\Longrightarrow\;
E = 0 \ \text{ para } r>b.
\]
\textbf{Agora sim há blindagem completa}: nenhum campo escapa para fora.\\
\textbf{Aplicação --- blindagem eletrostática.} É por isso que cabos coaxiais e
gabinetes de equipamentos sensíveis têm sua malha metálica \textbf{aterrada}:
uma blindagem flutuante (não aterrada) não impede o campo de vazar. O aterramento
é o que transforma o condutor em uma verdadeira gaiola de Faraday.}
\end{questao}

% BEGIN BANCO COMPLEMENTAR Potencial elétrico e superfícies equipotenciais
% =====================================================================
%  Banco complementar --- Potencial Eletrico e Superficies
%  Equipotenciais  (REVISADO)
%  Substitui a versao gerada por template: as 6 questoes N1 eram o
%  mesmo enunciado ("determine V a r de uma carga Q") com numeros
%  trocados, as 5 N2 eram "carga q atravessa uma ddp" repetido, e as
%  4 N3 eram "no ponto P as distancias as cargas q1 e q2 sao..."
%  igualmente clonado. As 6 N4 eram frases conceituais de uma linha.
%  O cap01c ja cobre: equipotenciais (MC), campo uniforme entre A e B,
%  tres equipotenciais, dipolo com linhas, linhas de forca com
%  equipotenciais, carga puntiforme, duas particulas A e B, campo a
%  20 cm, campo atmosferico, cargas +-2 uC, trabalho, eletron
%  acelerado, grafico V x d, superficies A/B/C, esfera condutora de
%  R=10 cm, tres cargas trazidas do infinito, duas esferas ligadas por
%  fio, quadrado de quatro cargas e esfera oca com carga central.
%  Distribuicao mantida: 6 N1 + 5 N2 + 4 N3 + 6 N4 = 21
% =====================================================================

\subsubsection*{Banco complementar --- Potencial elétrico e superfícies equipotenciais}

\begin{atencao}[Organização do banco]
O potencial é \textbf{escalar}: contribuições de várias cargas somam-se
algebricamente, com sinal, sem decomposição vetorial --- o que torna $V$ muito
mais fácil de calcular que $\vec E$. Em compensação, $V$ carrega menos
informação: pontos de $V=0$ podem ter campo intenso, e é o \emph{gradiente},
não o valor, que determina a força. O N4 explora essa distinção, a
independência do caminho e a estabilidade.
\end{atencao}

\blocofixacao

\begin{questao}[1]{}
Determine o potencial a $\SI{30}{\centi\meter}$ de uma carga
$Q=+\SI{2}{\micro\coulomb}$, adotando $V(\infty)=0$.
\dados{$\ke=\SI{9e9}{\newton\meter\squared\per\coulomb\squared}$}
\resposta{$V=\SI{60}{\kilo\volt}$}
\resolucao{
\[
V=\frac{\ke Q}{r}=\frac{\pot{9}{9}\times\pot{2}{-6}}{0{,}30}
=\frac{\pot{1.8}{4}}{0{,}30}=\pot{6}{4}\,\si{\volt}=\SI{60}{\kilo\volt}.
\]
\textbf{Atenção à potência de $r$:} o potencial cai com $1/r$, enquanto o campo
cai com $1/r^{2}$. Confundir as duas é o erro mais comum do tópico --- e o
resultado sai com a dimensão errada.}
\end{questao}

\begin{questao}[1]{}
Determine o potencial a $\SI{60}{\centi\meter}$ de uma carga
$Q=-\SI{4}{\micro\coulomb}$.
\resposta{$V=\SI{-60}{\kilo\volt}$}
\resolucao{
\[
V=\frac{\pot{9}{9}\times(-\pot{4}{-6})}{0{,}60}=-\pot{6}{4}\,\si{\volt}
=\SI{-60}{\kilo\volt}.
\]
\textbf{O sinal é do potencial, não do módulo.} Diferentemente do campo --- cuja
\emph{orientação} é dada por um vetor --- o potencial é um escalar com sinal
próprio. Cargas negativas produzem potencial negativo em toda parte.}
\end{questao}

\begin{questao}[1]{}
A que distância de $Q=+\SI{3}{\micro\coulomb}$ o potencial vale
$\SI{45}{\kilo\volt}$?
\resposta{$r=\SI{0.6}{\meter}$}
\resolucao{
\[
r=\frac{\ke Q}{V}=\frac{\pot{9}{9}\times\pot{3}{-6}}{\pot{4.5}{4}}
=\frac{\pot{2.7}{4}}{\pot{4.5}{4}}=\SI{0.6}{\meter}.
\]
\textbf{Superfície equipotencial:} \emph{todos} os pontos a
$\SI{0.6}{\meter}$ da carga estão a $\SI{45}{\kilo\volt}$. Em torno de uma
carga puntiforme, as equipotenciais são esferas concêntricas.}
\end{questao}

\begin{questao}[1]{}
Determine a carga que produz $\SI{-18}{\kilo\volt}$ a
$\SI{0.5}{\meter}$ de distância.
\resposta{$Q=-\SI{1}{\micro\coulomb}$}
\resolucao{
\[
Q=\frac{Vr}{\ke}=\frac{(-\pot{1.8}{4})(0{,}5)}{\pot{9}{9}}
=-\pot{1}{-6}\,\si{\coulomb}=-\SI{1}{\micro\coulomb}.
\]
O sinal negativo do potencial revelou diretamente o sinal da carga --- vantagem
prática de trabalhar com uma grandeza escalar.}
\end{questao}

\begin{questao}[1]{}
Sobre o potencial criado por várias cargas em um ponto, é correto afirmar:
\begin{alternativas}[2]
\item soma-se vetorialmente
\item soma-se algebricamente, com sinal
\item soma-se em módulo
\item depende do caminho até o ponto
\end{alternativas}
\resposta{(b)}
\resolucao{O potencial é \textbf{escalar}: basta somar as contribuições
$\ke q_k/r_k$ com seus sinais, sem ângulos nem componentes.\\
\textbf{Vantagem prática enorme:} calcular $V$ de dez cargas exige dez divisões
e uma soma; calcular $\vec E$ exigiria dez vetores decompostos em componentes.\\
A alternativa (d) descreve a propriedade oposta --- justamente por o campo
eletrostático ser \emph{conservativo}, o potencial de um ponto independe do
caminho, e é isso que permite defini-lo.}
\end{questao}

\begin{questao}[1]{}
Uma carga $q=+\SI{3}{\micro\coulomb}$ é colocada em um ponto onde
$V=\SI{200}{\volt}$. Determine sua energia potencial elétrica.
\resposta{$U=\SI{0.6}{\milli\joule}$}
\resolucao{
\[
U=qV=\pot{3}{-6}\times200=\pot{6}{-4}\,\si{\joule}=\SI{0.6}{\milli\joule}.
\]
\textbf{Distinção essencial:} $V$ é uma propriedade \emph{do ponto} --- existe
mesmo sem carga alguma ali. Já $U$ é uma propriedade \emph{do sistema}
carga-campo, e só faz sentido quando a carga está presente.\\
A relação $U=qV$ é o análogo elétrico de $E_p=mgh$: potencial faz o papel de
$gh$, e carga faz o papel de massa.}
\end{questao}

\blocoaplicacao

\begin{questao}[2]{}
Uma carga $q=+\SI{2}{\micro\coulomb}$ desloca-se de um ponto a
$\SI{50}{\volt}$ para outro a $\SI{20}{\volt}$.
\begin{itensq}
\item Determine o trabalho realizado pela força elétrica.
\item Determine a variação da energia potencial.
\end{itensq}
\resposta{$W=\SI{60}{\micro\joule}$; $\Delta U=\SI{-60}{\micro\joule}$}
\resolucao{(a)
\[
W=q\left(V_A-V_B\right)=\pot{2}{-6}(50-20)=\pot{6}{-5}\,\si{\joule}
=\SI{60}{\micro\joule}.
\]
(b) $\Delta U=-W=\SI{-60}{\micro\joule}$.\\
\textbf{Coerência dos sinais:} a carga é positiva e foi de maior para menor
potencial --- movimento espontâneo. A força elétrica realiza trabalho
\emph{positivo} e a energia potencial \emph{diminui}, convertendo-se em energia
cinética.\\
\textbf{Regra:} cargas positivas "descem" o potencial espontaneamente; cargas
negativas "sobem".}
\end{questao}

\begin{questao}[2]{}
Duas cargas estão fixas: $q_1=+\SI{3}{\micro\coulomb}$ a
$\SI{30}{\centi\meter}$ do ponto $P$, e $q_2=-\SI{2}{\micro\coulomb}$ a
$\SI{20}{\centi\meter}$ de $P$. Determine o potencial em $P$.
\resposta{$V=0$}
\resolucao{Somando algebricamente:
\[
V=\frac{\pot{9}{9}(\pot{3}{-6})}{0{,}30}
+\frac{\pot{9}{9}(-\pot{2}{-6})}{0{,}20}
=\pot{9}{4}-\pot{9}{4}=\SI{0}{\volt}.
\]
\textbf{Cuidado:} $V=0$ \textbf{não} significa campo nulo nem ausência de força.
As duas cargas produzem campos que \emph{não} se cancelam --- elas têm sinais
opostos, e seus campos apontam para o mesmo lado ao longo da linha que as une.\\
Uma carga colocada em $P$ teria energia potencial nula, mas sofreria força.
Essa distinção é aprofundada no bloco seguinte.}
\end{questao}

\begin{questao}[2]{}
Em um campo uniforme de $E=\SI{500}{\volt\per\meter}$, dois pontos estão
separados por $\SI{4}{\centi\meter}$ ao longo da direção do campo. Determine a
diferença de potencial.
\resposta{$U=\SI{20}{\volt}$}
\resolucao{
\[
U=E\,d=500\times0{,}04=\SI{20}{\volt},
\]
com o ponto \emph{a montante} do campo em potencial maior --- o campo aponta
sempre no sentido de $V$ decrescente.\\
\textbf{Se o deslocamento fosse perpendicular} ao campo, a diferença seria
\emph{nula}: os dois pontos estariam na mesma equipotencial. Só a componente do
deslocamento \emph{ao longo} do campo contribui.\\
\textbf{Unidades:} $\si{\volt\per\meter}$ e $\si{\newton\per\coulomb}$ são a
mesma coisa --- verifique multiplicando por metro e por coulomb.}
\end{questao}

\begin{questao}[2]{}
Um elétron é acelerado através de uma diferença de potencial de
$\SI{500}{\volt}$.
\begin{itensq}
\item Determine a energia adquirida, em joules e em elétron-volts.
\item Determine sua velocidade final, partindo do repouso.
\end{itensq}
\dados{$e=\pot{1.6}{-19}\,\si{\coulomb}$;
$m_e=\pot{9.11}{-31}\,\si{\kilogram}$}
\resposta{$\SI{500}{\electronvolt}=\pot{8}{-17}\,\si{\joule}$;
$v=\pot{1.33}{7}\,\si{\meter\per\second}$}
\resolucao{(a)
\[
W=eU=\pot{1.6}{-19}(500)=\pot{8}{-17}\,\si{\joule}.
\]
Em elétron-volts, a leitura é imediata: $\SI{500}{\electronvolt}$ --- \emph{por
definição}, $\SI{1}{\electronvolt}$ é a energia que a carga elementar adquire
sob $\SI{1}{\volt}$.\\
(b) Toda a energia vira cinética:
\[
v=\sqrt{\frac{2W}{m}}=\sqrt{\frac{2(\pot{8}{-17})}{\pot{9.11}{-31}}}
=\sqrt{\pot{1.756}{14}}=\pot{1.33}{7}\,\si{\meter\per\second}.
\]
\textbf{Verificação de consistência:} isso é $4{,}4\%$ da velocidade da luz ---
ainda no regime não relativístico, onde $\frac12mv^{2}$ é válido. Acima de cerca
de $\SI{50}{\kilo\electronvolt}$, a correção relativística passa a ser
necessária.}
\end{questao}

\begin{questao}[2]{}
Duas cargas, $q_1=+\SI{4}{\micro\coulomb}$ em $x=0$ e
$q_2=-\SI{1}{\micro\coulomb}$ em $x=\SI{0.5}{\meter}$, estão fixas. Determine o
ponto do segmento entre elas onde o potencial é nulo.
\resposta{$x=\SI{0.4}{\meter}$}
\resolucao{Impondo $V=0$ com $x$ medido a partir de $q_1$:
\[
\frac{\ke(4\mu)}{x}+\frac{\ke(-1\mu)}{0{,}5-x}=0
\;\Longrightarrow\;
\frac{4}{x}=\frac{1}{0{,}5-x}.
\]
\[
4(0{,}5-x)=x
\;\Longrightarrow\;
2=5x
\;\Longrightarrow\;
x=\SI{0.4}{\meter}.
\]
\textbf{Contraste importante com o campo:} o ponto de \emph{campo} nulo dessas
mesmas cargas fica \emph{fora} do segmento (a $\SI{1}{\meter}$ de $q_1$),
porque o campo é vetorial e só pode se cancelar onde os dois vetores forem
opostos. Já o potencial é escalar e se anula sempre que as contribuições tiverem
módulos iguais e sinais opostos --- o que ocorre dentro do segmento.\\
\textbf{Existe também um segundo ponto de $V=0$} fora do segmento, do lado da
carga menor, em $x=\SI{0.667}{\meter}$. A superfície completa de $V=0$ é uma
esfera.}
\end{questao}

\blocoaprofundamento

\begin{questao}[3]{}
Duas cargas $+Q$ e $-Q$, com $Q=\SI{2}{\micro\coulomb}$, estão separadas por
$\SI{40}{\centi\meter}$.
\begin{itensq}
\item Determine o potencial no ponto médio.
\item Determine o campo no ponto médio.
\item Explique a aparente contradição.
\end{itensq}
\resposta{$V=0$; $E=\SI{900}{\kilo\volt\per\meter}$}
\resolucao{(a) As duas contribuições têm módulos iguais e sinais opostos:
\[
V=\frac{\ke Q}{0{,}20}-\frac{\ke Q}{0{,}20}=0.
\]
(b) Os \emph{campos}, ao contrário, apontam ambos da carga positiva para a
negativa --- eles se \textbf{somam}:
\[
E=2\times\frac{\ke Q}{(0{,}20)^{2}}
=2\times\frac{\pot{9}{9}(\pot{2}{-6})}{0{,}04}
=\pot{9}{5}\,\si{\volt\per\meter}.
\]
(c) \textbf{Não há contradição.} O potencial não determina o campo --- quem o
determina é o \emph{gradiente} do potencial:
\[
\vec E=-\nabla V.
\]
Uma função pode valer zero em um ponto e ter derivada grande ali. É exatamente o
caso: caminhando ao longo do eixo, $V$ passa de positivo a negativo cruzando o
zero com inclinação acentuada.\\
\textbf{Analogia:} a altitude do nível do mar é zero, mas isso nada diz sobre a
inclinação do terreno naquele ponto. É a inclinação --- e não a altitude --- que
faz uma bola rolar.\\
\textbf{Consequência:} uma carga solta no ponto médio \emph{não} fica parada,
apesar de $V=0$ e $U=0$.}
\end{questao}

\begin{questao}[3]{}
O potencial ao longo de um eixo é dado por
$V(x)=\SI{200}{}-\SI{50}{}\,x$ (com $V$ em volts e $x$ em metros), no trecho
$0\le x\le\SI{4}{\meter}$.
\begin{itensq}
\item Determine o campo elétrico.
\item Determine o trabalho para levar $q=+\SI{2}{\micro\coulomb}$ de
      $x=1$ a $x=\SI{3}{\meter}$.
\item Identifique as superfícies equipotenciais.
\end{itensq}
\resposta{$E=\SI{50}{\volt\per\meter}$ no sentido $+x$;
$W=\SI{200}{\micro\joule}$}
\resolucao{(a)
\[
E_x=-\frac{\mathrm{d}V}{\mathrm{d}x}=-(-50)=\SI{+50}{\volt\per\meter}.
\]
Campo uniforme, apontando no sentido \emph{crescente} de $x$ --- ou seja, no
sentido em que $V$ \emph{decresce}, como sempre.\\
(b) $V(1)=\SI{150}{\volt}$ e $V(3)=\SI{50}{\volt}$:
\[
W=q\left(V_1-V_3\right)=\pot{2}{-6}(150-50)=\SI{200}{\micro\joule}.
\]
\emph{(Conferindo por força: $W=qE\,d=\pot{2}{-6}(50)(2)
=\SI{200}{\micro\joule}$ \checkmark.)}\\
(c) Como $V$ depende \emph{apenas} de $x$, as equipotenciais são
\textbf{planos perpendiculares ao eixo} $x$. O plano $x=\SI{1}{\meter}$ é toda
a superfície com $V=\SI{150}{\volt}$.\\
\textbf{Regra geral:} equipotenciais são sempre \textbf{perpendiculares} às
linhas de campo. Se não fossem, haveria componente de campo ao longo da
superfície, e mover uma carga sobre ela exigiria trabalho --- contradizendo a
definição de equipotencial.}
\end{questao}

\begin{questao}[3]{}
Uma carga $q=+\SI{1}{\micro\coulomb}$ é trazida do infinito até
$\SI{20}{\centi\meter}$ de uma carga fixa $Q=+\SI{5}{\micro\coulomb}$.
\begin{itensq}
\item Determine o trabalho realizado por um agente externo.
\item Determine a energia potencial do par.
\item Interprete o sinal.
\end{itensq}
\resposta{$W_{ext}=U=\SI{0.225}{\joule}$}
\resolucao{(a) e (b) Como a carga parte do repouso no infinito e chega em
repouso, todo o trabalho externo vira energia potencial:
\[
U=\frac{\ke Qq}{r}=\frac{\pot{9}{9}(\pot{5}{-6})(\pot{1}{-6})}{0{,}20}
=\frac{\pot{4.5}{-2}}{0{,}20}=\SI{0.225}{\joule}.
\]
\emph{(Equivalentemente, $U=qV$ com $V=\ke Q/r=\SI{225}{\kilo\volt}$.)}\\
(c) \textbf{Sinal positivo:} as cargas se repelem, então aproximá-las exige
trabalho \emph{contra} a força elétrica. O agente externo gasta energia, que
fica armazenada no sistema.\\
\textbf{Se as cargas tivessem sinais opostos}, $U$ seria negativa: elas se
atraem, e seria preciso \emph{segurá-las} durante a aproximação --- o agente
externo receberia energia.\\
\textbf{Verificação de coerência:} soltando a carga, ela é repelida e ganha
$\SI{0.225}{\joule}$ de energia cinética ao voltar ao infinito. A energia
potencial é integralmente recuperável --- marca de um campo conservativo.}
\end{questao}

\begin{questao}[3]{}
Compare, para duas cargas positivas iguais, o lugar geométrico dos pontos de
potencial nulo e o dos pontos de campo nulo.
\begin{itensq}
\item Determine onde $V=0$.
\item Determine onde $\vec E=\vec 0$.
\item Explique a diferença.
\end{itensq}
\resposta{$V$ nunca se anula (exceto no infinito); $\vec E$ se anula no ponto
médio}
\resolucao{(a) Ambas as cargas são positivas, logo cada uma contribui com
potencial \textbf{positivo} em todo ponto finito. A soma de dois positivos nunca
é zero:
\[
V=\frac{\ke Q}{r_1}+\frac{\ke Q}{r_2}>0\quad\text{sempre}.
\]
$V\to0$ apenas no limite $r\to\infty$.\\
(b) O campo é vetorial e \emph{pode} se cancelar. No ponto médio do segmento,
os dois vetores têm módulos iguais e sentidos opostos:
\[
\vec E=\vec 0\ \text{no ponto médio}.
\]
(c) \textbf{A diferença é a natureza das grandezas.}
\begin{itemize}
\item $V$ é escalar e, para cargas de mesmo sinal, as contribuições não podem se
cancelar --- só se somam.
\item $\vec E$ é vetorial e se cancela sempre que os vetores forem opostos, o que
não impõe restrição alguma sobre os sinais das cargas.
\end{itemize}
\textbf{Contraste com o caso de cargas opostas} ($+Q$ e $-Q$): ali ocorre
exatamente o inverso --- $V=0$ no ponto médio, mas $\vec E\neq\vec 0$.\\
\textbf{Síntese:} $V=0$ e $\vec E=\vec 0$ são condições \emph{independentes}.
Nenhuma implica a outra, em nenhum sentido.}
\end{questao}

\blocodesafio

\begin{questao}[4]{}
\textbf{(Demonstração --- trabalho e independência do caminho.)} Prove que o
trabalho da força elétrica entre dois pontos independe do caminho, e deduza a
relação com a diferença de potencial.
\begin{itensq}
\item Calcule o trabalho para o campo de uma carga puntiforme.
\item Generalize para uma distribuição qualquer.
\item Explique por que isso permite \emph{definir} o potencial.
\end{itensq}
\resposta{$W_{AB}=q(V_A-V_B)$, com
$V=\ke Q/r$}
\resolucao{(a) Para uma carga fixa $Q$ e uma carga de prova $q$, a força é
radial: $\vec F=\dfrac{\ke Qq}{r^{2}}\hat r$. Em um deslocamento
$\mathrm{d}\vec\ell$, apenas a componente radial contribui, pois
$\vec F\cdot\mathrm{d}\vec\ell=F\,\mathrm{d}r$:
\[
W_{AB}=\int_{r_A}^{r_B}\frac{\ke Qq}{r^{2}}\,\mathrm{d}r
=\ke Qq\left[-\frac{1}{r}\right]_{r_A}^{r_B}
=\ke Qq\left(\frac{1}{r_A}-\frac{1}{r_B}\right).
\]
\textbf{O resultado depende apenas de $r_A$ e $r_B$} --- as componentes
transversais do caminho não contribuem, por serem perpendiculares à força.\\
(b) \textbf{Generalização:} para $n$ cargas, a força total é a soma vetorial das
forças individuais, e o trabalho total é a soma dos trabalhos. Como cada parcela
independe do caminho, a soma também independe.\\
Definindo $V=\sum\ke Q_k/r_k$, obtém-se
\[
\boxed{W_{AB}=q\left(V_A-V_B\right)}
\]
(c) \textbf{Por que isso permite definir $V$.} Se o trabalho dependesse do
caminho, não haveria como atribuir um \emph{número} a cada ponto --- o "potencial"
seria ambíguo. A independência do caminho garante que
$V(P)=W_{\infty\to P}/q$ é bem definido.\\
\textbf{Formulação equivalente:} $\oint\vec E\cdot\mathrm{d}\vec\ell=0$ em
qualquer curva fechada. O campo eletrostático é \textbf{conservativo}, e essa é
exatamente a lei das malhas de Kirchhoff em sua forma de campo.\\
\textbf{Ressalva importante:} isso vale para campos \emph{eletrostáticos}.
Campos elétricos induzidos por variação de fluxo magnético \emph{não} são
conservativos, e para eles o potencial escalar deixa de existir --- assunto da
lei de Faraday.}
\end{questao}

\begin{questao}[4]{}
\textbf{(Projeto --- $V=0$ com $E\neq0$.)} Projete um par de cargas de sinais
opostos e módulos \emph{diferentes} tal que exista um ponto do segmento onde
$V=0$ mas $\vec E\neq\vec0$.
\begin{itensq}
\item Determine a condição geral.
\item Escolha valores e localize o ponto.
\item Verifique que o campo não é nulo ali.
\end{itensq}
\resposta{Condição: $|q_1|/r_1=|q_2|/r_2$; com $+4\,\si{\micro\coulomb}$ e
$-1\,\si{\micro\coulomb}$ a $\SI{0.5}{\meter}$, o ponto é
$x=\SI{0.4}{\meter}$}
\resolucao{(a) \textbf{Condição para $V=0$:}
\[
\frac{\ke q_1}{r_1}+\frac{\ke q_2}{r_2}=0
\;\Longrightarrow\;
\frac{|q_1|}{r_1}=\frac{|q_2|}{r_2},
\]
com $q_1$ e $q_2$ de sinais \emph{opostos}.\\
\textbf{Condição para $\vec E=\vec0$ no mesmo ponto:}
$\dfrac{|q_1|}{r_1^{2}}=\dfrac{|q_2|}{r_2^{2}}$.\\
As duas só coincidem se $r_1=r_2$ --- e então $|q_1|=|q_2|$, que é o caso
excluído pelo enunciado. \textbf{Com módulos diferentes, é impossível anular
ambos no mesmo ponto.}\\
(b) Tome $q_1=+\SI{4}{\micro\coulomb}$ em $x=0$ e
$q_2=-\SI{1}{\micro\coulomb}$ em $x=\SI{0.5}{\meter}$. De
$\dfrac{4}{x}=\dfrac{1}{0{,}5-x}$ vem $x=\SI{0.4}{\meter}$.\\
\emph{(Confira: $r_1=0{,}4$ e $r_2=0{,}1$, com
$4/0{,}4=1/0{,}1=10$ \checkmark.)}\\
(c) \textbf{Campo no ponto.} Ambos apontam no sentido $+x$ (afastando-se de
$q_1$, aproximando-se de $q_2$) e portanto se somam:
\[
E=\frac{\pot{9}{9}(\pot{4}{-6})}{0{,}16}
+\frac{\pot{9}{9}(\pot{1}{-6})}{0{,}01}
=\pot{2.25}{5}+\pot{9}{5}=\pot{1.125}{6}\,\si{\volt\per\meter}.
\]
Mais de um megavolt por metro --- longe de nulo.\\
\textbf{Interpretação:} uma carga colocada ali teria energia potencial nula, mas
sofreria força intensa. \textbf{$V$ e $\vec E$ carregam informações
diferentes:} $V$ é o valor, $\vec E$ é a taxa de variação.}
\end{questao}

\begin{questao}[4]{}
\textbf{(Estabilidade em um mínimo de potencial.)} Uma curva
$V(x)$ apresenta mínimo local em $x_0$.
\begin{itensq}
\item Analise a estabilidade de uma carga positiva colocada em $x_0$.
\item Repita para uma carga negativa.
\item Relacione com o teorema de Earnshaw.
\end{itensq}
\resposta{Positiva: \textbf{instável}; negativa: estável em $x$ --- mas nenhuma
é estável em três dimensões}
\resolucao{(a) A energia potencial de uma carga positiva é $U=qV$, com
$q>0$ --- ela tem a \emph{mesma} forma de $V$ e, portanto, também um mínimo em
$x_0$. Mínimo de energia significa \textbf{equilíbrio estável}... \emph{ao longo
de $x$}.\\
\emph{(Atenção ao raciocínio: é o mínimo de $U$, e não o de $V$, que define
estabilidade. Para carga positiva os dois coincidem; para negativa, invertem-se.)}\\
(b) Para carga negativa, $U=qV$ com $q<0$ \textbf{inverte} a curva: onde $V$
tem mínimo, $U$ tem \emph{máximo}. Equilíbrio \textbf{instável} em $x$.\\
(c) \textbf{Teorema de Earnshaw.} Em uma região sem cargas, o potencial
eletrostático satisfaz a equação de Laplace, $\nabla^{2}V=0$, ou seja
\[
\frac{\partial^{2}V}{\partial x^{2}}
+\frac{\partial^{2}V}{\partial y^{2}}
+\frac{\partial^{2}V}{\partial z^{2}}=0.
\]
Se $V$ tivesse um mínimo \emph{tridimensional} em $x_0$, as três segundas
derivadas seriam positivas e a soma não poderia ser zero. \textbf{Contradição.}\\
Logo, um mínimo em uma direção obriga a haver \emph{máximo} em outra --- todo
ponto crítico é um \textbf{ponto de sela}.\\
\textbf{Conclusão:} nenhuma carga pode ser mantida em equilíbrio estável apenas
por forças eletrostáticas, seja ela positiva ou negativa. A análise
unidimensional do item (a) é enganosa justamente por ignorar as outras duas
direções.\\
\textbf{Como se contorna na prática:} campos variáveis no tempo (armadilha de
Paul), campos magnéticos, ou realimentação ativa --- todas soluções que escapam
da hipótese eletrostática.}
\end{questao}

\begin{questao}[4]{}
\textbf{(Condutor em equilíbrio.)} Demonstre que toda a superfície de um
condutor em equilíbrio eletrostático é equipotencial.
\begin{itensq}
\item Prove a partir da condição de equilíbrio.
\item Mostre que o interior também está no mesmo potencial.
\item Deduza a orientação do campo na superfície.
\end{itensq}
\resposta{Campo interno nulo $\Rightarrow$ $V$ constante em todo o condutor; o
campo externo é \textbf{perpendicular} à superfície}
\resolucao{(a) \textbf{Condição de equilíbrio:} não há movimento de cargas
livres, logo o campo dentro do metal é \textbf{nulo}. Se houvesse campo, as
cargas se moveriam e o equilíbrio não existiria.\\
Tomando dois pontos $A$ e $B$ quaisquer do condutor e um caminho
\emph{interno} entre eles:
\[
V_A-V_B=\int_A^B\vec E\cdot\mathrm{d}\vec\ell=\int_A^B\vec 0\cdot\mathrm{d}\vec\ell=0.
\]
Portanto $V_A=V_B$ --- \textbf{todo o condutor é equipotencial}.\\
(b) O argumento não distinguiu superfície de interior: ele vale para
\emph{qualquer} par de pontos ligados por caminho dentro do metal. Superfície e
interior estão no mesmo potencial, e o condutor inteiro é um volume
equipotencial.\\
(c) \textbf{Orientação do campo externo.} Se houvesse componente
\emph{tangencial} à superfície, ela empurraria as cargas livres ao longo dela --
contradizendo o equilíbrio. Logo o campo externo é \textbf{perpendicular} à
superfície em todo ponto.\\
\textbf{Consequências, todas verificáveis:}
\begin{itemize}
\item A carga em excesso reside \emph{integralmente} na superfície externa.
\item A densidade superficial é maior onde a curvatura é maior (poder das
pontas).
\item Uma cavidade vazia dentro do condutor tem campo nulo --- é a blindagem
eletrostática.
\item Ligar dois condutores por fio iguala seus potenciais, e a carga se
redistribui proporcionalmente às capacitâncias.
\end{itemize}}
\end{questao}

\begin{questao}[4]{}
\textbf{(Do potencial ao campo.)} O potencial em uma região é
$V(x,y)=\SI{100}{}-\SI{20}{}\,x+\SI{5}{}\,y^{2}$ (volts, com $x,y$ em metros).
\begin{itensq}
\item Determine $\vec E(x,y)$.
\item Avalie em $(x,y)=(2;\,1)$.
\item Determine a forma das equipotenciais.
\end{itensq}
\resposta{$\vec E=(20;\,-10y)\,\si{\volt\per\meter}$; em $(2;1)$,
$\vec E=(20;\,-10)$ com $|E|=\SI{22.4}{\volt\per\meter}$}
\resolucao{(a) $\vec E=-\nabla V$, componente a componente:
\[
E_x=-\frac{\partial V}{\partial x}=-(-20)=\SI{+20}{\volt\per\meter};
\qquad
E_y=-\frac{\partial V}{\partial y}=-(10y)=-10y.
\]
Note que $E_x$ é \emph{uniforme}, enquanto $E_y$ cresce linearmente com $y$.\\
(b) Em $(2;\,1)$:
\[
\vec E=(20;\,-10)\,\si{\volt\per\meter};
\qquad
|E|=\sqrt{400+100}=\SI{22.4}{\volt\per\meter},
\]
formando $\ang{-26.6}$ com o eixo $x$.\\
\textbf{Observe:} o valor de $V$ no ponto ($V=100-40+5=\SI{65}{\volt}$) não
entrou no cálculo do campo. Só as \emph{derivadas} importam --- somar uma
constante a $V$ não altera $\vec E$.\\
(c) \textbf{Equipotenciais:} $100-20x+5y^{2}=\text{const}$, ou seja
\[
x=\frac{5y^{2}+100-C}{20},
\]
uma família de \textbf{parábolas} com eixo horizontal.\\
\textbf{Verificação de coerência:} o vetor $\vec E$ calculado deve ser
perpendicular à parábola que passa por $(2;1)$. A tangente à curva ali tem
inclinação $\mathrm{d}y/\mathrm{d}x=2$, e o campo tem inclinação
$-10/20=-1/2$ --- produto $-1$, portanto perpendiculares \checkmark.}
\end{questao}

\begin{questao}[4]{}
\textbf{(Equipotenciais e linhas de campo.)} Compare os dois conjuntos de
curvas para duas cargas positivas iguais.
\begin{itensq}
\item Descreva as equipotenciais perto de cada carga e muito longe do par.
\item Descreva as linhas de campo, incluindo o ponto médio.
\item Estabeleça as regras gerais que relacionam os dois conjuntos.
\end{itensq}
\resposta{Perto: esferas em torno de cada carga; longe: esferas em torno do par;
o ponto médio é ponto de sela do potencial}
\resolucao{(a) \textbf{Equipotenciais.}
\begin{itemize}
\item \emph{Muito perto de uma carga:} a contribuição dela domina, e as
superfícies são praticamente \textbf{esferas} centradas nessa carga.
\item \emph{Muito longe do par:} o sistema parece uma única carga $2Q$, e as
superfícies voltam a ser esferas, agora centradas no \textbf{centro do par}.
\item \emph{Na região intermediária:} as superfícies deixam de ser esferas e
assumem forma de "amendoim", chegando a se fundir em uma única superfície que
envolve as duas cargas.
\end{itemize}
(b) \textbf{Linhas de campo.} Saem radialmente de cada carga, curvam-se
afastando-se uma da outra e seguem para o infinito. \textbf{No ponto médio o
campo é nulo} --- as linhas não passam por ali. Não há linha ligando as duas
cargas, pois ambas são fontes.\\
(c) \textbf{Regras gerais.}
\begin{enumerate}
\item As linhas de campo são \textbf{perpendiculares} às equipotenciais em todo
ponto.
\item O campo aponta no sentido de $V$ \textbf{decrescente}.
\item Onde as equipotenciais estão \textbf{mais próximas}, o campo é mais
\textbf{intenso} --- pois $|E|=|\mathrm{d}V/\mathrm{d}n|$.
\item Equipotenciais nunca se cruzam (um ponto teria dois potenciais); linhas de
campo também não (haveria duas direções de força).
\item Onde $\vec E=\vec0$, as duas regras de perpendicularidade perdem sentido:
é um \textbf{ponto de sela} do potencial, e ali as equipotenciais podem se
cruzar.
\end{enumerate}
\textbf{O ponto médio é justamente esse caso excepcional:} $V$ tem mínimo ao
longo da linha que une as cargas e máximo na direção perpendicular --- exatamente
a estrutura de sela exigida pelo teorema de Earnshaw.}
\end{questao}

% END BANCO COMPLEMENTAR

\gabarito[1.4 Potencial elétrico e superfícies equipotenciais]

